% Random-sign IBLT, English version. Build: ./manuscript/random-sign/EN/build.sh
\documentclass[a4paper,11pt]{article}
\usepackage[margin=25mm]{geometry}
\usepackage[T1]{fontenc}
\usepackage{lmodern}
\usepackage{amsmath,amssymb,amsthm,booktabs,graphicx,xcolor}
\usepackage{array}
\newcommand{\exptag}[1]{\textsf{Exp-#1}}
\usepackage[colorlinks=true,linkcolor=blue!45!black,urlcolor=blue!45!black,citecolor=blue!45!black]{hyperref}
\usepackage{enumitem}
\usepackage{float}
\usepackage{xurl}
\usepackage{algorithm}
\usepackage[noend]{algpseudocode}
\setlist{nosep}
\renewcommand{\topfraction}{.85}
\renewcommand{\textfraction}{.12}
\renewcommand{\floatpagefraction}{.8}
\setlength{\emergencystretch}{2em}
\clubpenalty=10000
\widowpenalty=10000
\displaywidowpenalty=10000
\numberwithin{equation}{section}
\numberwithin{figure}{section}
\numberwithin{table}{section}
\numberwithin{algorithm}{section}
\newtheorem{theorem}{Theorem}[section]
\newtheorem{proposition}[theorem]{Proposition}
\newtheorem{lemma}[theorem]{Lemma}
\newtheorem{corollary}[theorem]{Corollary}
\newtheorem{definition}[theorem]{Definition}
\newtheorem{remark}[theorem]{Remark}
\newcommand{\E}{\mathbb E}
\newcommand{\Var}{\operatorname{Var}}
\newcommand{\Cov}{\operatorname{Cov}}
\newcommand{\Pois}{\operatorname{Pois}}
\newcommand{\Skellam}{\operatorname{Skellam}}
\newcommand{\given}{\,\middle|\,}
\let\plaineqref\eqref
\renewcommand{\eqref}[1]{Eq.~\plaineqref{#1}}
\title{Random-Sign IBLTs: Calibration-Free Self-Sizing After Decoding Failure}
\author{Min Wu\\ \normalsize\texttt{wumin@hdu.edu.cn}}
\begin{document}
\maketitle
\begin{abstract}
In a two-round set reconciliation protocol based on an Invertible Bloom Lookup Table (IBLT), when the fixed-capacity difference table of the first round suffers a decoding failure because the difference is overloaded, the receiver must estimate a safe upper bound on the unknown difference size $d$ from the sum of squared counts of the difference table, so as to set the second-round table size adaptively. The standard IBLT, however, uses the set direction of each element directly as the accumulation sign, which couples the hypergraph-overlap bias selected by failure with the sign composition of the two-sided difference; a one-sided difference amplifies the estimation bias in the same direction, forcing existing schemes to rely on tedious offline simulation grid calibration.

This paper proposes the \emph{random-sign IBLT}. By introducing a pseudorandom sign independent of the position mapping, the design decouples failure selection from the sign composition in principle, while fully preserving the number of cell fields, their widths, and the peeling recovery capability. We prove that the estimator is exactly unbiased on any fixed failed hypergraph structure; we prove a Gaussian limit for the estimation error conditioned on failure when the difference size is large, and derive a closed-form Gaussian upper bound; we further characterize a discrete limit distribution dominated by duplicate edge pairs at small difference scales, and construct a Skellam upper-bound rule that strictly controls the conditional miss rate for all $d\ge2$ and requires no offline calibration. Five hundred million paired trials and end-to-end real protocol experiments spanning several orders of magnitude of load show that the random-sign IBLT removes direction sensitivity while strictly meeting the preset risk, and that the total two-round communication cost approaches the ideal lower bound for a known difference size.
\end{abstract}


\section{Introduction}\label{sec:intro}
Many distributed systems require two network endpoints holding sets $A$ and $B$ to find their symmetric difference $A\triangle B$, the elements present on only one side. This problem is called set reconciliation. The Invertible Bloom Lookup Table (IBLT)\cite{gm,eppstein2011whatsdifference} is the standard data structure for this task. Each side builds a table of $M$ cells and hashes each element into the count, keySum, and checksum fields of $k$ cells. The receiver then subtracts the tables cell by cell. Shared elements contribute identical values to the same cells and cancel exactly, leaving only the contributions of the $d=|A\triangle B|$ differing elements. An element present only in $A$ has direction $+1$, and an element present only in $B$ has direction $-1$. When the difference size $d$ is small enough relative to the table size $M$ (the load $d/M$ is below the peeling threshold), the receiver repeatedly finds a pure cell containing a single differing element and recovers all differing elements through peeling. Once $d$ exceeds the table's capacity, peeling stalls and decoding fails.

An IBLT presents a basic tension: the table size $M$ must be fixed at construction, yet the difference size $d$ remains unknown until synchronization completes. In practice, $d$ for the same pair of nodes can vary by several orders of magnitude over time. During routine synchronization, the two sides are nearly identical and $d$ is small. After recovery from a network partition or a batch write, $d$ can grow sharply. Allocating a large table for the worst case wastes bandwidth in most routine synchronizations. Allocating a small table causes decoding failure when the difference is large. The two-round protocol\cite{wu2026selfsizing} addresses this tension. The first round sends a small table of fixed size, allowing routine small differences to complete in one round. If an oversized difference causes decoding failure, the count array of the difference table still retains the superimposed contributions of all $d$ differing elements. The receiver estimates an upper bound $U$ on $d$ from the sum of squared counts in this table. The second round allocates sufficient capacity based on $U$ and completes synchronization. The key requirement is to control the underestimate risk conditioned on failure (the conditional miss rate). If $U<d$ (a miss), insufficient capacity puts the second-round table at risk of another failure. An overly conservative $U$ wastes second-round bandwidth.

The method in \cite{wu2026selfsizing} uses the count sum and the sum of squared counts in the first-round difference table to construct an unbiased estimator $\widehat d$ and obtain an upper bound $U$ by inversion. Estimating $d$ after failure is difficult because decoding failure itself selects the sample, and this selection couples with the sign composition of the set difference (the ratio of differing elements from $A$ to those from $B$). When we estimate $d$ with this sum-of-squares method, each element's own contribution is fixed. The estimation error comes entirely from cross terms between differing elements that fall into the same cell. The standard IBLT uses each element's set-membership direction ($+1$ or $-1$) as its accumulation sign in count. Two elements with the same direction add when they share a cell, increasing the sum of squares. Elements with opposite directions cancel. Peeling fails because the position mapping of differing elements to cells has too much overlap. Together, these effects produce systematic bias in the failure sample. For a one-sided difference, all differing elements come from the same side. This is common in one-way synchronization, for example when a replica lacks only data newly written to the primary. All cross terms then have the same sign, so the high overlap selected by failure is amplified in the sum of squared counts, substantially biasing the estimate. In two-way synchronization where each side writes new data, the numbers of differing elements on the two sides are comparable. For this balanced difference, positive and negative cross terms largely cancel, leaving the post-failure mean estimate nearly unchanged. The bias magnitude varies continuously with the ratio of differing elements on the two sides, and its sign reverses with the number of hashes $k$ (Section~\ref{sec:model}). After decoding failure, the receiver does not know whether the difference comes from one side or both. The observed sum of squared counts therefore cannot reveal how much is due to the true $d$ and how much to the direction bias induced by failure selection. Thus, except in deep overload, the upper-bound rule for the standard IBLT in \cite{wu2026selfsizing} requires offline simulation calibration for each configuration.

Our core insight is that separating an element's set direction from its accumulation sign removes the source of this coupling. We therefore propose the random-sign IBLT. The change is lightweight: the two sides share a random-sign hash independent of the position mapping, assigning each key a random $+1$ or $-1$. On insertion, the element's random sign multiplies its contributions to all three fields: count, keySum, and checksum.

This modification supports both decoding and estimation. Shared elements receive the same random sign in both tables and still cancel exactly after subtraction. A cell containing one unrecovered differing element has count $+1$ or $-1$. Multiplying keySum by count cancels the random sign exactly (since $(\pm1)^2=1$) and recovers the original key. Computing the key's random sign and comparing it with count then reveals whether the key belongs to set $A$ or $B$. Under ideal purity, the random-sign IBLT therefore peels exactly the same elements as the standard IBLT at every step. The decoding threshold, failure probability, and number and byte width of cell fields remain unchanged. For estimation, each differing element's coding coefficient in the difference table becomes an independent uniform random sign, whether the difference is one-sided or balanced. Peeling success depends only on which cells the elements select (the mapping), while cancellation of the estimation error depends on the random signs. These are independent. For every mapping selected by failure, the random signs remain independent and uniform, so their contributions cancel in expectation. Using random signs to cancel cross terms is a classical idea from the AMS sketch and CountSketch\cite{alon1999spacecomplexity,charikar2004countsketch}. We implement it in a peelable data structure, enabling one table to support both difference recovery and difference-size estimation and allowing rigorous statistical analysis of the post-failure estimator.

\paragraph{Main results and contributions.}
We prove that the random-sign IBLT preserves the standard IBLT's decoding behavior under ideal purity and that the joint distribution of the difference-table statistic and the mapping is identical for all sign compositions. Section~\ref{sec:design} presents the encoding and peeling algorithm that achieves this decoupling. We then prove that the estimator $\widehat d$ is exactly unbiased on every fixed mapping, giving conditional mean $d$ after failure, and fully characterize the conditional distribution of the estimation error across all load scales (Sections~\ref{sec:measurement} and \ref{sec:inversion}). For large differences, we derive a Gaussian rule from the conditional Gaussian limit. For small differences, we characterize a discrete limit distribution and construct a Skellam rule that uniformly controls the conditional miss rate for all difference sizes without offline simulation calibration. We also provide an efficient planted sampling method for rare failures at low load. The experiments (Section~\ref{sec:exp}) verify that random-sign and standard encoding have identical decoding behavior. They show that the Skellam rule keeps the conditional miss rate strictly within the target risk across the full load range spanning several orders of magnitude. With this rule, the two-round protocol reliably completes synchronization at a total transmission cost close to the ideal lower bound for known $d$.

\paragraph{Structure of the paper.}
Section~\ref{sec:model} introduces the standard IBLT model and analyzes its post-failure estimation bias. Section~\ref{sec:design} presents the random-sign IBLT and the two-round protocol. Section~\ref{sec:measurement} proves exact unbiasedness given the mapping and derives the Gaussian rule for large differences. Section~\ref{sec:inversion} derives the discrete limit distribution for small differences and the Skellam rule for all scales. Section~\ref{sec:exp} reports the experimental evaluation. Sections~\ref{sec:related} and \ref{sec:conclusion} discuss related work, rule selection, and computational cost. The appendix contains complete proofs and supplementary experiments.

\section{Model and post-failure estimation bias}\label{sec:model}
We introduce the standard IBLT and the ideal peeling model, define the conditional miss rate required by the two-round protocol, characterize failure at low load, and explain how the standard IBLT's post-failure bias depends on sign composition. Appendix~\ref{app:notations} collects the notation used throughout the paper.

\subsection{The standard IBLT and the ideal peeling model}
The two sides hold sets $A$ and $B$ and seek their symmetric difference $A\triangle B$, whose size $d=|A\triangle B|$ is unknown in advance. The standard IBLT consists of $M$ cells, each with three fields: count, keySum, and checksum. To insert a key $x$, the cell hash selects $k$ distinct cells from $[M]$ and adds $1$, $x$, and $\mathrm{chk}(x)$ to their respective fields. Each side builds a table, and the tables are subtracted cell by cell. Shared keys contribute identical values to the same cells in both tables, so their contributions cancel in all three fields. Only the contributions of the $d$ differing keys remain. A key present only in $A$ has direction $s_i=+1$, and a key present only in $B$ has direction $s_i=-1$. We denote the numbers of keys in the two directions by $d_+$ and $d_-$. The ordered pair $(d_+,d_-)$ is the \emph{sign composition}, and $b=d_+-d_-=|A|-|B|$ is the difference between the set sizes. We use two extreme sign compositions throughout: a \emph{one-sided difference}, where all differing keys come from one side ($|b|=d$), and a \emph{balanced difference}, where both sides have the same number of differing keys ($d_+=d_-$, $b=0$). Both correspond to common synchronization scenarios. A one-sided difference arises in one-way synchronization when only one side inserts or deletes entries. A balanced difference arises in two-way synchronization when each side modifies a small number of entries and the insertions and deletions roughly cancel. Practical loads usually lie between these extremes, depending on how much each side changes. Decoding proceeds by peeling. We repeatedly find a pure cell containing one unrecovered differing key, recover the key by multiplying keySum by count, verify it with the checksum, and remove its contribution from all $k$ cells it selects. Removal may make other cells pure. Peeling continues until no pure cell remains.

Under ideal purity (a cell with count $\pm1$ contains exactly one unrecovered differing key), peeling success depends entirely on the mapping of differing keys to cells. View the $M$ cells as vertices and the $k$ cells selected by each differing key as a hyperedge. The difference table then corresponds to a random $k$-uniform hypergraph: the $d$ differing keys independently and uniformly choose $k$-subsets $S_i\subseteq[M]$, with different keys allowed to choose the same subset. We denote the unsigned mapping with key labels by $\mathcal G$ and the overlap between two edges by $H_{ij}=|S_i\cap S_j|$, the number of cells shared by differing keys $i$ and $j$. On $\mathcal G$, pure cells are exactly the unsigned degree-one vertices. Peeling repeatedly deletes edges incident to these vertices. The deletion order does not affect the terminal state, which is always the hypergraph's 2-core. A nonempty 2-core causes decoding failure, denoted by event $F$, with failure probability $\Pr(F)$. Thus $F$ depends only on the unsigned mapping $\mathcal G$. When the load $d/M$ is below the critical value $\alpha_c$ (about $0.818$ for $k=3$ and $0.772$ for $k=4$\cite{molloy,dw}), peeling succeeds with high probability. Rare failures arise from small cores such as duplicate edge pairs, and their probability decays polynomially in $M$. In the critical window of width of order $M^{-1/2}$ around $\alpha_c$, the failure probability rises from near zero to near one\cite{dm}. Above $\alpha_c$, peeling fails almost surely.

\subsection{The two-round protocol and the conditional miss rate}
Because the capacity $M$ must be fixed before $d$ is known, the two-round protocol\cite{wu2026selfsizing} first sends a small table of fixed size $M_1$. For a small difference, peeling succeeds in the first round. After a first-round failure, the receiver obtains an upper bound $U$ on $d$ from the difference table's count array and chooses the second-round size $m(U)$ based on $U$. The second round re-encodes and decodes using a fresh, independent cell hash.

Let $C_v$ denote the count of cell $v$ in the difference table. Following \cite{wu2026selfsizing}, we define the estimator $\widehat d$ of $d$ using the count sum $\sum_vC_v$ and sum of squares $\sum_vC_v^2$ saved before peeling begins:
\begin{equation}
\widehat d=\frac{1}{k(1-k/M)}\biggl(\sum_vC_v^2-\frac1M\Bigl(\sum_vC_v\Bigr)^2\biggr),\label{eq:dhat}
\end{equation}
and denote its unconditional standard deviation scale by $s(n)$:
\begin{equation}
s(n)=\sqrt{\frac{2n(n-1)}{M-1}}.\label{eq:s}
\end{equation}
For the standard IBLT, $\widehat d$ is exactly unbiased for $d$ and has variance $s(d)^2$ for every sign composition. This expectation is unconditional: it averages over all runs, regardless of failure $F$. If we evaluated the upper bound $U$ over all runs, adding a few standard deviations to $\widehat d$ would control the underestimate probability. In the two-round protocol, a successful first round ends the protocol immediately, without computing $U$. Only after decoding failure (event $F$) does the receiver compute $U$ to size the second-round table. If $U<d$, the second-round capacity falls below the prescribed margin, risking another failure. We call the event $U<d$ a \emph{miss} and its probability conditioned on first-round failure, $\Pr(U<d\mid F)$, the \emph{conditional miss rate}. The sizing rule must therefore satisfy
\begin{equation}
\Pr(U<d\mid F)\le u,\label{eq:target}
\end{equation}
where $u$ is the target risk. In deployment, we can choose $0.01$ or $0.05$, corresponding to a 1 \% or 5 \% underestimate risk.

The two-round protocol incurs a miss only when the first round fails and the upper bound underestimates $d$. This event has \emph{joint miss rate} $\Pr(F,\,U<d)=\Pr(F)\Pr(U<d\mid F)$. Controlling the conditional miss rate at $u$ bounds the joint miss rate by $u\Pr(F)$, both in overload, where failure is almost certain, and at low load, where it is rare. For failure samples without a miss ($U\ge d$), the second-round table has sufficient capacity margin to decode with high probability. Section~\ref{sec:protocol} gives the complete decomposition of the total two-round failure rate.

\subsection{Failure at low load and the planted distribution}\label{sec:plant}
\eqref{eq:target} must hold at all loads. Low load requires separate treatment. When $d$ is much smaller than $M$, the failure probability decays polynomially in $M$. Failure is then rare, making the distribution conditioned on $F$ difficult to analyze directly and requiring many simulated mappings for each failure sample. The structure of failure in this region is simple. Two edges with the same mapping form a duplicate edge pair. A single edge cannot form a 2-core, and two edges form a 2-core exactly when they are a duplicate edge pair. When $d=o(M)$, 2-cores with three or more edges are negligible relative to duplicate edge pairs. Failure is therefore approximately equivalent to the existence of a duplicate edge pair. Conditioning on $F$ is close to planting a duplicate edge pair in a random mapping. We thus define the \emph{planted distribution}: choose a pair of keys uniformly, assign both the same uniformly random $k$-subset, and let all remaining keys choose their mappings independently and uniformly.

\begin{proposition}[Failure is determined by duplicate edge pairs when $d=o(M)$]\label{prop:plant}
Fix $k\ge3$ and let $2\le d=o(M)$. The failure probability is asymptotically equal to the expected number of duplicate edge pairs: $\Pr(F)\sim\binom d2/\binom Mk$. The total variation distance between the distribution of $\mathcal G$ conditioned on $F$ and the planted distribution tends to zero.
\end{proposition}

Goodrich and Mitzenmacher\cite{gm} identify duplicate edge pairs as the main source of IBLT failure. The finite-length analysis of Yugawa and Wadayama\cite{yw} likewise shows that stopping sets of size 2 dominate. Proposition~\ref{prop:plant} extends this observation to a total variation approximation of the entire mapping for all $d=o(M)$. Lemma~\ref{lem:plant} in Appendix~\ref{app:graph} gives the quantitative error bound and proof. Proposition~\ref{prop:plant} supports both theory and experiments. It replaces conditioning on failure with the planted distribution, under which direct calculations are straightforward. The limit theorems in Sections~\ref{sec:measurement} and \ref{sec:inversion} use this approximation to handle conditioning on failure when $d=o(M)$. The risk decomposition in Section~\ref{sec:protocol} also uses it to estimate the first-round failure probability. For experiments, it enables direct sampling from the planted distribution, where every sample yields a failure. Section~\ref{sec:exp-setup} uses this planted sampling to evaluate low-load configurations beyond the reach of direct sampling and checks its finite-size accuracy where both methods are feasible.

\subsection{Post-failure bias of the standard IBLT}
Controlling \eqref{eq:target} is difficult because the failure event $F$ selects the mapping $\mathcal G$. Failure samples concentrate on mappings with large overlap $H_{ij}$, where differing keys share many cells. At low load, these mappings mainly contain duplicate edge pairs (Proposition~\ref{prop:plant}). In the critical window and at higher load, they contain large cores with dense overlap. Any statistic that depends on $H_{ij}$ therefore deviates from its unconditional distribution when conditioned on $F$.

In the standard IBLT, the difference-table count is $C_v=\sum_{i:v\in S_i}s_i$, and each differing key's coding coefficient is its fixed set direction $s_i\in\{+1,-1\}$. In the sum of $C_v^2$, the diagonal terms contribute exactly $kd$. The estimation error comes entirely from the cross terms $s_is_jH_{ij}$ between keys sharing cells. Excess overlap raises $\widehat d$ for same-sign pairs and lowers $\widehat d$ for opposite-sign pairs. The conditional relative bias of $\widehat d$ thus factors exactly into two terms.
\begin{proposition}[Bias decomposition of the standard IBLT]\label{prop:projection}
In the standard IBLT, for any sign composition $(d_+,d_-)$,
\begin{equation}
\frac{\E[\widehat d\mid F]}{d}-1=c_F\Bigl(\frac{b^2}{d}-1\Bigr),\qquad c_F=\frac{\E[H_{12}\mid F]-k^2/M}{k(1-k/M)}.\label{eq:projection}
\end{equation}
\end{proposition}
The proof is in Appendix~\ref{app:std}. Since $F$ is invariant under permutations of key labels, $\E[H_{ij}\mid F]$ is identical for all $i\ne j$. The formula uses differing keys $1$ and $2$ as a representative pair. The factor $c_F$ depends only on failure selection of the mapping and measures the excess of inter-edge overlap in the failure sample over its unconditional mean $k^2/M$. The factor $b^2/d-1$ depends only on the sign composition $(d_+,d_-)$. For a one-sided difference ($|b|=d$), every intersecting edge pair has sign product $+1$, so $b^2/d-1=d-1$. All excess overlap contributes with the same sign, fully amplifying its effect on $\widehat d$. For a balanced difference ($b=0$), there are $d/2$ fewer same-sign pairs than opposite-sign pairs, giving a relative bias of $-c_F$, whose magnitude is $1/(d-1)$ of the one-sided value.
\begin{figure}[!htbp]
\centering\includegraphics[width=\linewidth]{../figure/fig-bias.pdf}
\caption{Post-failure relative bias of the standard IBLT versus $(b/d)^2$ ($M=16384$). (a) $k=3$. (b) $k=6$. Each color denotes a failure probability. All configurations lie in the critical window: $d/M$ is 0.812--0.822 for $k=3$ and 0.633--0.639 for $k=6$. Points are measurements, and lines follow \eqref{eq:projection}.}
\label{fig:bias}
\end{figure}

Figure~\ref{fig:bias} shows this decomposition at $M=16384$ for five failure probabilities (0.10 to 0.90) each for $k=3$ and $k=6$. Within each configuration, all five sign compositions share the same batch of mappings and therefore identical peeling outcomes. Since $b^2/d-1=d(b/d)^2-1$, the relative bias at each failure probability is linear in $(b/d)^2$, with slope determined by $c_F$. We determine $c_F$ from the one-sided point. The remaining predictions of \eqref{eq:projection} differ from the measurements by at most twice the standard error of the difference. For $k=3$ at failure probability 0.10, the one-sided relative bias is $+0.70\%$, corresponding to a bias of about 0.64 standard deviations $s(d)$. As the failure probability increases, the failure sample approaches the full sample, weakening selection and reducing $|c_F|$. The sign of $c_F$ also depends on $k$. For $k=6$, $c_F<0$, so the one-sided failure sample is biased low, with relative bias $-0.43\%$ at failure probability 0.10 (Figure~\ref{fig:bias}(b)). The bias changes sign with $k$ and varies with sign composition according to \eqref{eq:projection}. A fixed correction therefore cannot offset it (Figure~\ref{fig:composition}).

In practice, the receiver does not know $(d_+,d_-)$ in advance. The difference table's count array alone cannot reveal whether a larger observed $\widehat d$ reflects a larger true $d$ or failure-induced bias for a one-sided difference. The upper-bound rule in \cite{wu2026selfsizing} therefore requires simulation calibration by load interval outside deep overload. By \eqref{eq:projection}, the bias arises because the standard IBLT uses the set direction $s_i$ as the accumulation sign in count, with sign composition determining the factor that amplifies $c_F$. For a balanced difference, this factor is $-1$, leaving a relative bias of $-c_F$. This suggests encoding differing keys with independent, uniform random signs so that every cross term has zero mean for each mapping, while preserving peeling's ability to recover each key's original set direction. The post-failure estimator would then be unbiased for every sign composition. We develop this idea in the random-sign IBLT in Section~\ref{sec:design}.

\section{The random-sign IBLT}\label{sec:design}
We introduce the random-sign IBLT. The design multiplies each key by a random sign shared by both sides during insertion and deletion. The cell fields, field widths, and peeling process remain unchanged. We first give the encoding and decoding rules and prove ideal decoding equivalence with the standard IBLT and independence of the coding coefficients from the mapping. We then explain how the two-round protocol uses the lower bound $|b|$ to determine the second-round capacity.

\subsection{Encoding rules, decoding equivalence, and sign independence}
The two sides agree in advance on a sign hash $\varepsilon:\mathcal U\to\{-1,+1\}$ independent of the cell hash. Each sign is generated directly from the key and a shared seed, requiring no per-key storage. Algorithm~\ref{alg:encode} gives the complete encoding and decoding procedure. The lines marked $\star$ show all changes from the standard IBLT: applying the sign during insertion and deletion and recovering the direction on output. The number of cell fields, field widths, and peeling control flow remain unchanged. We describe encoding and decoding in turn.

\begin{algorithm}[!htbp]
\caption{Encoding and decoding of the random-sign IBLT. Lines marked $\star$ differ from the standard IBLT; field arithmetic is modulo $2^{64}$.}
\label{alg:encode}
\begin{algorithmic}[1]
\Procedure{Insert}{$x$}
\For{each cell $v$ of $x$}
\State $\mathrm{count}_v\gets\mathrm{count}_v+\varepsilon(x)$\Comment{$\star$ standard IBLT adds $1$}
\State $\mathrm{keySum}_v\gets\mathrm{keySum}_v+\varepsilon(x)\,x$\Comment{$\star$ standard IBLT adds $x$}
\State $\mathrm{checksum}_v\gets\mathrm{checksum}_v+\varepsilon(x)\,\mathrm{chk}(x)$\Comment{$\star$ standard IBLT adds $\mathrm{chk}(x)$}
\EndFor
\EndProcedure
\Procedure{Delete}{$x$}
\State as in \Call{Insert}{$x$}, subtract the same quantities from the three fields
\EndProcedure
\Procedure{Peel}{difference table}
\State add cells with count $\pm1$ to the candidate queue
\While{the candidate queue is nonempty}
\State remove cell $v$; if $C_v\ne\pm1$, process the next candidate
\State $x\gets C_v\,\mathrm{keySum}_v$
\If{$\mathrm{checksum}_v=C_v\,\mathrm{chk}(x)$ and the cells of $x$ include $v$}
\State output key $x$ and direction $C_v\,\varepsilon(x)$\Comment{$\star$ standard IBLT outputs $C_v$}
\State if the direction is $+1$, call \Call{Delete}{$x$}, otherwise call \Call{Insert}{$x$}
\State add cells whose count becomes $\pm1$ to the candidate queue
\EndIf
\EndWhile
\State \Return whether all three fields of every cell are zero
\EndProcedure
\end{algorithmic}
\end{algorithm}

Encoding follows lines 1--7 of Algorithm~\ref{alg:encode}. To insert a key $x$, we add $\varepsilon(x)$, $\varepsilon(x)\,x$, and $\varepsilon(x)\,\mathrm{chk}(x)$ to the three fields of each of its $k$ cells. Deletion subtracts the same quantities. We accumulate keySum and checksum modulo $2^{64}$. In the set model, these two fields also support the bitwise XOR implementation commonly used in the standard IBLT (Appendix~\ref{app:xor}). Our analysis and experiments use the additive form. Subtracting the two tables cell by cell cancels the identical contributions of each shared key in all three fields. A differing key $i$ has coding coefficient $s_i\varepsilon_i$ in the difference table, where $\varepsilon_i=\varepsilon(x_i)$. Thus, cell $v$ in the difference table has fields
\begin{equation}
C_v=\sum_{i:v\in S_i}s_i\varepsilon_i,\qquad
\mathrm{keySum}_v=\sum_{i:v\in S_i}s_i\varepsilon_i\,x_i,\qquad
\mathrm{checksum}_v=\sum_{i:v\in S_i}s_i\varepsilon_i\,\mathrm{chk}(x_i).\label{eq:fields}
\end{equation}
We treat $\varepsilon_i$ as iid uniform signs independent of the mapping $\mathcal G$.

Decoding follows lines 8--17 and uses the standard IBLT peeling process. A pure cell $v$ containing a single unrecovered differing key $x_i$ has count $C_v=s_i\varepsilon_i\in\{-1,+1\}$. Multiplying $\mathrm{keySum}_v$ by $C_v$ cancels the sign $s_i\varepsilon_i$ and yields the candidate key $x_i$ (line 12). Once the candidate passes the checksum and cell-membership checks (line 13), we multiply $C_v$ by $\varepsilon(x_i)$. Since $\varepsilon_i^2=1$, this recovers the original set direction $s_i=C_v\,\varepsilon(x_i)$ (line 14). We then call \textsc{Delete} or \textsc{Insert} according to the recovered direction $s_i$ (line 15). Both operations apply $\varepsilon(x_i)$ again in all $k$ cells of $x_i$. The contribution subtracted from each field is therefore exactly the key's contribution multiplied by its coding coefficient $s_i\varepsilon_i=C_v$. After each peeling step, the three fields still equal the encoded sum of the unrecovered keys. As in the standard IBLT, peeling succeeds when all three fields of every cell are zero (line 17).

Figure~\ref{fig:structure} compares the standard and random-sign IBLTs under the same mapping in a small one-sided example ($s_1=\dots=s_4=+1$). Both encodings have the same pure cells: the degree-$1$ cells 1, 4, 7, and 8, shaded in the figure. All have counts in $\{\pm1\}$. In pure cell 4, for example, the random sign gives $C_4=-1$. Multiplying by $\varepsilon_2=-1$ recovers the true direction of $x_2$ as $s_2=C_4\varepsilon_2=+1$. The two encodings follow identical peeling processes. Their sums of squared counts differ substantially. The single-key diagonal terms sum to $kd=12$ in both encodings. In the shared cells (cells 2, 3, 5, and 6), the standard IBLT accumulates contributions with the same sign. Each shared cell has count $2$, all cross terms are positive, and $\sum_vC_v^2=20$. In the random-sign IBLT, contributions along solid lines ($\varepsilon_i=+1$) cancel those along dashed lines ($\varepsilon_i=-1$) in the shared cells. Each shared cell in this example contains exactly one positive and one negative key. All cross terms are negative, the counts cancel to $0$, and $\sum_vC_v^2=4$. The next two propositions generalize the decoding equivalence and sign independence illustrated in Figure~\ref{fig:structure}.

\begin{figure}[!htbp]
\centering\includegraphics[width=.92\linewidth]{../figure/fig-structure.pdf}
\caption{The standard IBLT and the random-sign IBLT under the same mapping ($k=3$, $M=8$, one-sided difference $d=4$). The shaded cells have count $\pm1$; in (b), solid and dashed lines correspond to $\varepsilon_i=+1$ and $\varepsilon_i=-1$, respectively.}
\label{fig:structure}
\end{figure}

Random signs allow a degree-three cell with coding coefficients $(+1,+1,-1)$ to have count $+1$. This cannot occur in the standard IBLT under a one-sided difference. The candidate key read from such a cell fails the checksum and cell-membership checks, so peeling skips the cell until it contains a single key. Sign cancellation thus preserves correct key deletion and the stopping condition. Proposition~\ref{prop:decode} formalizes this property.

\begin{proposition}[Ideal decoding equivalence]\label{prop:decode}
Assume the pure-cell test has no false positives and field arithmetic does not overflow. At every intermediate state, the random-sign IBLT can remove exactly the keys removable by unsigned degree-one peeling. Under the same cell hash, the random-sign and standard IBLTs therefore leave the same 2-core. The failure event $F$ depends only on the unsigned mapping $\mathcal G$. The decoding threshold, failure probability, and recovered result are identical to those of the standard IBLT.
\end{proposition}

\begin{proof}
We induct on the number of peeling steps. By the field invariant, every degree-one cell has count $\pm1$ and passes the check. Ideal purity excludes cells of any other degree, so the two processes can follow the same cell-selection order. The full proof is in Appendix~\ref{app:proof-decode}.
\end{proof}

In an implementation, the finite 64-bit checksum gives a nonzero false-acceptance probability. Under an ideal hash, each check of a non-pure cell falsely accepts with probability at most $2^{-64}$. Section~\ref{sec:protocol} accounts for this term separately in the risk upper bound for the two-round protocol. Section~\ref{sec:exp-decode} verifies the implemented decoding equivalence in five hundred million paired trials with real IBLTs. Table~\ref{tab:decode-paired} in the appendix summarizes success-failure agreement and recovery checks for each configuration. Section~\ref{sec:exp-cost} reports decoding time.

Random signs retain their distribution after conditioning on failure. Failure selects mappings with large overlap, so conditioning usually changes the distribution of random variables associated with the mapping. This selection causes the bias of the standard IBLT under a one-sided difference. By Proposition~\ref{prop:decode}, failure depends only on the mapping. The signs come from a sign hash independent of the cell hash and remain iid uniform after we fix the mapping. Proposition~\ref{prop:composition} states this property precisely.

\begin{proposition}[Signs are independent of the mapping]\label{prop:composition}
For any fixed direction vector $(s_i)$, the coding coefficients $(s_i\varepsilon_i)$ are iid uniform signs independent of the mapping $\mathcal G$. Thus, failure $F$ is independent of the coding coefficients. Given $F$, or even the entire $\mathcal G$, the coding coefficients remain iid uniform signs. Statistics constructed from $(C_v)$, $\mathcal G$, and $F$ have the same joint distribution for all sign compositions $(d_+,d_-)$.
\end{proposition}

\begin{proof}
For fixed $(s_i)$, the map $(\varepsilon_i)\mapsto(s_i\varepsilon_i)$ is a bijection on $\{\pm1\}^d$ that preserves the uniform distribution. Since $(\varepsilon_i)$ is independent of $\mathcal G$, $(s_i\varepsilon_i)$ is an iid uniform sign sequence independent of $\mathcal G$. The joint distribution of $(s_i\varepsilon_i,\mathcal G)$ therefore does not depend on $(s_i)$. By Proposition~\ref{prop:decode}, $F$ is a function of $\mathcal G$, so conditioning on $F$ or $\mathcal G$ preserves the distribution of the coding coefficients. By \eqref{eq:fields}, $(C_v)$ is determined by $(s_i\varepsilon_i)$ and $\mathcal G$. Statistics constructed from $(C_v)$, $\mathcal G$, and $F$ thus have a joint distribution independent of $(s_i)$ and identical for all $(d_+,d_-)$.
\end{proof}

Section~\ref{sec:exp-e2e} and Table~\ref{tab:vldb} in the appendix illustrate this property numerically in the protocol. For a fixed configuration and calibration factor, random-sign encoding gives similar conditional miss rates under one-sided and balanced differences. Figure~\ref{fig:encoding} compares the post-failure distributions of the standard and random-sign encodings.

Propositions~\ref{prop:decode} and \ref{prop:composition} underpin our analysis. Decoding success depends only on $\mathcal G$, and the estimator's randomness comes from the independent signs. We can therefore analyze the post-failure distribution by first fixing the mapping and then studying the quadratic form in iid signs on the overlap graph. The results hold for all sign compositions. The joint distribution of the coding coefficients $s_i\varepsilon_i$ is independent of the direction vector $(s_i)$, so we can treat $s_i\varepsilon_i$ and $\varepsilon_i$ identically in the analysis. We write both as $\varepsilon_i$ below.

\subsection{\texorpdfstring{The two-round protocol flow and the lower bound $|b|$}{The two-round protocol flow and the lower bound |b|}}\label{sec:flow}
In the two-round protocol, the receiver subtracts the tables cell by cell and runs peeling on the difference table. For the standard IBLT, the count sum $\frac1k\sum_vC_v=b=d_+-d_-$ directly gives the deterministic lower bound $d\ge|b|$. For the random-sign IBLT, $\frac1k\sum_vC_v=\sum_i\varepsilon_i$ depends only on the random signs and no longer encodes the set-size difference $b$. The sender therefore attaches the set size $|A|$ (8 bytes) to the first-round table, allowing the receiver to compute $b=|A|-|B|$ directly.

If first-round peeling succeeds, the protocol ends immediately. After a first-round failure, the receiver computes an upper bound $U$ from the raw counts $(C_v)$ of the difference table. It incorporates the known lower bound $|b|$ by taking $U_b=\max(U,|b|)$ (for a one-sided difference, $|b|=d$). The second-round table has size $m(U_b)=\lceil U_b/(0.8\alpha_c)\rceil$. We re-encode and decode using a fresh cell hash independent of the first round. Whenever $U_b\ge d$, the second-round load is at most $0.8\alpha_c$, leaving a margin below the peeling threshold. The second round therefore succeeds with high probability. The analysis applies to any constant target load below $\alpha_c$; we use 0.8. Sections \ref{sec:measurement} and \ref{sec:inversion} analyze the post-failure distribution of the sum of squared counts and give the rule for computing $U$.

\section{Unbiased estimation after failure and the Gaussian upper bound}\label{sec:measurement}
We use the independence of signs and mapping established in Section~\ref{sec:design} to analyze $\widehat d$ after decoding failure and construct an upper bound for $d$. Cancellation between random signs of different keys makes $\widehat d$ exactly unbiased for every fixed mapping (Proposition~\ref{prop:moments}). It therefore remains unbiased even after arbitrarily rare failures. Controlling the conditional miss rate at the target risk $u$ also requires the distribution of post-failure fluctuations of $\widehat d$ around $d$. When $d\gg\sqrt M$, the standardized estimation error converges to the standard Gaussian distribution (Theorem~\ref{thm:clt}), yielding the Gaussian rule (Definition~\ref{def:gauss}).

\subsection{Conditional unbiasedness and conditional variance given the mapping}
We expand the sum of squares and the total sum in~\eqref{eq:dhat} by key. Each differing key occupies $k$ cells and satisfies $\varepsilon_i^2=1$, so the diagonal contributions of the two terms sum to exactly $d$. The estimation error consists of the cross terms between different keys:
\begin{equation}
\widehat d-d=\frac{2}{k(1-k/M)}\sum_{i<j}\varepsilon_i\varepsilon_j\Bigl(H_{ij}-\frac{k^2}{M}\Bigr).\label{eq:rs-expand}
\end{equation}
The right-hand side is a quadratic form in the random signs. Each coefficient $H_{ij}-k^2/M$ measures the overlap relative to its unconditional mean and depends entirely on the mapping $\mathcal G$. For the standard IBLT, the same expansion replaces the sign products with fixed values $s_is_j$. Selection of overlaps by failure then translates directly into bias (Proposition~\ref{prop:projection}). With random signs, each $\varepsilon_i\varepsilon_j$ takes $\pm1$ with equal probability even on a fixed mapping where failure occurs (Proposition~\ref{prop:composition}). The cross terms therefore have mean zero on every mapping. Intuitively, failure selects mappings with large overlap and amplifies the cross-term coefficients, while each term retains equally likely signs. The error remains centered at zero, and failure changes only its scale. Proposition~\ref{prop:moments} expresses this intuition through the conditional mean and variance given the mapping.

\begin{proposition}[Conditional unbiasedness and conditional moments given the mapping]\label{prop:moments}
For every fixed mapping $\mathcal G$,
\begin{equation}
\E[\widehat d\mid\mathcal G]=d,\qquad
\Var(\widehat d\mid\mathcal G)=\frac{4}{k^2(1-k/M)^2}\sum_{i<j}\Bigl(H_{ij}-\frac{k^2}{M}\Bigr)^2\le2d(d-1).\label{eq:condmoments}
\end{equation}
Thus, for any mapping event $\mathcal S$ of positive probability, $\E[\widehat d\mid\mathcal S]=d$. In particular, $\E[\widehat d\mid F]=d$ holds for every finite $M$ and every sign composition, even for arbitrarily rare failures. Averaging over the random mapping $\mathcal G$ gives $\E\widehat d=d$ and $\Var\widehat d=s(d)^2$.
\end{proposition}

\begin{proof}[Proof sketch]
Given $\mathcal G$, $(\varepsilon_i)$ are iid uniform signs. Taking expectations term by term in~\eqref{eq:rs-expand} and using $\E[\varepsilon_i\varepsilon_j\mid\mathcal G]=0$ proves conditional unbiasedness. The sign products for distinct unordered pairs $(i,j)$ are pairwise uncorrelated, yielding the conditional variance in~\eqref{eq:condmoments}. The Cauchy--Schwarz inequality gives the upper bound $2d(d-1)$. The full proof is in Appendix~\ref{app:proof-moments}.
\end{proof}

For a fixed mapping, $\widehat d$ is a Rademacher quadratic form with the Gram matrix of the centered incidence vectors. Its conditional mean equals the trace, and its conditional variance comes only from the off-diagonal terms. These are the same identities used by the Hutchinson trace estimator\cite{hutchinson1990trace}. Proposition~\ref{prop:moments} expresses them in terms of overlap coefficients, establishing conditional unbiasedness for any mapping event.

Figure~\ref{fig:composition} illustrates this unbiasedness numerically. The post-failure bias of the standard encoding changes direction with $k$, while the mean for the random-sign encoding stays near zero. Section~\ref{sec:exp-decode} and Appendix~\ref{app:dist} compare conditional moments across additional configurations.

\subsection{\texorpdfstring{The Gaussian limit after failure when $d\gg\sqrt M$}{The Gaussian limit after failure when d >> sqrt(M)}}
Proposition~\ref{prop:moments} establishes the center of the estimator after failure. Constructing an upper bound also requires the shape of the error distribution around this center. Failure affects $\widehat d$ only through the mapping. By~\eqref{eq:condmoments}, the sum of squared overlaps $\sum_{i<j}(H_{ij}-k^2/M)^2$ determines the conditional error scale. Failure selects mappings with different values of this sum and may therefore change the error distribution. We standardize the error by the unconditional standard deviation $s(d)$:
\begin{equation}
X=\frac{\widehat d-d}{s(d)},\label{eq:X}
\end{equation}
and study the conditional distribution of $X$ given $F$.

We first fix a mapping. By~\eqref{eq:rs-expand}, the error is a weighted sum of sign products $\varepsilon_i\varepsilon_j$. The main contributions come from key pairs sharing a cell: the intersecting edge pairs with $H_{ij}\ge1$ in $\mathcal G$. Their expected number is about $k^2d^2/(2M)$. Sign products for distinct key pairs are pairwise uncorrelated, and no small set of terms dominates the sum. The number of intersecting edge pairs grows with $d^2/M$ and tends to infinity when $d\gg\sqrt M$. The error then sums many $\pm1$ perturbations, and its distribution approaches a Gaussian. When $d$ is of the same order as $\sqrt M$, only finitely many intersecting edge pairs contribute. The error consists of a few $\pm1$ terms, and its distribution remains discrete. We distinguish these regimes by the scale parameter
\begin{equation}
c=\frac{d}{\sqrt M}.\label{eq:c}
\end{equation}
The expected number of intersecting edge pairs is about $k^2c^2/2$, so $c$ plays the role of the collision scale in the birthday paradox\cite{arratia1989chen,dw}. It and the load $d/M$ are independent parameters.

Conditioning on failure preserves the sign structure. On any mapping where failure occurs, $(\varepsilon_i)$ remain iid uniform signs (Proposition~\ref{prop:composition}). The Gaussian shape persists as long as the mapping still contains many intersecting edge pairs. This also holds for rare failures. Under a subcritical load, a typical failure is caused by a duplicate edge pair. This raises a single $H_{ij}$ to $k$ and leaves the overall structure of intersecting edge pairs almost unchanged.

\begin{theorem}[Gaussian limit after failure]\label{thm:clt}
Fix $k\ge3$ and let $d/\sqrt M\to\infty$, allowing any limiting behavior of $d/M$. Then
\begin{equation}
\sup_x\bigl|\Pr(X\le x\mid F)-\Phi(x)\bigr|\longrightarrow0.\label{eq:clt}
\end{equation}
This includes the critical window, rare failures under a fixed subcritical load, and arbitrary overload.
\end{theorem}

\begin{proof}[Proof sketch]
The proof has two steps. First, we fix the mapping and take randomness only over the signs. Suppose the number of intersecting edge pairs is of the same order as $M$ and each key participates in overlaps of size at most $O(\log M)$. The higher cumulants of the standardized quadratic form~\eqref{eq:rs-expand} then tend to zero, and the conditional distribution of $X$ converges to $\mathcal N(0,1)$ uniformly over such mappings. Second, we show that, conditional on $F$, the mapping satisfies these conditions with probability tending to $1$. We consider three cases according to the limit of $d/M$. When $d/M$ has a positive finite limit, $\Pr(F)\ge\binom Mk^{-1}$. The probability that a mapping violates the conditions is superpolynomially small and remains negligible after conditioning. When $d=o(M)$, Proposition~\ref{prop:plant} shows that the mapping distribution given failure is asymptotically identical to the planted distribution in TV distance. The Gaussian limit holds for the planted distribution. When $d>M$, failure occurs with certainty, so the conditional and unconditional distributions coincide. The full proof is in Appendix~\ref{app:clt}.
\end{proof}

Figure~\ref{fig:encoding} compares the distributions of $X$ for failure samples and all samples in a configuration with $c\approx104$. For the standard encoding, we compute $X$ in the same way from its counts $C_v=\sum_{i:v\in S_i}s_i$. With random-sign encoding, the two distributions almost coincide. With standard encoding, the failure distribution is shifted as a whole.

\subsection{The Gaussian rule and its conditional miss rate}\label{sec:gauss}
By Theorem~\ref{thm:clt}, $\widehat d$ is approximately Gaussian after failure, with mean $d$ and standard deviation $s(d)$. If $d$ were known, we could set the margin using the fact that $\widehat d$ falls below $d+s(d)\Phi^{-1}(u)$ with probability about $u$. However, the scale $s(d)$ (\eqref{eq:s}) depends on the unknown $d$. Direct substitution of $s(\widehat d)$ would make the margin fluctuate with the observation. A low $\widehat d$ gives a small $s(\widehat d)$, reducing the margin precisely for samples that need a larger one. The risk analysis would then need to account for this additional error.

Test inversion is a classical method for constructing confidence sets\cite{neyman1937}. We reverse the reasoning and test candidate values individually. For each candidate $n$, we assume the true value is $n$. The scale $s(n)$ is then known, and $\widehat d$ falls below $n+s(n)\Phi^{-1}(u)$ with probability about $u$. An observation below this lower bound indicates that $n$ is too large, so we exclude $n$. The largest remaining candidate gives the upper bound $U$. The risk is controlled by a single event: $U<d$ implies that the true value $d$ is excluded. Under the true value $d$, this exclusion has probability about $u$. Definition~\ref{def:gauss} follows this construction, replacing $\Phi^{-1}(u)$ with a more accurate quantile for finite $M$.

\begin{definition}[Gaussian rule]\label{def:gauss}
Given an observation $\widehat d$ and a target risk $u\in(0,1/2)$, we test each candidate integer $n\ge2$ as the true value, for which the scale $s(n)$ is known. Let $\chi^2_{M-1,u}$ denote the lower $u$ quantile of the chi-square distribution with $M-1$ degrees of freedom. Set the acceptance lower bound to
\begin{equation}
t(n)=n+s(n)\,\frac{\chi^2_{M-1,u}-(M-1)}{\sqrt{2(M-1)}}.\label{eq:gauss}
\end{equation}
We accept candidate $n$ if $\widehat d\ge t(n)$ and reject it otherwise. We set $U$ to the largest accepted candidate, returning $+\infty$ if none is accepted.
\end{definition}

For fixed $M,k$ as $d\to\infty$, $\widehat d/d$ converges in distribution to $\chi^2(M-1)/(M-1)$\cite{wu2026selfsizing}, which has mean 1. After centering the chi-square variable at its mean and scaling it to unit variance, its lower $u$ quantile is $(\chi^2_{M-1,u}-(M-1))/\sqrt{2(M-1)}$. Equation~\eqref{eq:gauss} uses this quantile in place of $\Phi^{-1}(u)$ to capture skewness more accurately for finite $M$. The two quantiles agree as $M\to\infty$. In the experimental range $M\ge64$, $u\ge0.01$, $t(n)$ is strictly increasing for $n\ge2$ (Appendix~\ref{app:gauss-mono}), so we can find $U$ by monotone binary search. Outside this range, we check candidates individually as specified in the definition. The rule depends only on the observation $\widehat d$ and the known parameters $(M,k,u)$.

\begin{corollary}[Conditional miss rate of the Gaussian rule]\label{cor:gauss}
Under the conditions of Theorem~\ref{thm:clt}, for any fixed $u\in(0,1/2)$, $\limsup\Pr(U<d\mid F)\le u$.
\end{corollary}

\begin{proof}[Proof sketch]
We take $n=d$. Theorem~\ref{thm:clt} and $(t(d)-d)/s(d)\to\Phi^{-1}(u)$ give a rejection probability of $u+o(1)$ for the true candidate. The conclusion follows from $\{U<d\}\subseteq\{\widehat d<t(d)\}$. The full proof is in Appendix~\ref{app:clt}.
\end{proof}

By Corollary~\ref{cor:gauss}, the Gaussian rule has conditional miss rate asymptotically at most $u$ throughout the range $d\gg\sqrt M$ of Theorem~\ref{thm:clt}. This range includes the critical window and the overload region. Figure~\ref{fig:coverage} shows that the Gaussian rule's miss/u approaches the target level as $c$ grows. When $d$ is of the same order as $\sqrt M$ or smaller, the Gaussian approximation no longer holds, and the actual conditional miss rate exceeds $u$. Section~\ref{sec:inversion} explains this behavior and gives the Skellam rule, whose conditional miss rate is asymptotically at most $u$ uniformly over all $d\ge2$.

\section{The Skellam rule for all loads}\label{sec:inversion}
The Gaussian rule's conditional miss rate guarantee requires $d\gg\sqrt M$. If we size the IBLT table for the expected load, $d$ can be far below $\sqrt M$. Failure is then rare, but the two-round protocol must still use the post-failure observation to size the second-round table. In this range, the conditional miss rate of the Gaussian rule can be several times the target $u$ (Table~\ref{tab:gauss-c} in the appendix).

The reason lies in the shape of the post-failure error at low load. Failure then has essentially one source: two differing keys select exactly the same $k$ cells. These cells contain only these two keys. The table has no pure cell, so peeling cannot start. Under random signs, the two keys are equally likely to have the same or opposite signs. If we place each key in the table alone, it leaves a count of $\pm1$ in each of its $k$ cells. Together, their squared counts sum to $2k$. Dividing by $k$ gives their exact contribution of $2$. With opposite signs, the counts in all $k$ shared cells cancel to $0$. The entire sum of squares of $2k$ disappears, leaving the estimate $2$ below the true value. With equal signs, each shared cell has count $\pm2$ and squared count $4$. The sum of squares becomes $4k$, leaving the estimate $2$ above the true value. This jump is equally likely to be upward or downward, and its size is independent of $M$. The remaining keys also occasionally share a cell, almost always exactly one. Each such overlap shifts the estimate by $2/k$, again with either direction equally likely. The number of such key pairs is approximately Poisson\cite{arratia1989chen,barbour1992poisson} and remains finite when $c$ is bounded. The post-failure error thus consists of one large jump of size $2$ and several small jumps of size $2/k$.

The Gaussian approximation cannot capture these two types of jumps. When $d$ is small relative to $\sqrt M$, the standard deviation $s(d)$ of the Gaussian rule tends to zero. The downward jump of size $2$ remains unchanged and lies arbitrarily many standard deviations into the Gaussian tail. The probability of rejecting the true value therefore tends to one half. When $c$ is of order one, only finitely many small jumps occur. The error still has a lattice distribution, with a visible departure from the Gaussian in the lower tail corresponding to the target risk (Appendix~\ref{app:limit}). We construct a null distribution that retains both types of jumps and approaches the Gaussian used by the Gaussian rule as $c$ grows (Figure~\ref{fig:aux}). The resulting Skellam rule controls the conditional miss rate for all $d\ge2$. Appendices~\ref{app:inversion} and~\ref{app:search} give the full proof and numerical implementation.

\subsection{The lattice-preserving test statistic \texorpdfstring{$\widetilde d$}{d~}}\label{sec:raw}
To keep these jumps exactly on lattice points, we need a discrete test statistic. The correction term $(\sum_vC_v)^2/M$ in $\widehat d$ shifts its values off the lattice. We therefore use the normalized sum of squares without subtracting the mean term,
\begin{equation}
\widetilde d=\frac1k\sum_{v=1}^M C_v^2.\label{eq:raw-energy}
\end{equation}
Expanding the sum of squares gives $d$ from the diagonal terms, leaving only overlap contributions to the error:
\begin{equation}
\widetilde d-d=\frac2k\sum_{i<j}H_{ij}\,\varepsilon_i\varepsilon_j.\label{eq:raw-expand}
\end{equation}
Every term on the right is an integer multiple of $2/k$, so $\widetilde d-d$ takes only integer multiples of $2/k$. The expansion shows both types of jumps directly. A duplicate edge pair has $H_{ij}=k$ and contributes exactly $2\varepsilon_i\varepsilon_j$, or $\pm2$ with equal probability. An edge pair sharing a single cell has $H_{ij}=1$ and contributes $\pm2/k$.

Without the correction term, $\widetilde d$ remains exactly unbiased for every mapping.

\begin{proposition}[Conditional unbiasedness of $\widetilde d$]\label{prop:raw-unbiased}
For every fixed mapping $\mathcal G$, $\E[\widetilde d\mid\mathcal G]=d$. Thus, for any mapping event $\mathcal S$ with positive probability, $\E[\widetilde d\mid\mathcal S]=d$. In particular, $\E[\widetilde d\mid F]=d$ holds for all finite $M$ and all sign compositions, however rare the failure.
\end{proposition}

\begin{proof}
Given $\mathcal G$, $H_{ij}$ is constant. By Proposition~\ref{prop:composition}, the signs remain iid and uniform, so $\E[\varepsilon_i\varepsilon_j\mid\mathcal G]=0$ for $i\ne j$. Taking expectations term by term in~\eqref{eq:raw-expand} gives the result.
\end{proof}

This property also follows from the random signs. In the standard IBLT, the cross terms have fixed sign products $s_is_j$. The same expansion gives the unconditional mean of the sum of squares as $d+\frac kM(b^2-d)$, which depends on the unknown $b$. The correction term in $\widehat d$ removes this dependence. With random signs, the cross terms have mean zero for every mapping. The sum of squares is therefore unbiased without the correction term.

$\widetilde d$ and $\widehat d$ differ only by the squared sum of the counts, weighted by $k/M$. This correction depends only on the random signs and is the classical AMS second-moment estimator\cite{alon1999spacecomplexity} (Appendix~\ref{app:ams}). After standardization by $s(d)$, the difference tends to zero, so Theorem~\ref{thm:clt} also holds for $\widetilde d$ (Appendix~\ref{app:clt}).

\subsection{Constructing the null distribution from the two jumps}\label{sec:aux}
For each candidate $n$, we construct a reference distribution for $\widetilde d$ after failure when the true value is $n$. This is the null distribution, which we use in Section~\ref{sec:skellam} to test each candidate. We construct it by counting the occurrences of the two types of jumps.

The large jump occurs exactly once and contributes $2\varepsilon_1\varepsilon_2$. Its direction is determined by the sign product of the duplicate edge pair.

The small jumps come from the $n-2$ keys outside the duplicate edge pair. They form $\binom{n-2}2$ pairs. Each pair shares a cell with probability approximately $k^2/M$, producing a small jump. Upward and downward jumps are equally likely, so each has expected count
\begin{equation}
\lambda_n=\frac12\binom{n-2}2\frac{k^2}M=\frac{k^2(n-2)(n-3)}{4M}.\label{eq:lambda}
\end{equation}
Key pairs sharing a cell are sparse. By Poisson splitting\cite{kingman1993poisson}, the upward and downward counts are approximately independent $\Pois(\lambda_n)$ counts. The upward count minus the downward count gives the net number of steps and follows a Skellam distribution\cite{skellam1946}, denoted $Z_n\sim\Skellam(\lambda_n,\lambda_n)$. Each step has size $2/k$.

Adding both types of jumps to the candidate $n$ gives the null distribution.
\begin{definition}[null distribution]\label{def:null}
For $\lambda\ge0$, let
\begin{equation}
J_\lambda=2\varepsilon_1\varepsilon_2+\frac2kZ,\qquad Z\sim\Skellam(\lambda,\lambda),\label{eq:J}
\end{equation}
where $\varepsilon_1\varepsilon_2$ takes the values $\pm1$ with equal probability and is independent of $Z$. We call $2\varepsilon_1\varepsilon_2$ the \textit{$\pm2$ term} and $\frac2kZ$ the \textit{Skellam term}. For a candidate integer $n\ge2$, we take $\lambda_n$ from~\eqref{eq:lambda}. The null distribution for candidate $n$ is the distribution of
\begin{equation}
n+J_{\lambda_n}=n+2\varepsilon_1\varepsilon_2+\frac2kZ_n\label{eq:aux}
\end{equation}
with CDF denoted by $G_n$.
\end{definition}
In Section~\ref{sec:uniform}, we prove that the post-failure distribution of $\widetilde d$ approaches $G_d$ uniformly over all $d\ge2$ (Theorems~\ref{thm:cross} and~\ref{thm:uniform}).

When $n=2,3$, fewer than two keys remain, so $\lambda_n=0$ and the null distribution assigns probability one half to each of $n\pm2$. In general, its mean is $n$. Since $\Var Z_n=2\lambda_n$, its variance is $4+2(n-2)(n-3)/M$.

\begin{figure}[!htbp]
\centering\includegraphics[width=\linewidth]{../figure/fig-aux.pdf}
\caption{Null distributions at three values of $c=d/\sqrt M$ ($k=3$, with $\lambda_d$ set to its limit $k^2c^2/4$). Blue and orange correspond to opposite and equal signs of the duplicate edge pair. The dashed line is the Gaussian used by the Gaussian rule.}
\label{fig:aux}
\end{figure}

The shape of this family explains why it applies at all loads. Figure~\ref{fig:aux} shows how the null distribution varies with $c$. When $c$ is small, the Skellam term is almost zero and the distribution is concentrated at $\pm2$. The Gaussian used by the Gaussian rule concentrates between these points, which lie beyond its tails. When $c\approx1$, the opposite-sign and equal-sign cases each spread into a lattice distribution with spacing $2/k$. These distributions overlap, but their combined shape still differs visibly from the Gaussian. As $c$ grows further, the width of the Skellam term far exceeds $2$, and the $\pm2$ term becomes a small contribution. The null distribution then approaches the Gaussian used by the Gaussian rule. This family thus connects the small-$c$ region to the region where the Gaussian rule applies. Section~\ref{sec:uniform} gives the rigorous approximation guarantee.

\subsection{The null distribution approximates the post-failure distribution at all loads}\label{sec:uniform}
Definition~\ref{def:null} constructs the null distribution from a counting argument. To use it to test candidates, we establish that it describes the post-failure distribution of $\widetilde d$ for all $d$. The receiver does not know the load region containing $d$ and cannot use the load to select a rule. The same family $G_d$ must therefore apply for both small and large $d$. We analyze these two cases and then combine them.

When $d$ is of the same order as $\sqrt M$ or smaller, overlaps are sparse. The model underlying $G_d$ consists of a duplicate edge pair and scattered single-cell overlaps among the remaining keys. The actual mapping may have additional structure: the duplicate edge pair may share a cell with a remaining key, two keys may share more than one cell, or one key may share cells with two other keys. Each case is much rarer than a single-cell overlap, and its probability tends to zero as $M$ grows. After excluding these cases, the actual mapping has exactly the structure underlying $G_d$. The actual distribution is therefore close to $G_d$ at every lattice point (Theorem~\ref{thm:cross}), as illustrated by the first two shapes in Figure~\ref{fig:aux}.

When $d\gg\sqrt M$, intersecting edge pairs are numerous and interwoven, so the sparse structure above no longer applies. The null distribution $G_d$ still applies because both distributions approach the same Gaussian. The statistic $\widetilde d$ is a sum of many small terms, each with equally likely signs, and approaches a Gaussian by Theorem~\ref{thm:clt}. The null distribution $G_d$ also approaches a Gaussian, as the third shape in Figure~\ref{fig:aux} shows, with width matching $s(d)$. The two distributions differ in detail but share the same Gaussian limit. Theorem~\ref{thm:uniform} combines these cases to give a uniform approximation over all $d\ge2$.

\begin{theorem}[Discrete limit distribution when $c$ is bounded]\label{thm:cross}
Fix $k\ge3$ and a finite constant $C$.
\begin{enumerate}[label=(\roman*)]
\item For $2\le d\le C\sqrt M$, the total variation distance between the conditional distribution of $\widetilde d$ given $F$ and the null distribution $G_d$ tends to zero uniformly and is at most
\begin{equation}
O_k\Bigl(\frac dM+\frac{d^3}{M^2}\Bigr).\label{eq:tv-error}
\end{equation}
\item If $d/\sqrt M\to c\in(0,\infty)$, the conditional distribution of $\widetilde d-d$ given $F$ converges in total variation to the distribution of $J_{k^2c^2/4}$, whose variance is $4+2c^2$.
\item The limiting distributions in (ii) satisfy
\begin{equation*}
J_{k^2c^2/4}\Rightarrow2\varepsilon_1\varepsilon_2\quad(c\to0),\qquad\frac{J_{k^2c^2/4}}{\sqrt2\,c}\Rightarrow\mathcal N(0,1)\quad(c\to\infty).
\end{equation*}
\end{enumerate}
\end{theorem}

Figure~\ref{fig:crosscdf} compares the lower tail of this discrete limit distribution with the Gaussian approximation, using planted samples with $c$ ranging from 0.25 to 8. Section~\ref{sec:exp-dist} also reports the distance between the empirical and limit distributions at fixed $c$ as $M$ grows.

\begin{proof}[Proof sketch]
We prove (i) in two steps. First, we replace the actual failure with a duplicate edge pair. By Proposition~\ref{prop:plant}, we can replace the mapping conditional on failure with the planted distribution, incurring error $O_k(d/M^2+M^{1-k})$ in the range of the theorem (Lemma~\ref{lem:plant}). Second, we exclude the three cases above from the planted mapping. Their probabilities are $O_k(d/M)$, $O_k(d^2/M^2)$, and $O_k(d^3/M^2)$, respectively. The remaining structure consists of a duplicate edge pair and mutually disjoint single-cell overlaps. The duplicate edge pair gives the $\pm2$ term. By the Chen--Stein local dependence bound\cite[Theorem~1]{arratia1989chen}, the overlap count is approximately $\Pois(2\lambda_d)$. Independent signs split this count into two independent $\Pois(\lambda_d)$ counts. Their difference gives the Skellam term, so $\widetilde d-d$ is close to $J_{\lambda_d}$. The Poisson approximation has error $O_k((d^2+d^3)/M^2)$. Since $2\le d\le M$ and $k\ge3$, each of $d/M^2$, $d^2/M^2$, and $M^{1-k}$ is at most $d/M$. Also, $(d^2+d^3)/M^2$ is at most $2d^3/M^2$. Combining these bounds gives~\eqref{eq:tv-error}.
(ii) By (i), it suffices to show that $J_{\lambda_d}$ converges in total variation to $J_{k^2c^2/4}$. When $d/\sqrt M\to c$, we have $\lambda_d\to k^2c^2/4$. The required convergence follows from continuity in the Poisson parameter. Independence of the two terms gives $\Var J_\lambda=4+8\lambda/k^2$. Substituting $\lambda=k^2c^2/4$ gives $4+2c^2$.
(iii) As $c\to0$, $\lambda\to0$ and the Skellam term tends to zero in probability, leaving only the $\pm2$ term of $J_\lambda$. As $c\to\infty$, $\Var J_{k^2c^2/4}\sim2c^2$ and the Skellam term satisfies the CLT. The $\pm2$ term divided by $\sqrt2\,c$ tends to zero. Appendix~\ref{app:inversion} gives the full argument.
\end{proof}

\begin{theorem}[Uniform CDF approximation]\label{thm:uniform}
Fix $k\ge3$. The Kolmogorov distance between the conditional distribution of $\widetilde d$ given $F$ and the null distribution $G_d$ satisfies
\begin{equation}
\sup_{d\ge2}\ \sup_{x\in\mathbb R}\bigl|\Pr(\widetilde d\le x\mid F)-G_d(x)\bigr|\longrightarrow0
\quad(M\to\infty).\label{eq:uniform}
\end{equation}
\end{theorem}

\begin{proof}[Proof sketch]
When $d\gg\sqrt M$, Gaussian limits let us handle a dense intersection graph without extending the matching structure. By Theorem~\ref{thm:clt}, the conditional distribution of $\widetilde d$, standardized by $s(d)$, approaches the standard Gaussian. The null distribution is that of $d+J_{\lambda_d}$. When $\lambda_d\to\infty$, the Skellam term of $J_{\lambda_d}$ satisfies the CLT. The ratio of $\Var J_{\lambda_d}=4+2(d-2)(d-3)/M$ to $s(d)^2$ tends to one, and the $\pm2$ term divided by $s(d)$ tends to zero. Both distributions approach the same continuous Gaussian CDF, so their Kolmogorov distance tends to zero. This argument applies to $d=o(M)$, a positive load ratio, and unbounded overload. We obtain uniformity by a subsequence argument. Take a subsequence along which $d/\sqrt M$ converges. Theorem~\ref{thm:cross} handles finite limits, and the Gaussian argument above handles infinite limits (Appendix~\ref{app:inversion}).
\end{proof}

The same family of null distributions $G_d$ therefore describes the post-failure distribution of $\widetilde d$ for all $d\ge2$. For small $d$, the approximation holds pointwise on the lattice. For large $d$, both distributions approach the same Gaussian. We can therefore use this family without knowing the load region. In Section~\ref{sec:skellam}, we use it to construct the Skellam rule.

\subsection{The Skellam rule controls the conditional miss rate at all loads}\label{sec:skellam}
The Skellam rule uses the test inversion of the Gaussian rule (Section~\ref{sec:gauss}), replacing the Gaussian lower tail with the lattice lower tail of $G_n$. If the observation falls in the lower $u$ tail of $G_n$, candidate $n$ is too large to be the true value, so we reject candidate $n$. We take $U$ to be the largest accepted candidate.

\begin{definition}[Skellam rule]\label{def:skellam}
Let $\widetilde d$ be the observation and $u\in(0,1/2)$ the target risk. For each candidate integer $n\ge2$, we accept $n$ if $G_n(\widetilde d)>u$ and reject it otherwise. We define the upper bound as the supremum of all accepted candidates,
\begin{equation}
U=\sup\bigl\{n\ge2:\ G_n(\widetilde d)>u\bigr\},\label{eq:skellam-U}
\end{equation}
with $U=+\infty$ when the accepted candidates are unbounded.
\end{definition}

The null distribution for candidate $2$ takes only the values $0$ and $4$. Since $\widetilde d\ge0$, we always accept candidate $2$, so $U\ge2$. For practical table sizes, $U$ is also finite (for $u=0.01$, $M>18$ suffices). We can compute it exactly by a finite search, at a cost roughly proportional to $\lambda_d$ per computation (Appendix~\ref{app:search}).

The behavior at small and large loads explains why the Skellam rule applies throughout. When $c$ is small, the Skellam term is almost zero. The null distribution $G_n$ assigns half its mass to $n-2$ and half to $n+2$. We accept candidate $n$ exactly when the smaller value $n-2$ does not exceed the observation, giving $U\approx\widetilde d+2$. The rule thus adds back a potentially lost large jump. When the duplicate edge pair has opposite signs, the observation is $d-2$ and $U$ equals $d$ exactly. With equal signs, the observation is $d+2$ and $U$ is about $d+4$. Neither case produces a miss. In this region, the scale $s(n)$ of the Gaussian rule tends to zero, so its upper bound $U$ stays close to the observation. It underestimates $d$ when the duplicate edge pair has opposite signs, giving a conditional miss rate that tends to one half. When $c$ is large, the Skellam term dominates and $G_n$ approaches a Gaussian. The Skellam rule then approaches the Gaussian rule.

Theorem~\ref{thm:coverage} gives the main result of this section. For sufficiently large tables, whatever load region contains $d$, the probability that the Skellam rule underestimates $d$ after failure is at most the target risk $u$. This guarantee holds for any sign composition. The rule uses only $\sum_vC_v^2$ from the same table and the known parameters $(M,k,u)$. The Gaussian rule provides this guarantee only when $d\gg\sqrt M$ (Corollary~\ref{cor:gauss}). In Section~\ref{sec:protocol}, we apply this guarantee to the risk decomposition of the two-round protocol.

\begin{theorem}[Uniform asymptotic upper bound on the conditional miss rate]\label{thm:coverage}
Fix $k\ge3$ and $u\in(0,1/2)$. For any sign composition, the Skellam rule satisfies
\begin{equation}
\limsup_{M\to\infty}\ \sup_{d\ge2}\ \Pr(U<d\mid F)\le u.\label{eq:coverage}
\end{equation}
\end{theorem}

\begin{proof}[Proof sketch]
If the observation has distribution $G_n$ exactly, we reject $n$ when it falls in a lower tail with total probability at most $u$. By the definition of the CDF, the rejection probability is at most $u$, including for discrete distributions. The rejection region is a lower set, so the rejection probabilities under the actual distribution and $G_n$ differ by at most their Kolmogorov distance. A miss implies rejection of the true value: $\{U<d\}\subseteq\{G_d(\widetilde d)\le u\}$. Taking $n=d$ and applying Theorem~\ref{thm:uniform} gives~\eqref{eq:coverage}. By Proposition~\ref{prop:composition}, the joint distribution of the coding coefficients is independent of the sign composition. The result therefore holds for any sign composition. Appendix~\ref{app:inversion} gives the full proof.
\end{proof}

Theorem~\ref{thm:coverage} is asymptotic in the table size. Two effects arise on finite tables. First, $G_n$ is a lattice distribution, so the rejection region contains only whole probability masses. The actual conditional miss rate can therefore be strictly below $u$. Second, the null distribution omits overlaps between the duplicate edge pair and the remaining keys. These occur at a rate of about $(d-2)k^2/M$ and can push the conditional miss rate slightly above $u$ on small tables. Section~\ref{sec:exp-cover} evaluates finite-table behavior experimentally (Figures~\ref{fig:coverage} and~\ref{fig:confirm}). At $M=64$, the largest miss/u is 1.08, an excess beyond sampling error. All configurations with $M\ge256$ meet the target within sampling error. The effect of the omitted overlaps decreases as $1/M$, and extra capacity also decreases as $M$ grows. When the first-round table must be small, the second-round sizing margin $1/(0.8\alpha_c)$ absorbs slight underestimation (Section~\ref{sec:protocol}). If more margin is needed, we can reduce the target load from 0.8, using a larger second-round table to lower the two-round failure rate.

\subsection{Combining the two rules using the first-round observation}\label{sec:switch}
The two rules have limitations in different ranges. The Gaussian rule underestimates $d$ only when $c$ is bounded, where the Skellam rule is inexpensive to compute. The cost of the Skellam rule grows roughly with $\lambda_d$, or $c^2$, while the Gaussian rule meets the target for large $c$. A system can therefore combine the rules using the first-round observation. We use the Gaussian rule when $\widetilde d$ is at or above a switch point and the Skellam rule otherwise. We compute $\widetilde d$ directly from the first-round counts, so switching adds no cost.

The combined rule retains the uniform guarantee of Theorem~\ref{thm:coverage} if the switch point grows faster than $\sqrt{M_1}$ as $M_1$ increases (Appendix~\ref{app:switch}). When $d/\sqrt{M_1}$ is bounded, $\widetilde d$ falls below the switch point with probability tending to one, so we use the Skellam rule. When $d/\sqrt{M_1}\to\infty$, both rules compare an asymptotically Gaussian statistic with the same lower $u$ quantile and give asymptotically identical conclusions. For example, a switch point of $1.5\alpha_cM_1$ corresponds to a first-round load of 1.5. Above load 1.5, sizing requires only one monotone binary search, and the protocol takes the same time in the overload region as with the Gaussian rule (Section~\ref{sec:exp-cost}).

\subsection{From the conditional miss rate to the failure probability of the two-round protocol}\label{sec:protocol}
We now connect the post-failure conditional miss rate to the overall success rate of two-round reconciliation. The protocol follows Section~\ref{sec:flow}. We denote first-round failure by $F_1$ and second-round failure by $F_2$. Let $E$ denote a checksum error: a non-pure cell passes the check, a key that passes the check lies outside the set difference, or all fields are zero at termination despite a nonempty core.

\begin{theorem}[Two-round risk decomposition]\label{thm:protocol}
Under the ideal hash model, let $\Pr_m(F)$ denote the probability of peeling failure for $d$ differing keys in a fresh table with $m$ cells. If the reserve size is at least $m(d)$, then
\begin{equation}
\Pr\big((F_1\cap F_2)\cup E\big)
\le\Pr(F_1,\ U_b<d)+\Pr(F_1)\max_{m\ge m(d)}\Pr\nolimits_m(F)+\Pr(E).\label{eq:risk}
\end{equation}
\end{theorem}

\begin{proof}
By the union bound, it suffices to bound $F_1\cap F_2$ and $E$ separately. The upper bound $U_b$ is a function of the first-round counts and the set sizes. The second-round mapping is independent of it. Under ideal purity, $F_2$ depends only on the second-round mapping and $d$. Hence $\Pr(F_1\cap F_2)=\E[\mathbf 1_{F_1}\Pr_{m(U_b)}(F)]$. On $\{U_b<d\}$, we bound the conditional failure probability by one. On $\{U_b\ge d\}$, we have $m(U_b)\ge m(d)$. Taking the supremum gives~\eqref{eq:risk}.
\end{proof}

The three terms correspond to three sources of error. The first is a sizing error: after first-round failure, the upper bound underestimates $d$. The second-round table is sized using $U_b$ and has insufficient capacity. This term is the joint miss rate $\Pr(F_1)\Pr(U_b<d\mid F_1)$. For the Skellam rule, Theorem~\ref{thm:coverage} bounds it asymptotically by $u\Pr(F_1)$, uniformly over all $d\ge2$. For a one-sided difference, $U_b\ge|b|=d$, so this term is zero. The Gaussian rule can have a conditional miss rate far above $u$ when $c$ is bounded, although failure itself is then rare. By Proposition~\ref{prop:plant}, when $d=o(M_1)$, $\Pr(F_1)$ is asymptotically equal to the expected number of duplicate edge pairs, $\binom d2/\binom{M_1}k$. This is about $2.8\times10^{-6}$ for $k=3$, $M_1=1024$, and $d=\sqrt{M_1}=32$, so the joint miss rate remains small.

The second term comes from the randomness of peeling. Even with sufficient second-round capacity, the random hypergraph can leave a nonempty 2-core with small probability. The second-round load is at most $0.8\alpha_c$, so Lemma~\ref{lem:propL} gives a failure probability that decays polynomially in $d$. The third term is a verification error: the checksum identifies a non-pure cell as pure or indicates termination on a nonempty core, leading to incorrect recovery. Each check passes falsely with probability at most $2^{-64}$. The number of checks is at most the number of cells plus $k$ times the number of deleted keys. The keySum and count fields of a nonempty core can both cancel to zero. The probability that all checksums are then also zero is at most $2^{-64}$, and this termination error can occur at most once per round. Thus $\Pr(E)$ is of order $2(M_1+m+kd)\cdot2^{-64}$, which is negligible at practical scales. With 32-bit checksums, this term is of order $2(M_1+m+kd)\cdot2^{-32}$, at most $10^{-4}$ for $M_1=16384$ and load at most 3. Since $\widetilde d$ uses only the count field, checksum width does not affect the first term. At deployment, we allocate the target failure probability between the sizing and checksum terms.

The bound in~\eqref{eq:risk} adds the three terms and conservatively treats every miss in the first term as a two-round failure. The second-round table has spare capacity, so it still succeeds with high probability under slight underestimation. The actual contribution of such misses is on the same scale as the second term. The protocol experiments in Section~\ref{sec:exp-e2e} also show that two-round failures arise almost entirely from the second term. Both rules have two-round failure rates comparable to that of the oracle upper bound ($U=d$).

\section{Experimental evaluation}\label{sec:exp}
The experiments answer three questions. First, whether the random-sign encoding removes the dependence of the post-failure estimate on the sign composition without changing decoding (Section~\ref{sec:exp-decode}). Second, what shape the lower tail of the post-failure statistic has (Section~\ref{sec:exp-dist}), and how this shape determines the conditional miss rate and the extra capacity of the two rules at finite size (Section~\ref{sec:exp-cover}). Third, what the reliability and communication cost are when the upper bound is placed into the two-round protocol of a real IBLT (Section~\ref{sec:exp-e2e}), and what the computational cost is (Section~\ref{sec:exp-cost}). Finally, Section~\ref{sec:exp-deploy} summarizes the guarantees and costs of the various rule choices and gives deployment recommendations.

The experiments are divided into two kinds, statistical experiments (\exptag{STAT}) and protocol experiments (\exptag{PROTO}). The statistical experiments study only the distribution of the post-failure statistic, with configurations in Section~\ref{sec:exp-setup}; the protocol experiments run the two-round protocol on a complete IBLT, with configurations in Section~\ref{sec:exp-e2e}. A few figures additionally take a single configuration to display the shape of the distribution, with its parameters given in the figure caption.

\subsection{Configurations and sampling of the statistical experiments}\label{sec:exp-setup}
The statistical experiments do not need a complete IBLT: the sampler independently and uniformly generates $d$ random hyperedges and runs peeling, and on the same mapping simultaneously computes $X=(\widehat d-d)/s(d)$ of the standard encoding under one-sided and balanced differences, together with the exact integers $\sum_vC_v$ and $\sum_vC_v^2$ of the random-sign encoding, and then computes the upper bound of the two rules from the same observation.

Each configuration is determined by $(k,M,d)$. Two axes describe the position of a configuration: $c=d/\sqrt M$ (\eqref{eq:c}) and the first-round load $d/(\alpha_cM)$, controlling different effects. The parameter $c$ determines the statistical scale of the estimation error: when $c$ is small the post-failure $\widehat d-d$ is dominated by the $\pm2$ offset of a few duplicate edge pairs, and when $c$ is large it tends to a continuous distribution; the load determines the position of the peeling phase transition, subcritical below 0.95, in the critical window from 0.95 to 1.05, and overloaded above 1.05. For fixed $k,M$ the two axes satisfy $d/(\alpha_cM)=c/(\alpha_c\sqrt M)$, so only changing $M$ can move the load while keeping $c$ fixed, thereby separating the two effects.

The statistical experiments take $k=3,4,5$ and $M\in\{256,1024,4096,16384\}$. For each pair $(k,M)$, $d$ takes the union of three kinds of values: the 8 values making $c\in\{0.5,1,2,4,8,16,32,64\}$, the 5 values making the load $\in\{0.95,1,1.05,1.5,3\}$, and the small differences $d=2,3,6$. After removing the points with load exceeding 3 there are 189 configurations, each collecting about 20,000 failures. $M=64$ gives 42 configurations by the same rule, used separately to test the boundary of applicability of the rules. The positions of these configurations on the two axes are shown in Figure~\ref{fig:exp-grid} of the appendix.

\paragraph{Planted sampling.}
Failure samples have two sources. When the load is near or above the critical value, failure is common, and it suffices to randomly generate a mapping and keep the failure samples; this direct sampling gives the true failure-conditional distribution. When the load is low the failure probability is small, and direct sampling is too expensive. By Proposition~\ref{prop:plant}, when $d=o(M)$ failure is almost always caused by a pair of keys with completely identical mappings, so such samples can be constructed directly: randomly choose a pair of differing keys to fall on the same group of $k$ cells, and draw the mappings and signs of the remaining keys independently. In the statistical experiments the 84 configurations with high failure rate use direct sampling, and the remaining 105 use this planted sampling, marked separately in the figures (Appendix~\ref{app:sampling-impl}).

The planted-sampling control (\exptag{PLANT}) compares the two kinds of samples on 12 planted configurations where direct sampling is still affordable: among direct failures the proportion caused by a duplicate edge pair is 0.988 to 1.000, the KS distance of $X$ is 0.0048 to 0.0079, comparable to the sampling error at this sample size, and the miss/u of the Skellam rule also agrees. The planted approximation deviates only near the critical load when $M$ is small, and the planted configurations of the statistical experiments are not in this region (Appendix~\ref{app:planted}).

\subsection{The random-sign encoding removes the dependence on the sign composition}\label{sec:exp-decode}
This subsection verifies two properties of the random-sign encoding: the post-failure estimate no longer depends on the sign composition, and the result of decoding is identical to that of the standard encoding.

\paragraph{Conditional bias.}
Under the random-sign encoding, the post-failure $\widehat d$ is unbiased for every sign composition; under the standard encoding, the selection of the mapping by failure introduces into $\widehat d$ a bias that depends on the sign composition. Figure~\ref{fig:encoding} compares three conditions on the same batch of 4 million mappings, with the success or failure of peeling identical in each case: the failure samples of the standard encoding under a one-sided difference are shifted right as a whole by about $0.6s(d)$ and are almost unbiased under a balanced difference; the failure samples of the random-sign encoding almost coincide with all samples (numerical values in Appendix~\ref{app:dist}).

\begin{figure}[!htbp]
\centering\includegraphics[width=\linewidth]{../figure/fig-encoding.pdf}
\caption{$X=(\widehat d-d)/s(d)$ under three conditions on the same batch of mappings ($k=3$, $M=16384$, $d=13314$, failure rate 0.125).}
\label{fig:encoding}
\end{figure}

The bias of the standard encoding depends on both the sign composition and $k$, and the receiver cannot correct it in advance. Figure~\ref{fig:composition}(a) compares the two encodings at five sign compositions: the solid lines of the standard encoding are near zero at the balanced point and deviate more as the difference tilts to one side, with the direction of the deviation determined by $k$; for a one-sided difference it is $+0.638$ at $k=3$ and $-0.387$ at $k=6$ (mechanism in Appendix~\ref{app:std}). The receiver knows $k$ and therefore the direction of the bias, but it does not know $(d_+,d_-)$ in advance and cannot determine the size of the correction. The dashed lines of the random-sign encoding stay near zero, and after magnifying the vertical axis by about 40 times (Figure~\ref{fig:composition}(b)) there is still no discernible difference among the five sign compositions. This agrees with Proposition~\ref{prop:composition}: the coding coefficients $s_i\varepsilon_i$ are iid random signs, and their distribution is independent of $(d_+,d_-)$. The 84 direct-sampling configurations of the statistical experiments give the same conclusion (Appendix~\ref{app:dist}).

\begin{figure}[!htbp]
\centering\includegraphics[width=\linewidth]{../figure/fig-composition.pdf}
\caption{The mean of the post-failure $X$ as the sign composition varies ($M=16384$, failure probability about 0.1, with the two encodings sharing the same batch of mappings on the same configuration). Color indicates $k$. (a) Solid lines with filled markers are the standard encoding, dashed lines with open markers are the random-sign encoding; (b) the random-sign encoding data in (a) magnified about 40 times on the vertical axis, with error bars of two standard errors.}
\label{fig:composition}
\end{figure}

\paragraph{Decoding equivalence.}
The random-sign encoding does not change the set of keys recovered by decoding (Proposition~\ref{prop:decode}), so it inherits the threshold and the failure probability of the standard IBLT. The decoding-equivalence test (\exptag{DEC}) verifies this with paired trials on a complete IBLT: the two encodings share the same batch of keys and cell hash, and each trial checks whether the success or failure of the first round, each accepted key, and the set recovered at termination agree with an independent degree-one peeling reference. In more than 500 million paired trials the number of disagreements is zero (Appendix~\ref{app:decode-exp}). Table~\ref{tab:decode} gives the failure probability near the threshold: after normalizing the load by the respective $\alpha_c$, the five values of $k$ almost coincide, and the transition band narrows as $M$ grows, consistent with the critical window of width of order $M^{-1/2}$ in Section~\ref{sec:model}.

\begin{table}[!htbp]
\centering
\caption{Failure probability of first-round decoding. Columns are the load $d/(\alpha_cM)$; the two encodings agree in every trial, so the values apply to both.}
\label{tab:decode}
\begin{tabular}{rr@{\hspace{1.2em}}ccc@{\hspace{1.2em}}ccc@{\hspace{1.2em}}ccc}
\toprule
& & \multicolumn{3}{c}{$M=1024$} & \multicolumn{3}{c}{$M=16384$} & \multicolumn{3}{c}{$M=65536$}\\
\cmidrule(lr){3-5}\cmidrule(lr){6-8}\cmidrule(lr){9-11}
$k$ & $\alpha_c$ & 0.98 & 1 & 1.02 & 0.98 & 1 & 1.02 & 0.98 & 1 & 1.02\\
\midrule
3 & 0.818 & 0.298 & 0.715 & 0.958 & 0.000 & 0.648 & 1.000 & 0.000 & 0.619 & 1.000\\
4 & 0.772 & 0.266 & 0.743 & 0.977 & 0.000 & 0.654 & 1.000 & 0.000 & 0.620 & 1.000\\
5 & 0.702 & 0.268 & 0.756 & 0.975 & 0.000 & 0.661 & 1.000 & 0.000 & 0.629 & 1.000\\
6 & 0.637 & 0.284 & 0.740 & 0.971 & 0.000 & 0.665 & 1.000 & 0.000 & 0.627 & 1.000\\
10 & 0.461 & 0.383 & 0.759 & 0.966 & 0.000 & 0.663 & 1.000 & 0.000 & 0.630 & 1.000\\
\bottomrule
\end{tabular}
\end{table}

The additional requirement of the random-sign encoding on the decoder is the same as when the standard IBLT handles a balanced difference. The random signs make a cell of degree greater than one possibly show a count of $\pm1$ (Section~\ref{sec:design}); all such candidates are rejected at the checksum step, with not a single false acceptance. The frequency with which the random-sign encoding encounters such candidates is independent of the sign composition and equals the frequency of the standard encoding under a balanced difference (Figure~\ref{fig:decode} of the appendix). The cell fields and byte widths of the two encodings are the same; the random-sign encoding additionally sends the set size $|A|$ (8 bytes) with the table, used to recover the lower bound $|b|$.

\subsection{The discrete limit distribution accurately describes the post-failure lower tail at finite size}\label{sec:exp-dist}
The limit distribution of Section~\ref{sec:inversion} almost coincides with the empirical post-failure lower tail without any fitting. Figure~\ref{fig:crosscdf} takes six configurations with $c$ from 0.25 to 8 at $M=16384$ and compares the empirical CDF of the post-failure $(\widetilde d-d)/s(d)$, the limit distribution of Theorem~\ref{thm:cross}, and $N(0,1)$: on all six configurations the empirical CDF almost coincides with the limit distribution, while $N(0,1)$ deviates clearly when $c$ is small. The empirical distribution is computed exactly from integer sums of squares, and the limit distribution uses the $G_d$ of the main text, with no fitted parameters. The finite-size error decreases as $M$ grows; at $c=1,2,4$ the KS distance drops from within 0.034 at $M=256$ to within 0.012 at $M=16384$ (Table~\ref{tab:limit-ks} of the appendix).

\begin{figure}[!htbp]
\centering\includegraphics[width=\linewidth]{../figure/fig-crossover-cdf.pdf}
\caption{The post-failure lower tail of $(\widetilde d-d)/s(d)$ at different $c$ ($k=3$, $M=16384$; 400,000 planted samples per configuration for $c=1,2,4$, and 100,000 for the rest). The limit distribution is computed directly from Theorem~\ref{thm:cross}, with zero fitting. Below the title are the KS distances from the empirical CDF to the limit distribution and to $N(0,1)$, with gray lines at 1\% and 5\%.}
\label{fig:crosscdf}
\end{figure}

The steps of the lower tail come from the duplicate edge pairs and are filled in as $c$ grows. When $c$ is small, failure is almost always caused by one duplicate edge pair; when the two edges have opposite signs the observation is 2 below $d$. This $-2$ offset is large relative to $s(d)$; at $c=0.25$ it lies at a standardized error of about $-5.7$, and about four tenths of the failure samples fall on this step, the large step at $c=0.25$ in Figure~\ref{fig:crosscdf}. As $c$ grows, the overlap among the remaining edges brings many $\pm2/k$ offsets that fill in the steps, and the lower tail gradually approaches the Gaussian (Appendix~\ref{app:limit}).

This discrete structure determines the behavior of the two rules at finite size. The Gaussian rule treats the discrete left tail as a continuous tail: when $c$ is small $s(d)$ is small, the $-2$ offset caused by opposite signs of the duplicate edge pair lies far in the Gaussian tail, and the probability of rejecting the true $d$ exceeds $u$. The $G_n$ of the Skellam rule counts the probability of the observation value itself and accepts the candidate when the observation falls on a lattice point, so it is conservative in the small-$c$ and low-load regions where the values are sparse. The conditional miss rate results of the next subsection agree with these two points.

\subsection{Conditional miss rate and extra capacity at finite size}\label{sec:exp-cover}
This subsection tests the risk of the two rules at finite table size and measures the extra capacity paid by the upper bound. Theorem~\ref{thm:coverage} guarantees that the conditional miss rate of the Skellam rule is asymptotically at most $u$ for all $d\ge2$, and the Gaussian rule has the same guarantee only when $c\to\infty$ (Corollary~\ref{cor:gauss}). The conditional miss rate is the proportion of misses ($U<d$) among failure samples, reported as miss/$u$; a value not exceeding 1 means the risk is within the target. The extra capacity is $(U-d)/d$.

\paragraph{Conditional miss rate.}
The Skellam rule keeps the risk near the target over the whole region. The circles in Figure~\ref{fig:coverage} all fall near miss/u $=1$ on the 189 configurations, and the maximum 1.04 is within the sampling error of 20,000 failures per configuration. Further, the 12 configurations closest to the risk upper bound are retested with independent samples (\exptag{WORST}), covering bounded $c$, the critical window, and the overload region: in Figure~\ref{fig:confirm} the estimated miss/u values do not exceed 0.99, and the largest confidence upper bound is $1.024u$, meeting the $1.1u$ criterion set before sampling (Appendix~\ref{app:cover-more}). The low-load configurations with $c<8$ are conservative, most clearly at $u=0.01$, for the reason of the discrete structure of Section~\ref{sec:exp-dist}: the lower tail concentrates on a few lattice points, and the rule cannot use up the risk budget exactly between lattice points.

\begin{figure}[!htbp]
\centering\includegraphics[width=\linewidth]{../figure/fig-coverage.pdf}
\caption{miss/u versus $c=d/\sqrt M$ in the statistical experiments ($k=3,4,5$, $M=256$--16384, about 20,000 failures per configuration; configurations with zero miss are drawn on the lower boundary). The top row is direct sampling, the bottom row is planted sampling; circles are the Skellam rule, triangles are the Gaussian rule.}
\label{fig:coverage}
\end{figure}

\begin{figure}[!htbp]
\centering\includegraphics[width=\linewidth]{../figure/fig-confirm.pdf}
\caption{Estimated miss/u of the Skellam rule and two-sided 95\% intervals (Clopper--Pearson) on the \exptag{WORST} configurations. Configurations are ordered by load, with the load in parentheses; the dashed line is miss/u $=1$.}
\label{fig:confirm}
\end{figure}

The Gaussian rule meets the target only when $c$ is large. The triangles in Figure~\ref{fig:coverage} are near 1 when $c\ge8$ and rise quickly as $c$ decreases, already exceeding the target by more than 8 times at $u=0.01$, $c=1$ (Table~\ref{tab:gauss-c} of the appendix). This excess does not disappear as the table size grows, consistent with the finite limit of Theorem~\ref{thm:cross}. The reason is again the discrete structure: when $c$ is very small, the post-failure observation falls near $d+2$ and $d-2$ with probability one half each (Proposition~\ref{prop:smalld}), the Gaussian tail gives the latter an extremely small probability, the rule therefore accepts the candidate $d-2$, and the miss rate tends to 1/2.

\paragraph{Extra capacity.}
The extra capacity is determined by the table size and decreases as $M^{-1/2}$, with the two rules almost identical. At $u=0.01$ and $c\ge8$, the median extra capacity of the Skellam rule drops from 25\% at $M=256$ to 2.6\% at $M=16384$, because $U-d$ is about 2.3--2.8 times $s(d)$ and $s(d)/d\approx\sqrt{2/M}$. In the two-round protocol this extra capacity appears as the increase of the second-round table size over the oracle; in Figure~\ref{fig:protocol}(a) it falls along the dotted line of slope $-1/2$ (Section~\ref{sec:exp-e2e}). When $c$ is bounded the difference itself is small and few extra elements appear; at $c=4$ the median of $U-d$ is 15 (Appendix~\ref{app:cover-more}).

For engineering use we recommend $M\ge1024$. The null distribution of Section~\ref{sec:inversion} omits the overlap between the duplicate edge pair and the remaining edges, whose occurrence rate $(d-2)k^2/M$ grows as the table size decreases. At $M=64$ this omission is no longer negligible, and the largest miss/u of the Skellam rule is 1.08, exceeding the sampling error. At $M=256$ no miss/u exceeding the sampling error is observed, but the occurrence rate of the omitted structure is 4 times that at $M=1024$, and the median extra capacity also reaches 25\%.

\subsection{The two-round protocol on a real IBLT}\label{sec:exp-e2e}
Theorem~\ref{thm:protocol} decomposes the two-round failure probability into three terms: the joint miss rate, the peeling failure when the capacity is sufficient, and the checksum misjudgment. This subsection runs the two-round protocol on a complete IBLT, tests the relative sizes of the first two terms in the real protocol, measures the communication cost of the random-sign IBLT, and compares it with the existing calibration rule.

The protocol experiments (\exptag{PROTO}) take $k=3,4,5$, first-round table size $M_1\in\{256,1024,4096,16384\}$, and first-round loads 0.95, 1, 1.05, 1.5, 3, with the difference taking two sign compositions, one-sided (all on the $A$ side) and balanced, for 120 configurations in total; the two sides additionally hold 1000 common elements. Each configuration runs until the first round accumulates 20,000 failures. After each first-round failure, three upper bounds determine the second-round table size separately and the second round runs: the Skellam rule, the Gaussian rule, and the oracle that takes the true $d$ directly. The two rules take $U_b=\max(U,|b|)$ as in Section~\ref{sec:flow}, the second-round table size is $\max\{\lceil U_b/(0.8\alpha_c)\rceil,64\}$, and a fresh cell hash is used. The three upper bounds share all keys and hashes within the same trial, so the difference comes only from the upper bound itself and the three can be compared pairwise trial by trial (implementation details in Appendix~\ref{app:protocol}).

The reliability of the two-round protocol is comparable to that of the oracle. In Figure~\ref{fig:protocol}(b), the final failure rates of the two rules fall in the same range as the oracle, none exceeding 0.13\% when $M_1\ge1024$, and at $u=0.01$ the maximum is not higher than the oracle (Table~\ref{tab:e2e} of the appendix); the two rules are statistically indistinguishable. The final failure comes mainly from the peeling failure of the second round itself, consistent with the expectation that the second term dominates in Theorem~\ref{thm:protocol}: at $u=0.01$, no more than 6\% of the final failures follow a miss. The reason is that the second-round table size leaves a margin $1/(0.8\alpha_c)$, which absorbs most slight underestimates. The same reason also explains why the rules are no worse than the oracle: the upper bound given by the rules is usually larger than $d$, the second-round table is slightly larger than the oracle, and the peeling failure is therefore smaller.

The extra communication cost is determined only by $M_1$ and decreases as $M_1^{-1/2}$. In Figure~\ref{fig:protocol}(a), the configurations with different $k$, load, and sign composition at the same $M_1$ almost coincide and fall along the dotted line of slope $-1/2$. At $u=0.01$, the increase of the second-round table size of the Skellam rule over the oracle drops from 25.6\% at $M_1=256$ to 2.6\% at $M_1=16384$, consistent with the extra capacity of Section~\ref{sec:exp-cover}; converted to total bytes, the overhead drops from 14.7\% to 1.5\%.

\begin{figure}[!htbp]
\centering\includegraphics[width=\linewidth]{../figure/fig-protocol.pdf}
\caption{The two-round protocol as the first-round table size $M_1$ varies (120 configurations, 13.94 million trials per $u$ in total). (a) The increase of the second-round table size over the oracle after a first-round failure, one point per configuration, with the dotted line of slope $-1/2$; (b) the final two-round failure rate of each configuration at $u=0.01$, with zero configurations drawn at 0 below the axis break.}
\label{fig:protocol}
\end{figure}

Figure~\ref{fig:scale} gives the same trend at fixed load (first-round target load 1.5, balanced, $u=0.01$): as the table size grows, the miss/u of the Skellam rule gradually approaches 1 from below, while the Gaussian rule stays near 1 (in the top row each point merges all first-round failures of the same configuration, Appendix~\ref{app:cover-more}); the total byte overhead relative to the oracle drops from 16.8\% at $M_1=256$ to 1.7\% at $M_1=16384$, tending to become negligible as the table size grows.

\begin{figure}[!htbp]
\centering\includegraphics[width=\linewidth]{../figure/fig-scale.pdf}
\caption{The two-round protocol at fixed load as the table size varies (balanced, $u=0.01$, target load 1.5). One column per $k$; the top row merges all first-round failures of the same configuration (about 330,000 per point); the bottom row is the communication overhead relative to the oracle of the same configuration.}
\label{fig:scale}
\end{figure}

\paragraph{Comparison with the existing calibration rule.}
The existing calibration rule depends on the sign composition under the standard encoding, which is the consequence in the protocol of the bias shown in Section~\ref{sec:exp-decode}. We plug the calibration factor $q_{0.01}$ of \cite{wu2026selfsizing} into the same protocol, taking $U=\lceil\widehat d/q_{0.01}\rceil$, with the rest of the settings the same as the two rules of this paper (Table~\ref{tab:vldb} of the appendix). Under the standard encoding, at load 0.95--1.0 the same factor gives a conditional miss rate of $0.05$--$0.40u$ for a one-sided difference and $0.86$--$0.94u$ for a balanced one, differing by an order of magnitude; the random-sign encoding brings both to $0.80$--$1.01u$. The calibration factor must be calibrated by simulation separately for each set of parameters, and the literature gives values only for $k=3$, $M_1\le4096$, and $u=0.01$; on these parameters it is close to the Skellam rule in final failure rate and total bytes. The Skellam rule is computed directly from the theoretical distribution of Section~\ref{sec:inversion}, requiring no calibration for any $k$, $M$, load, or $u$.

\subsection{Computational cost}\label{sec:exp-cost}
Compared with the standard encoding, the extra computational time comes almost entirely from the sizing of the Skellam rule, and the overhead of the random-sign encoding itself is negligible. With the Gaussian rule as the baseline, the Skellam rule increases the total protocol time by less than half when the load does not exceed 1.5; at load 3, sizing becomes the main cost and the total time is 2.8--3.8 times that of the Gaussian rule; combining the two rules by the first-round observation avoids this part of the cost (Section~\ref{sec:switch}).

The random-sign encoding makes any difference look like a balanced difference to the decoder; under a one-sided difference the candidates checked per recovered key increase from 1.00 to 1.77 ($k=3$), and the resulting increase in decoding time is less than 1\% of the total protocol time.

The sizing cost of the Skellam rule grows with the load. Its exact CDF requires computing and looking up Poisson probability tables, which depend only on $(M,k)$ and the candidate $n$, and the same receiver can cache and reuse them across synchronizations. Table~\ref{tab:proto} gives the time of the complete protocol at $M_1=16384$, including first-round encoding and decoding, sizing, and the second round: at load 1.0--1.05 the Skellam rule is 14--23\% more than the Gaussian rule, and at load 1.5 it is 39--43\% more; the sizing cost grows roughly as $c^2$, and at load 3 the total time of the Skellam rule is 57--88 ms. The conditions of the computational-cost measurement (\exptag{COST}), the per-sizing time under different cache conditions, and the memory use are in Appendix~\ref{app:timing}.

\begin{table}[!htbp]
\centering
\caption{Total protocol time per trial (ms, median of 10 repetitions; $M_1=16384$, one-sided difference; the Skellam rule caches the Poisson probability tables between successive trials).}
\label{tab:proto}
\begin{tabular}{lrrrrrr}
\toprule
& \multicolumn{3}{c}{Skellam rule} & \multicolumn{3}{c}{Gaussian rule}\\
\cmidrule(lr){2-4}\cmidrule(lr){5-7}
Load & $k=3$ & $k=4$ & $k=5$ & $k=3$ & $k=4$ & $k=5$\\
\midrule
1.0 & 12.7 & 16.6 & 17.6 & 10.9 & 14.0 & 15.4\\
1.05 & 15.9 & 19.2 & 21.5 & 13.1 & 15.6 & 18.4\\
1.5 & 19.7 & 24.0 & 27.4 & 14.0 & 16.8 & 19.7\\
3 & 57.1 & 84.1 & 88.1 & 18.8 & 22.4 & 25.8\\
\bottomrule
\end{tabular}
\end{table}

\subsection{Rule choice and deployment}\label{sec:exp-deploy}
At deployment we recommend combining the two rules by the first-round observation (Section~\ref{sec:switch}). Table~\ref{tab:sizing-options} summarizes the guarantees and costs of using each rule alone and in combination. The Gaussian rule needs only a microsecond-scale scalar search, and its guarantee range is $d\gg\sqrt M$, suited to applications whose workload can guarantee this condition and that pursue low latency; the table size alone is not enough to rule out bounded $c$ and a fixed small $d$. The Skellam rule holds for all $d\ge2$ and does not need to know which load region $d$ falls in, at the cost of millisecond-scale computation that looks up Poisson tables, which becomes the main cost of the protocol when the first round is heavily overloaded (Section~\ref{sec:exp-cost}). The combined use takes the Gaussian rule at high load and the Skellam rule otherwise, keeping the uniform guarantee unchanged and the time in the overload region the same as the Gaussian rule.

\begin{table}[!htbp]
\centering
\caption{Choice of the sizing rule.}
\label{tab:sizing-options}
\begin{tabular}{@{}p{.19\linewidth}p{.24\linewidth}p{.29\linewidth}p{.18\linewidth}@{}}
\toprule
Rule & Range of conditional miss rate asymptotically at most $u$ & Suitable scenario & Computational cost\\
\midrule
Gaussian rule & $d\gg\sqrt M$ & workload guarantees the range, low latency valued & microsecond-scale, scalar search\\
Skellam rule & all $d\ge2$, uniformly & $d$ may range from far below $\sqrt M$ to overload & millisecond-scale, cached Poisson tables\\
Combined use & all $d\ge2$, uniformly & needs the uniform guarantee while controlling overload time & microsecond-scale at high load, otherwise as the Skellam rule\\
\bottomrule
\end{tabular}
\end{table}

\section{Related work}\label{sec:related}
\paragraph{IBLT and set reconciliation.} IBLT\cite{gm} and the invertible Bloom filter\cite{eppstein2011whatsdifference} list set differences by peeling. Decoding success is determined by the 2-core of a random hypergraph\cite{molloy,dw,dm}. At low load, duplicate edge pairs dominate failure\cite{gm,yw}, and the probability of a linear-sized core at subcritical load is exponentially small\cite{kim2008poisson,dn08}. Subsequent work improves the hypergraph structure\cite{lazaro2021irregular}, develops rateless codes\cite{yang2024rateless,lazaro2023ratecompatible}, or predicts and breaks the 2-core after decoding failure\cite{gadurek2026breaking}. Belazzougui, Kucherov, and Walzer\cite{belazzougui2024better} use finite-field markers to distinguish the direction of the difference and make the failure probability exponentially small. These works seek to reduce failure or resume decoding after failure. We read the same table after failure to determine the next-round capacity.

\paragraph{Estimating the difference size.} A common approach to estimating $d$ is to send a separate estimator, such as the strata estimator\cite{eppstein2011whatsdifference}, min-wise hashing\cite{broder2000minwise}, the odd sketch\cite{mitzenmacher2014odd}, or GXBits\cite{jia2024gxbits}. PBS\cite{gong2021pbs} first estimates $d$ using 128 independent ToW sketches, then sets capacity to a fixed multiple of the estimate. It transmits and accounts for the estimator and reconciliation structure separately. We obtain the estimate from the same counts used for decoding, avoiding a separate estimator. This reuse requires analysis conditional on decoding failure: the estimate is needed only after failure, and failure selects the mapping.

\paragraph{Random-sign sketches and second-moment estimation.} AMS\cite{alon1999spacecomplexity} estimates $F_2$ using a sum of random signs. CountSketch\cite{charikar2004countsketch} supports point queries with random-sign counters. Thorup and Zhang\cite{thorup2004tabulation} analyze its second-moment extension and give a centered second-moment estimate for a single counter array. Fast-AGMS\cite{cormode2005fastagms} uses the same structure for distributed queries. $\widehat d$ has the same algebraic form as these estimators. Its conditional unbiasedness given the mapping also follows the structure of the Hutchinson trace estimator\cite{hutchinson1990trace}. These sketches carry no decodable keySum and do not address the distribution conditional on decoding failure. Fischer and Nakos\cite{fischer2025sparse} combine hypergraph peeling and a random-sign CountSketch in a single measurement. An independent short sketch supplies the tail norm, which their analysis treats as known. We encode the shared random signs in the IBLT's count, keySum, and checksum. The same table supports peeling and, after failure, provides an estimate independent of the difference direction, together with a post-failure upper bound valid at all loads.

\paragraph{Reading the same table after failure.} Prior work\cite{wu2026selfsizing} reads count energy after first-round failure to set the second-round capacity. Its analysis models sign composition with iid signs and obtains a conditional mean equal to $d$. Its upper-bound rule uses a chi-square approximation under deep overload and per-configuration simulation calibration at other loads. We encode random signs in all three IBLT fields, making these properties hold exactly for every fixed mapping. We also derive a post-failure distribution and upper bound valid at all loads.

\paragraph{Conditional limits and discrete approximations.} For bounded $c$ (\eqref{eq:c}), we derive the limit distribution using the Chen--Stein method\cite{arratia1989chen} for the Poisson approximation of sparse overlap. The Skellam distribution\cite{skellam1946} describes the difference between two types of counts. At low load, we replace the failure-conditioned distribution with the planted distribution, which is the distribution size-biased by the duplicate edge pair count\cite{arratia2019sizebias}. When $d\gg\sqrt M$, we analyze the random-sign quadratic form given the mapping using the cumulant method on a dependency graph\cite{janson1988cumulant}. This requires only concentration bounds\cite{warnke2016typical,combes2015mcdiarmid} and degree control for the overlap graph. We prove a Gaussian limit uniform over deterministic coefficients using an edge-exposure martingale CLT\cite{hall1980martingale}. A general quadratic form CLT is given in\cite{dejong}. The upper-bound rule uses classical test inversion\cite{neyman1937}.

\section{Discussion and conclusion}\label{sec:conclusion}
The random-sign IBLT lets one table simultaneously perform decoding and difference-size estimation, with the cells and the decoding capability both unchanged. Whether decoding succeeds is determined only by the mapping, and the randomness of the estimate comes from the signs; this separation makes the estimate exactly unbiased for every fixed mapping and identically distributed for all sign compositions, and makes the post-failure distribution fully characterizable: Gaussian when $d\gg\sqrt M$, and a discrete distribution of a $\pm2$ term plus a Skellam term when $c=d/\sqrt M$ is bounded. The same family of distributions describes all scales, and the Skellam rule that takes it as the null distribution makes the conditional miss rate asymptotically at most the target risk uniformly over all $d\ge2$.

At deployment we recommend combining the two rules by the first-round observation, which keeps the uniform guarantee for all $d\ge2$ with the time in the overload region the same as the Gaussian rule; the guarantees and costs of the various choices are in Section~\ref{sec:exp-deploy}. On the deployment parameters, the Skellam rule meets the target on all experimental configurations with $M\ge256$, and for engineering use we recommend $M\ge1024$ (Section~\ref{sec:exp-cover}). The checksum can be shortened to 32 bits in exchange for narrower cells, leaving the risk of sizing unaffected, and the probability of a verification error is still at most $10^{-4}$ at $M_1=16384$ and load at most 3 (Section~\ref{sec:protocol}). Storing the counts of several independent signs per cell reduces the absolute error but does not change the post-failure conditional scale, and one still performs inversion by the post-failure distribution (Appendix~\ref{app:multi}). The random-sign encoding does not increase the space overhead of the cells, and the increase in decoding time is less than 1\% of the total protocol time (Section~\ref{sec:exp-cost}).

Three questions remain for future work. First, before building the tables, reject more candidates with a rigorous finite-parameter tail upper bound, reducing the computational cost of the Skellam rule. Second, turn the errors of Theorem~\ref{thm:cross} and Theorem~\ref{thm:clt} into computable constants, giving a finite risk bound for a specified $M$. Third, extend to multiset differences whose weights can exceed one: there the random-sign count gives an unbiased estimate of the $F_2$ of the difference vector, and the decoding and post-failure analysis of this extension remain to be completed.


\bibliographystyle{plain}
\bibliography{references}

\appendix
\section{Notation}\label{app:notations}
This appendix collects the notation used in the main text and the principal proofs, ordered by section for easy reference. Except where quantities are compared across different scales, the dependence of each quantity on $M$ is suppressed; we fix $k\ge3$ and let $M\to\infty$, and a subscript $k$ denotes a constant depending only on $k$. The symbol $\Rightarrow$ denotes convergence in distribution. Section~\ref{sec:design} uniformly writes the coding coefficient $s_i\varepsilon_i$ as $\varepsilon_i$, and this convention is followed below.

\begin{table}[!htbp]
\centering\footnotesize
\caption{Principal notation.}
\begin{tabular}{@{}l p{0.72\linewidth}@{}}
\toprule
Notation & Meaning\\
\midrule
\multicolumn{2}{@{}l}{\bfseries Model and peeling}\\
$A$, $B$ & sets held by the two sides; $|A|$ and $|B|$ are the set sizes\\
$M$, $M_1$ & number of cells; number of cells in the first-round table\\
$k$ & number of cells selected by each key\\
$x$, $x_i$ & key; the $i$-th differing key\\
$\mathrm{chk}(x)$ & checksum of key $x$\\
$d$ & difference size $|A\triangle B|$\\
$d_+$, $d_-$ & number of keys in $A$ only and in $B$ only\\
$b$ & difference of the two set sizes, $d_+-d_-=|A|-|B|$\\
$S_i$ & the $k$-subset selected by key $i$\\
$\mathcal G$ & unsigned mapping with key labels\\
$H_{ij}$ & overlap $|S_i\cap S_j|$ of two edges\\
$F$ & peeling failure on a single table; $\Pr(F)$ is the failure probability\\
$F_1$, $F_2$ & peeling failure in the first and second rounds\\
$E$ & checksum misjudgment event\\
$\alpha_c$ & peeling threshold (about $0.818$ for $k=3$ and $0.772$ for $k=4$)\\
\midrule
\multicolumn{2}{@{}l}{\bfseries Encoding and measurement}\\
$s_i$ & set direction of key $i$, taking values $\pm1$\\
$\mathcal U$, $\varepsilon(\cdot)$ & universe of keys, and the sign hash shared by the two sides\\
$\varepsilon_i$ & random sign $\varepsilon(x_i)$ of key $i$; in the analysis the coding coefficient $s_i\varepsilon_i$ is written as $\varepsilon_i$\\
$C_v$ & count of cell $v$, $\sum_{i:v\in S_i}\varepsilon_i$ (in the standard encoding, $\sum_{i:v\in S_i}s_i$)\\
$\widehat d$ & estimator of $d$, see~\eqref{eq:dhat}; used by the Gaussian rule\\
$s(n)$ & standardization scale $\sqrt{2n(n-1)/(M-1)}$, see~\eqref{eq:s}\\
$X$ & standardized estimation error $(\widehat d-d)/s(d)$, see~\eqref{eq:X}\\
$\widetilde d$ & estimator of $d$, $\frac1k\sum_vC_v^2$, see~\eqref{eq:raw-energy}; used by the Skellam rule\\
$c_F$ & normalized amount by which the overlap exceeds its unconditional mean in the failure sample, see~\eqref{eq:projection}\\
$\Phi$, $\mathcal N(0,1)$ & CDF and distribution of the standard Gaussian\\
\midrule
\multicolumn{2}{@{}l}{\bfseries Inversion and protocol}\\
$u$ & target risk, taken as $0.01$ and $0.05$ in the experiments\\
$n$ & candidate difference size, an integer $n\ge2$\\
$U$, $U_b$ & upper bound on $d$ given by the selected rule; $\max(U,|b|)$ after enforcing the lower bound $|b|$\\
$t(n)$ & acceptance lower bound of the Gaussian rule, see~\eqref{eq:gauss}\\
$\chi^2_{M-1,u}$ & lower $u$ quantile of the chi-square distribution with $M-1$ degrees of freedom\\
$\lambda_n$ & parameter of the Skellam term in the null distribution, see~\eqref{eq:lambda}\\
$Z_n$ & $\Skellam(\lambda_n,\lambda_n)$, the difference of two independent $\Pois(\lambda_n)$\\
$J_\lambda$ & $2\varepsilon_1\varepsilon_2+\frac2k\Skellam(\lambda,\lambda)$, see~\eqref{eq:J}\\
$G_n$ & CDF of the null distribution $n+J_{\lambda_n}$ of candidate $n$ (Definition~\ref{def:null})\\
$c$ & $d/\sqrt M$ (\eqref{eq:c}); its limit in Theorem~\ref{thm:cross}\\
$m(\cdot)$ & second-round table-size rule $m(U_b)=\lceil U_b/(0.8\alpha_c)\rceil$\\
$\Pr_m(F)$ & failure probability on a new table with $m$ cells\\
\midrule
\multicolumn{2}{@{}l}{\bfseries Proofs}\\
$\Omega_2$ & sum of squared overlaps $\sum_{i<j}H_{ij}^2$\\
$\mathcal G_M$ & set of good graphs in Appendix~\ref{app:graph}, whose complement has probability an arbitrary fixed negative power of $M$\\
\bottomrule
\end{tabular}
\end{table}

\section{Post-failure bias of the standard IBLT}\label{app:std}
This appendix proves Proposition~\ref{prop:projection} and gives a numerical verification.

In the standard IBLT the count is $C_v=\sum_{i:v\in S_i}s_i$, and the coefficient is the set direction. Expanding $\widehat d$,
\begin{equation}
\widehat d-d=\frac{2}{k(1-k/M)}\sum_{i<j}s_is_j\Bigl(H_{ij}-\frac{k^2}{M}\Bigr).\label{eq:std-expand}
\end{equation}
Failure changes the distribution of the overlap: the overlap of same-sign pairs raises $\widehat d$, and the overlap of opposite-sign pairs lowers $\widehat d$, which yields Proposition~\ref{prop:projection}.

\begin{proof}[Proof of Proposition~\ref{prop:projection}]
By~\eqref{eq:std-expand},
\begin{equation*}
\E[\widehat d\mid F]-d=\frac{2}{k(1-k/M)}\sum_{i<j}s_is_j\,\E\Bigl[H_{ij}-\frac{k^2}{M}\Bigm|F\Bigr].
\end{equation*}
$F$ depends only on the unsigned mapping and is invariant under permutation of the key labels; the random mapping $\mathcal G$ is exchangeable over the labels, so $\E[H_{ij}\mid F]$ is the same for all $i<j$ and equals $k^2/M+k(1-k/M)c_F$. The number of same-sign pairs minus the number of opposite-sign pairs is $\binom{d_+}2+\binom{d_-}2-d_+d_-=(b^2-d)/2$. Substituting gives $\E[\widehat d\mid F]-d=c_F(b^2-d)$, and dividing by $d$ yields~\eqref{eq:projection}. The proof uses only these two properties of $F$, so it holds for any event that depends only on the unsigned mapping, is invariant under label permutation, and has positive probability.
\end{proof}

The simulations used in Figure~\ref{fig:bias} of Section~\ref{sec:model} and Figure~\ref{fig:composition} of Section~\ref{sec:exp-decode} directly verify Proposition~\ref{prop:projection}. The key point is that the five sign compositions on the same configuration share the same batch of mappings: $F$ depends only on the unsigned mapping, so the failure samples of the five sign compositions are identical sample by sample and hence share the same $c_F$. The other factor on the right, $b^2/d-1$, depends only on the sign composition and equals $d-1$ under a one-sided difference. Hence the conditional bias of any sign composition equals the one-sided bias multiplied by $(b^2-d)/(d^2-d)$: once $c_F$ is determined from a single one-sided point, the bias of every other level can be predicted with no fitting. Figure~\ref{fig:bias} plots exactly this prediction. The horizontal axis is $(b/d)^2$; the straight line of each color is determined by the one-sided point of that configuration at $(b/d)^2=1$, and the remaining four points fall on the line, confirming~\eqref{eq:projection}.

The full verification covers 20 configurations with $k=3,4,5,6$ and failure probability from 0.10 to 0.90 on the $M=16384$ grid, each configuration having five levels of sign composition $d_-/d\in\{0,0.1,0.25,0.4,0.5\}$. Figure~\ref{fig:bias} shows ten of these configurations with $k=3$ and $k=6$; the standard-encoding curves in Figure~\ref{fig:composition}(a) are four of the configurations with failure probability about 0.1, including $k=4$ and $k=5$. The 20 configurations have four non-one-sided sign composition levels each, 80 predictions in total. In units of $s(d)$, the difference between prediction and measurement is at most 0.0064, all within two standard errors (the standard error accounts for the sampling error of both the one-sided mean and the measured value).

\section{Encoding properties and proofs of exact moments}\label{app:encoding}
We describe the XOR implementation of keySum and checksum and prove the finite identities in Sections~\ref{sec:design} and~\ref{sec:measurement}. We then summarize their exact verification for small parameter values and discuss the AMS observation obtained from the count sum.

\subsection{XOR implementation}\label{app:xor}
In the set model, each side inserts each key at most once. We can therefore accumulate keySum and checksum by bitwise XOR, as in a common implementation of the standard IBLT, and multiply only the count by $\varepsilon(x)$. Under XOR, $-x=x$. Thus, in each difference-table cell, keySum is the XOR of the differing keys that select the cell, and checksum is the XOR of their $\mathrm{chk}$ values. Both are independent of the coding coefficients. In a pure cell, keySum and checksum equal $x_i$ and $\mathrm{chk}(x_i)$ directly, so we check whether checksum equals $\mathrm{chk}(x)$. To delete $x_i$, we subtract $C_v$ from the count and XOR the other two fields once more with $x_i$ and $\mathrm{chk}(x_i)$, respectively. The induction argument for the additive form in the next subsection also applies here. We recover the direction as $C_v\,\varepsilon(x)$, and peeling follows the same process and gives the same result as in the additive form. In a multiset, the same key may be inserted repeatedly. XOR cancels repeated contributions while the count retains them, so multisets require the additive form. If peeling modifies the difference table in place, the receiver can accumulate the two scalars $\sum_vC_v^2$ and $\sum_vC_v$ during cell-by-cell subtraction, without retaining the original difference table.


\subsection{\texorpdfstring{Proof of Proposition~\ref{prop:decode} (ideal decoding equivalence)}{Proof of ideal decoding equivalence}}\label{app:proof-decode}
\begin{proof}
We induct on the number of peeling steps. After each step, the three fields of every cell equal the encoded sum of the remaining keys: the count is $\sum s_i\varepsilon_i$, keySum is $\sum s_i\varepsilon_i x_i$, and checksum is $\sum s_i\varepsilon_i\,\mathrm{chk}(x_i)$. Each sum ranges over the remaining keys that select the cell. This invariant holds initially by~\eqref{eq:fields}.

Suppose that at some step cell $v$ has degree one in the unsigned hypergraph, with unique key $x_i$. Its count is $s_i\varepsilon_i\in\{\pm1\}$, count times keySum equals $x_i$, and checksum equals count times $\mathrm{chk}(x_i)$. The cell hash of $x_i$ also contains $v$, so $v$ passes the pure-cell test. Conversely, a cell of any other degree that passes the test is a false positive, which the ideal-purity assumption excludes. Thus the cells removable under random-sign encoding are exactly the unsigned degree-one cells. To delete $x_i$, we subtract its three-field contribution multiplied by the coding coefficient $s_i\varepsilon_i$ from its $k$ cells. This preserves the invariant and removes the corresponding edge from the unsigned hypergraph. With the same cell-selection order, the two processes reach the same endpoint. Peeling always ends at the 2-core, independently of the selection order, so $F$ depends only on $\mathcal G$.

At termination, we check all fields. Under random signs, a nonempty core can have all-zero counts, so counts alone cannot certify success. The keySum fields can cancel at the same time. Success is misjudged only if all checksums are also zero. This event is included in $E$ in Section~\ref{sec:protocol}.
\end{proof}

\subsection{\texorpdfstring{Proof of Proposition~\ref{prop:moments} (conditional unbiasedness and conditional moments)}{Proof of conditional unbiasedness and conditional moments}}\label{app:proof-moments}
\begin{proof}
Using $H_{ii}=k$ and $\varepsilon_i^2=1$, we obtain the sum of squared counts and the squared sum:
\begin{align*}
\sum_vC_v^2&=\sum_{i,j}\varepsilon_i\varepsilon_jH_{ij}=kd+2\sum_{i<j}\varepsilon_i\varepsilon_jH_{ij},\\
\frac1M\Bigl(\sum_vC_v\Bigr)^2&=\frac{k^2}M\Bigl(\sum_i\varepsilon_i\Bigr)^2=\frac{k^2}Md+2\sum_{i<j}\varepsilon_i\varepsilon_j\frac{k^2}M.
\end{align*}
Dividing the first identity by $k$ gives~\eqref{eq:raw-expand}. Subtracting the second identity from the first gives $kd(1-k/M)+2\sum_{i<j}\varepsilon_i\varepsilon_j(H_{ij}-k^2/M)$. Dividing this by $k(1-k/M)$ gives~\eqref{eq:rs-expand}. For a fixed mapping, $\E[\varepsilon_i\varepsilon_j]=0$ ($i\ne j$), so the conditional expectations of $\widehat d-d$ and $\widetilde d-d$ are both zero.

We now express the coefficients of the cross terms as inner products. Let $w_i\in\mathbb R^M$ be the incidence vector of key $i$ minus its mean $k/M$. Then $\langle w_i,w_j\rangle=H_{ij}-k^2/M$ and $\|w_i\|^2=k(1-k/M)$, so~\eqref{eq:rs-expand} becomes
\begin{equation*}
\widehat d-d=\frac{2}{k(1-k/M)}\sum_{i<j}\langle w_i,w_j\rangle\,\varepsilon_i\varepsilon_j.
\end{equation*}
Given $\mathcal G$, the inner products are constant. All randomness comes from the sign products $\varepsilon_i\varepsilon_j$. The expectation of a product of sign products is nonzero if and only if each subscript occurs an even number of times, in which case the expectation is one.

For the variance, we expand the square and consider two unordered pairs $\{i,j\}$ and $\{i',j'\}$. When the pairs coincide, the product is $\varepsilon_i^2\varepsilon_j^2=1$. When they differ, they share at most one subscript, so some subscript occurs exactly once and the expectation is zero. Thus only the diagonal terms remain:
\begin{equation*}
\Var(\widehat d\mid\mathcal G)=\frac{4}{k^2(1-k/M)^2}\sum_{i<j}\langle w_i,w_j\rangle^2,
\end{equation*}
which gives the equality in~\eqref{eq:condmoments}.

For the third central moment, we expand the cube and consider ordered triples of unordered pairs. If all three pairs coincide, the product is $\varepsilon_i\varepsilon_j$. If exactly two coincide, the product is the sign product of the third pair. Both cases have expectation zero. If the pairs are distinct, every subscript occurs an even number of times exactly when the pairs form a triangle $\{i,j\},\{i,h\},\{j,h\}$. Each triangle corresponds to $3!=6$ ordered triples. Hence
\begin{equation*}
\E[(\widehat d-d)^3\mid\mathcal G]=\frac{8\cdot6}{k^3(1-k/M)^3}\sum_{i<j<h}\Bigl(H_{ij}-\frac{k^2}{M}\Bigr)\Bigl(H_{ih}-\frac{k^2}{M}\Bigr)\Bigl(H_{jh}-\frac{k^2}{M}\Bigr).
\end{equation*}

For the two upper bounds, we use Cauchy--Schwarz: $\langle w_i,w_j\rangle^2\le\|w_i\|^2\|w_j\|^2=[k(1-k/M)]^2$. Substitution into the variance formula shows that each of the $\binom d2$ unordered pairs contributes at most $4$, giving $\Var(\widehat d\mid\mathcal G)\le2d(d-1)$. The component of the vector $\sum_i\varepsilon_iw_i$ on cell $v$ is $C_v-\frac1M\sum_uC_u$, so
\begin{equation*}
\widehat d=\frac{\|\sum_i\varepsilon_iw_i\|^2}{k(1-k/M)}\ge0,
\end{equation*}
and the triangle inequality gives $\|\sum_i\varepsilon_iw_i\|\le\sum_i\|w_i\|=d\sqrt{k(1-k/M)}$, so $\widehat d\le d^2$.

Under the random mapping $\mathcal G$, the conditional mean is always $d$. The law of total variance therefore gives $\Var\widehat d=\E[\Var(\widehat d\mid\mathcal G)]$. The intersection size $H_{ij}$ of two independent uniform $k$-subsets has a hypergeometric distribution with mean $k^2/M$ and variance $k^2(1-k/M)^2/(M-1)$. Substitution into the variance formula shows that each unordered pair contributes $4/(M-1)$, giving $\Var\widehat d=2d(d-1)/(M-1)=s(d)^2$.
\end{proof}

\subsection{Exact verification for small parameter values}\label{app:exact}
We verify the finite identities in Sections~\ref{sec:design} and~\ref{sec:measurement} by exhaustive enumeration for small parameter values. For the six triples $(M,k,d)$ in Table~\ref{tab:exact}, we enumerate all mappings and signs, covering 155,640 states. All arithmetic uses integers and rationals, so the results are exact. In every state, common elements cancel in the difference table. Signed peeling recovers the correct keys and directions, and the unrecovered keys form exactly the unsigned 2-core. For every mapping, $\E[\widehat d\mid\mathcal G]=d$, the conditional variance and third central moment match the formulas above, and $\Var(\widehat d\mid\mathcal G)\le2d(d-1)$. The distribution of $\widehat d$ over the random signs is the same for every direction vector. Table~\ref{tab:exact} summarizes the unconditional and post-failure statistics.

\begin{table}[!htbp]
\centering
\caption{Exact enumeration results for small parameter values.}
\label{tab:exact}
\begin{tabular}{@{}crccccc@{}}
\toprule
$(M,k,d)$ & states & $p_F$ & $\Var\widehat d$ & $\E[\widehat d\mid F]$ & $\Var(\widehat d\mid F)$ & $\Pr(\widehat d=0\mid F)$\\
\midrule
$(4,3,2)$ & 40 & $1/4$ & $4/3$ & 2 & 4 & $1/2$\\
$(4,3,3)$ & 160 & 1 & 4 & 3 & 4 & 0\\
$(4,3,4)$ & 560 & 1 & 8 & 4 & 8 & $27/512$\\
$(5,3,3)$ & 1,760 & $2/5$ & 3 & 3 & 4 & 0\\
$(5,3,4)$ & 11,440 & 1 & 6 & 4 & 6 & $93/8000$\\
$(6,3,4)$ & 141,680 & $2029/4000$ & $24/5$ & 4 & $11256/2029$ & $75/4058$\\
\bottomrule
\end{tabular}
\end{table}

In every row of Table~\ref{tab:exact}, $\Var\widehat d$ equals $s(d)^2=2d(d-1)/(M-1)$ and $\E[\widehat d\mid F]$ equals $d$, consistent with conditional unbiasedness given the mapping in~\eqref{eq:condmoments}. Failure preserves the mean and increases the scale. For the three parameter sets with $p_F<1$, $\Var(\widehat d\mid F)$ is $3$ times, $4/3$ times, and about $1.16$ times the unconditional variance, respectively. When $p_F=1$, failure is certain and the two variances coincide. After failure, $\widehat d$ equals zero with positive probability. For $(4,3,2)$, this probability is $1/2$.

The values in the $(4,3,2)$ row come from a duplicate edge pair. Each key independently and uniformly chooses one of the $\binom43=4$ $3$-subsets. When the two keys choose the same subset, each of these three cells has degree two. There is no pure cell, and peeling stops immediately, so $p_F=1/4$. The counts of the three cells then all equal $\varepsilon_1+\varepsilon_2$, and the remaining cells have zero counts. When the two signs are equal, $\sum_vC_v^2=12$ and $\sum_vC_v=\pm6$, so~\eqref{eq:dhat} gives $\widehat d=4$. When they differ, all counts are zero and $\widehat d=0$. The two cases each have probability $1/2$. Thus $\E[\widehat d\mid F]=2=d$ and $\Var(\widehat d\mid F)=4$. This variance is $3$ times the unconditional variance $4/3$, and $\Pr(\widehat d=0\mid F)=1/2$. This is the smallest instance of the duplicate edge pair structure in Section~\ref{sec:inversion}.


\subsection{An AMS observation independent of failure}\label{app:ams}
Each key enters $k$ cells with the same sign. Thus, under random signs, the normalized count sum equals the sign sum in every sample:
\begin{equation}
\frac1k\sum_vC_v=\sum_{i=1}^d\varepsilon_i.\label{eq:ams}
\end{equation}
The right-hand side depends only on the signs, and its square is the standard AMS second-moment estimator with unit weights\cite{alon1999spacecomplexity}. Rearranging~\eqref{eq:dhat} and~\eqref{eq:raw-energy}, we relate $\widetilde d$ and $\widehat d$ through this AMS term:
\begin{equation}
\widetilde d=\Bigl(1-\frac kM\Bigr)\widehat d+\frac kM\Bigl(\frac1k\sum_vC_v\Bigr)^2.\label{eq:raw-decomp}
\end{equation}

\begin{proposition}[AMS on the same table]\label{prop:ams}
The full distribution of $(\frac1k\sum_vC_v)^2$ is independent of the mapping and hence of ideal failure. Specifically, $\frac1k\sum_vC_v$ takes the value $d-2j$ with probability $2^{-d}\binom dj$, $j=0,\ldots,d$. The squared statistic has mean $d$ and variance $2d(d-1)$, which is $M-1$ times the unconditional variance $s(d)^2$ of $\widehat d$. By~\eqref{eq:raw-decomp}, $\widetilde d$ is a weighted sum of this squared statistic and $\widehat d$, with deterministic weights. Given the mapping, the conditional covariance of the squared statistic with $\widehat d$ is $\frac{4}{k(1-k/M)}\sum_{i<j}(H_{ij}-k^2/M)$.
\end{proposition}

\begin{proof}
The right-hand side of~\eqref{eq:ams} depends only on the signs and is independent of $\mathcal G$. It is therefore independent of every event that depends only on the mapping. $\sum_i\varepsilon_i=d-2j$ if and only if exactly $j$ signs are negative, with probability $2^{-d}\binom dj$. $\E(\sum_i\varepsilon_i)^2=d$, and the fourth moment is $3d^2-2d$, so the variance is $2d(d-1)$. \eqref{eq:raw-decomp} follows by directly rearranging~\eqref{eq:dhat} and~\eqref{eq:raw-energy}. For the covariance formula, we expand both statistics as quadratic forms in the signs and use the fact that sign products for distinct unordered pairs are uncorrelated. We multiply $(\sum_i\varepsilon_i)^2-d=2\sum_{i<j}\varepsilon_i\varepsilon_j$ by~\eqref{eq:rs-expand} and take expectations pair by pair. Averaging over failure gives $2d(d-1)c_F$, which is generally nonzero.
\end{proof}

This observation connects the design to classical AMS. The count array is a random-sign linear sketch whose sum gives exactly the AMS sketch. The AMS sketch is independent of failure, so its distribution remains unchanged however rare failure is. This independence comes at a cost in precision: the AMS estimator has variance $M-1$ times that of $\widehat d$. The estimator $\widehat d$ replaces the square of a single sum with the sum of squares over the $M$ cells, reducing the variance to $2d(d-1)/(M-1)$. Its correlation with failure requires the post-failure analysis. The covariance in Proposition~\ref{prop:ams} shows that the two estimators are generally correlated after failure, so we must allocate risk separately when combining them.

We verify Proposition~\ref{prop:ams} using the exhaustive enumeration procedure in Appendix~\ref{app:exact} for the four parameter sets $(4,3,2)$, $(4,3,4)$, $(5,3,3)$, and $(6,3,4)$, covering 144,040 states. For every mapping, $(\frac1k\sum_vC_v)^2$ has the same distribution, with mean $d$ and variance $2d(d-1)$.

\section{Hypergraph lemmas}\label{app:graph}
This appendix proves the hypergraph results used by both Theorem~\ref{thm:clt} and the discrete limit distribution: the small-core counting, the fact that failure is determined by a duplicate edge pair when $d=o(M)$, the low-load lemma at fixed subcritical load, and the concentration of the overlap graph. These results depend only on the unsigned mapping $\mathcal G$, where $\mathcal G$ consists of $d$ labeled independent uniform $k$-subsets, with duplicate edges allowed.

Fix $k\ge3$. Let $\lambda=\binom d2/\binom Mk$ be the expected number of duplicate edge pairs, let $Q$ be the planted distribution that first picks a pair of edge labels uniformly and then conditions on this pair mapping identically, and let $b_k=\lceil3(k-2)/2\rceil$. Write $K=\binom Mk$, let $N=\sum_{i<j}\mathbf1_{\{S_i=S_j\}}$ be the number of duplicate edge pairs, and let $D=\{N\ge1\}$.

\subsection{Small-core counting and duplicate edge pairs}
\begin{lemma}[Small-core counting]\label{lem:smallcore}
A 2-core subgraph with $e$ edges uses at most $v_e=\lfloor ke/2\rfloor$ vertices. For $d\le A_0M$ and $e\le\varepsilon M$ (with $\varepsilon<1/k$), the first-moment upper bound on its occurrence probability satisfies
\begin{equation*}
B_e=\binom de\binom M{v_e}\Bigl(\binom{v_e}k\Big/\binom Mk\Bigr)^e\le\bigl[C(e/M)^{(k-2)/2}\bigr]^e,
\end{equation*}
where $C$ depends only on $k$ and $A_0$. When the factor in $d$ is kept, $B_e\le(C_kd/M)^e(e/M)^{\lceil(k-2)e/2\rceil}$.
\end{lemma}

The proof of Theorem 1 in \cite{gm} uses the same first-moment counting to obtain a failure probability $O(t^{-k+2})$; the analogous counting of stopping sets appears in \cite{yw,orlitsky2005stopping}. Here we keep the factor in $d$, for use by Lemma~\ref{lem:plant} when $d=o(M)$.

\begin{proof}
For a fixed set $I$ of $e$ edge labels and a set $V$ of $v_e$ vertices, the event $A_{I,V}=\{S_i\subseteq V,\ i\in I\}$ contains every core lying in $V$. Summing over $I$ and $V$ gives the first moment $B_e$ in the lemma. Using $\binom de\le(3d/e)^e$, $\binom M{v_e}\le(3M/v_e)^{v_e}$, $\binom{v_e}k/\binom Mk\le(v_e/M)^k$, together with $ke-v_e=\lceil ke/2\rceil$ and $v_e\le ke/2$,
\[
B_e\le3^{e+v_e}\Bigl(\frac de\Bigr)^e\Bigl(\frac{v_e}M\Bigr)^{\lceil ke/2\rceil}\le\Bigl(\frac{C_kd}M\Bigr)^e\Bigl(\frac eM\Bigr)^{\lceil(k-2)e/2\rceil},
\]
which is the bound that keeps the factor in $d$. Substituting $d\le A_0M$ and dropping the ceiling since $e/M<1$ gives the first bound. Let $a=(k-2)/2$; keeping the ceiling at $e=3$ gives $B_3=O(M^{-b_k})$. For $e\ge4$, the ratio of consecutive terms is at most $Ce^a((e+1)/M)^a$. For fixed $r\ge0$, choose $\varepsilon$ small enough that the ratio of consecutive terms of $e^rB_e$ is at most $1/2$ on $4\le e<\varepsilon M$, so that
\[
\sum_{e=3}^{\lfloor\varepsilon M\rfloor}e^rB_e=O(M^{-b_k}).
\]
This sum also covers the intermediate range with $e\to\infty$ and $e=o(M)$.
\end{proof}


The following lemma is the quantitative form of Proposition~\ref{prop:plant}: when $d=o(M)$, failure is almost always caused by a pair of duplicate edges, so the failure-conditioned mapping is close to the planted distribution $Q$ in TV. Duplicate edge pairs are the bottleneck of IBLT failure, a mechanism pointed out in \cite{gm}; $Q$ has density $N/\lambda$ with respect to the original distribution, the size-biased distribution of $N$ \cite{arratia2019sizebias}.

\begin{lemma}[failure is determined by a duplicate edge pair when $d=o(M)$]\label{lem:plant}
Let $2\le d=o(M)$; $d\to\infty$ is not required. Then
\begin{align*}
\Pr(F\setminus\{\text{there is a duplicate edge pair}\})&\le C_k\Bigl(\frac{d^3}{M^{b_k+3}}+\frac{d^4}{M^{2k}}\Bigr)=o(\lambda),\\
\bigl\|\mathcal L(\mathcal G\mid F)-Q\bigr\|_{TV}&\le C_k\Bigl\{\lambda+dM^{k-b_k-3}+\frac{d^2}{M^k}\Bigr\}\to0.
\end{align*}
\end{lemma}

\begin{proof}
A single edge cannot form a core; two edges form a core if and only if they map identically. Hence $D\subseteq F$, and on $F\setminus D$ the core has at least three edges. Since $d=o(M)$, eventually $d\le\varepsilon M$, so every core of size $e\le d$ falls within the range of Lemma~\ref{lem:smallcore}; use the bound that keeps the factor in $d$. The $e=3$ term is $O_k(d^3M^{-3-b_k})$. For $e\ge4$, use the envelope without the ceiling $(C_kd/M)^e(e/M)^{(k-2)e/2}$, whose ratio of consecutive terms is at most $C_k(d/M)((e+1)/M)^{(k-2)/2}$, tending to zero uniformly on $e\le d=o(M)$, so $\sum_{e=4}^dB_e\le C_kd^4/M^{2k}$. Adding the two gives the first bound. From $\lambda\asymp_kd^2/M^k$ and $b_k\ge k-1$, the ratio of the first bound to $\lambda$ is $O_k(dM^{k-b_k-3}+d^2/M^k)\to0$.

A specified pair of edges maps identically with probability $1/K$, so $\E N=\lambda$. Two distinct edge pairs are simultaneously duplicated with probability $K^{-2}$: if they share an edge, three edges map identically; if they do not, the two pairs are independent. Hence $\E[N(N-1)]\le\lambda^2$, and from $\mathbf1_D\ge N-N(N-1)/2$ we get $\lambda-\lambda^2/2\le\Pr(D)\le\lambda$. The density of $Q$ with respect to the original distribution is exactly $N/\lambda$; by the triangle inequality,
\[
\bigl\|\mathcal L(\mathcal G\mid D)-Q\bigr\|_{TV}\le\frac12\E\Bigl|\frac{\mathbf1_D}{\Pr(D)}-\frac{\mathbf1_D}{\lambda}\Bigr|+\frac1{2\lambda}\E(N-\mathbf1_D)=\frac{\lambda-\Pr(D)}{\lambda}\le\frac\lambda2.
\]
Since $D\subseteq F$, $\|\mathcal L(\mathcal G\mid F)-\mathcal L(\mathcal G\mid D)\|_{TV}=\Pr(F\setminus D)/\Pr(F)\le\Pr(F\setminus D)/\Pr(D)$. As $\lambda\to0$, $\Pr(D)\ge\lambda/2$, and by the first bound this term is $O_k(dM^{k-b_k-3}+d^2/M^k)$. Adding the two bounds gives the second bound. The proof uses only the small-core counting, so it holds for all $2\le d=o(M)$.
\end{proof}

\subsection{failure structure at fixed subcritical load}
\begin{lemma}[Low load]\label{lem:propL}
Let $d/M\to\alpha\in(0,\alpha_c)$. Then $\Pr(F\setminus P_2)=O(M^{-b_k})$, $\Pr(F)=\lambda\{1+O(M^{-\delta_k})\}$, and $\|\mathcal L(\mathcal G\mid F)-Q\|_{TV}=O(M^{-\delta_k})$, where $P_2$ denotes the event that the actual core consists of exactly two duplicate edges and $\delta_k=\min\{k-2,\,b_k-(k-2)\}\ge1$.
\end{lemma}

\begin{proof}
We first show that for every fixed $\varepsilon>0$ and $\beta>0$, $\Pr(|\mathrm{core}|\ge\varepsilon M)=o(M^{-\beta})$. This is a variant of a known result: Kim \cite{kim2008poisson}, Theorem 1.7, proves for the binomial hypergraph $H(n,p;k)$ that $\Pr(|\mathrm{core}|\ge\delta n)$ is exponentially small at subcritical load, and \cite{gm} also states the same conclusion citing \cite{dn08}, whose model assumes a given bounded degree distribution. The mapping in this paper allows duplicate edges, and the probability of a duplicate edge at linear load is only polynomially small, so we prove the required bound directly by message passing of fixed depth, following the finite-round neighborhood analysis and concentration used for LDPC codes \cite[Theorem~2]{richardson2001capacity}, and handle the cell degree, which has no fixed upper bound in this model, separately. Let $Z_T$ be the number of edges all of whose incoming messages are $1$ at round $T$; then $|\mathrm{core}|\le Z_T$ on every graph. Local convergence to a Poisson branching tree gives $\E Z_T/M\to\alpha q_T^k$, where $q_0=1$ and $q_{t+1}=1-\exp(-k\alpha q_t^{k-1})$. When $\alpha<\alpha_c$, this recursion has no positive fixed point, so one can choose $T$ with $\alpha q_T^k<\varepsilon/4$. On the set $H$ of graphs whose maximum degree does not exceed $\lceil\log M\rceil$, changing $h$ edge coordinates changes $Z_T$ by at most $L_Mh$, where $L_M$ is of polylogarithmic order. By the typical bounded differences inequality \cite[Proposition~2]{combes2015mcdiarmid}, $\Pr(Z_T\ge\varepsilon M)\le\Pr(H^c)+\exp\{-cM/L_M^2\}=o(M^{-\beta})$.

On $F\setminus P_2$, the actual core has at least three edges. The part of size below $\varepsilon M$ is controlled by Lemma~\ref{lem:smallcore}, and the rest by the linear core bound above, giving $\Pr(F\setminus P_2)=O(M^{-b_k})$. The remaining steps are the same as the second paragraph of the proof of Lemma~\ref{lem:plant}: the Bonferroni inequality gives $\Pr(F)=\lambda+O(\lambda^2+M^{-b_k})$; $\|\mathcal L(\mathcal G\mid D)-Q\|_{TV}\le\lambda/2$; and since $P_2\subseteq D\subseteq F$, $\|\mathcal L(\mathcal G\mid F)-\mathcal L(\mathcal G\mid D)\|_{TV}\le\Pr(F\setminus P_2)/\Pr(D)$. From $\lambda\asymp_kM^{2-k}$, the two terms combine into $O(M^{-\delta_k})$. This is a total variation bound on the entire labeled mapping.
\end{proof}


\subsection{Superpolynomial concentration of the overlap graph}
\begin{lemma}[Superpolynomial graph concentration]\label{lem:mdconc}\label{lem:conc}
Fix $k\ge3$, and let $d/M$ lie in a fixed compact interval of $(0,\infty)$. Write $m_{ij}=H_{ij}$, $\Omega=\sum_{i<j}m_{ij}$, $\Omega_2=\sum_{i<j}m_{ij}^2$, $\ell_i=l_i=\sum_{j\ne i}H_{ij}$, and $\Omega_\ell=\sum_i\ell_i^2$. There is a deterministic set of graphs $\mathcal G_M$, whose complement has probability smaller than any fixed negative power of $M$, on which these three statistics deviate from their exact means by at most $O(\sqrt M(\log M)^4)$, and
\begin{equation*}
\Omega_2=\binom d2\left\{\frac{k^2}M+\frac{k^2(k-1)^2}{M(M-1)}\right\}+O(\sqrt M(\log M)^4),\qquad
\max_i l_i\le k\lceil\log M\rceil.
\end{equation*}
\end{lemma}

\begin{proof}
Write $D_v$ for the degree of cell $v$, let $L=\lceil\log M\rceil$ and $\mathcal D_L=\{\max_vD_v\le L\}$. Since $D_v\sim\mathrm{Bin}(d,k/M)$, $\Pr(\mathcal D_L^c)\le M(eC/L)^L$, which is superpolynomially small.

\emph{Lipschitz bound.} For two graphs in $\mathcal D_L$, successively replace the edges with differing labels by the edges of the other graph. The incidence set of every intermediate graph is contained in the union of the two endpoint graphs, so no vertex degree exceeds $2L$. Replacing one edge in such an intermediate graph affects only the $O(L)$ edge labels overlapping the old or new edge, and each $m_{ij}\le k$, so $\Omega$ and $\Omega_2$ each change by $O(L)$. For the replaced edge $\ell_i=O(L)$, whose square changes by $O(L^2)$; the $\ell_j$ of the other $O(L)$ affected labels change by $O(1)$, each square changing by $O(L)$. Thus $\Omega_\ell$ changes in total by $O(L^2)$. Hence all three are $CL^2$-Lipschitz in the Hamming distance over edge labels on $\mathcal D_L$.

\emph{Concentration.} For functions that are Hamming Lipschitz on the good set, the typical bounded differences inequality \cite{warnke2016typical} gives concentration; here we use the form of Combes \cite[Proposition~2]{combes2015mcdiarmid}: extend $f$ from $\mathcal D_L$ as a Lipschitz extension, truncate it to a deterministic nonnegative polynomial range, and then apply McDiarmid's inequality over the independent edges. The bad set has superpolynomially small probability, and both $f$ and its truncated extension are at most $O(M^3)$, so the change in the expectation is also superpolynomially small; hence
\[
\Pr\bigl(|f-\E f|>C\sqrt M(\log M)^4\bigr)\le\Pr(\mathcal D_L^c)+2e^{-c(\log M)^4},\qquad f=\Omega,\Omega_2,\Omega_\ell.
\]

\emph{Deterministic means.} For the overlap $H$ of independent uniform $k$-subsets, write $h_1=\E H=k^2/M$ and $h_2=\E H^2=k^2/M+k^2(k-1)^2/[M(M-1)]$. Exactly,
\[
\E\Omega=\tbinom d2h_1,\qquad\E\Omega_2=\tbinom d2h_2,\qquad\E\Omega_\ell=d\{(d-1)h_2+(d-1)(d-2)h_1^2\}.
\]
The last identity holds because, given $S_i$, the $m_{ij}$ are independent and their conditional first and second moments do not depend on $S_i$. Taking the good set as the intersection of $\mathcal D_L$ with the three deviation events gives the stated concentration bound. Since $l_i\le kL$, the lemma follows.
\end{proof}

\section{The Gaussian limit after failure and the Gaussian rule}\label{app:clt}
This appendix proves Theorem~\ref{thm:clt} and Corollary~\ref{cor:gauss}. Write $D_v$ for the degree of cell $v$ and $l_i=\sum_{j\ne i}H_{ij}$. The structure of the proof is as follows: first prove a Gaussian limit for the quadratic form in the random signs on a deterministic good graph (Lemma~\ref{lem:cumulant}), uniformly over good graphs; then show that the good graphs still carry almost all of the probability after conditioning on failure; for the $d=o(M)$ and overload regimes use the Gaussian limit under deterministic coefficients (Lemma~\ref{lem:plantclt}), which is uniform in the coefficients and holds under both the random mapping and the planted distribution; finally branch on the limit of $d/M$ and assemble the pieces by a subsequence argument.

\subsection{CLT for the random signs on a good graph}
\begin{lemma}\label{lem:cumulant}
Fix $k\ge3$. Let a deterministic sequence of graphs satisfy $\Omega_2/M\to\beta\in(0,\infty)$ and $\max_il_i=O(\log M)$. Then $\Omega_2^{-1/2}\sum_{i<j}H_{ij}\varepsilon_i\varepsilon_j\Rightarrow N(0,1)$, and the convergence is uniform over graphs satisfying the same set of constant bounds.
\end{lemma}

\begin{proof}
There are at most $\Omega_2=O(M)$ nonzero summands, each bounded by $k$. Two summands are dependent only when they share a key, so the maximum degree of the dependency graph is at most $2\max_il_i=O(\log M)$. The covariance of distinct summands is zero (Appendix~\ref{app:encoding}), so the variance of the sum is exactly $\Omega_2$. For a fixed order $r\ge3$, a joint cumulant is nonzero only when the group of summands is connected in the dependency graph \cite{janson1988cumulant}; the number of connected groups is $O_r(M(\log M)^{r-1})$, and the cumulant of each group is bounded. Dividing by $\Omega_2^{r/2}\asymp M^{r/2}$ gives $O_r(M^{1-r/2}(\log M)^{r-1})\to0$. The first cumulant is zero and the second is one; the standard Gaussian is determined by its moments, and the method of moments gives convergence. These bounds depend only on the range of $\beta$ and the constant in $\max_il_i$, so they are uniform over graphs satisfying the same set of constant bounds; if the convergence of the CDF were not uniform, one could extract a deterministic good graph along which convergence fails, contradicting the argument above.

The difference between $\widehat d-d$ and $\frac{2}{k(1-k/M)}$ times this sum is
\begin{equation*}
-\frac{k}{M(1-k/M)}\Bigl[\Bigl(\sum_i\varepsilon_i\Bigr)^2-d\Bigr],
\end{equation*}
which depends only on the signs and has $L^2$ norm $O(d/M)$, that is $O(M^{-1/2})$ after dividing by $s(d)\asymp d/\sqrt M$, for every graph. On the good graph of Lemma~\ref{lem:conc}, $\Omega_2$ equals its exact mean plus $O(\sqrt M(\log M)^4)$, so the ratio of $\frac{4\Omega_2}{k^2(1-k/M)^2}$ to $s(d)^2$ tends to one, and the normalization by $\sqrt{\Omega_2}$ in this lemma can be replaced by normalization by $s(d)$. Hence the standardized $X$ of $\widehat d$ has the same uniform Gaussian limit on the good graph. By \eqref{eq:raw-decomp}, the difference between $\widetilde d$ and $\widehat d$ is $\frac kM[(\sum_i\varepsilon_i)^2-\widehat d]$. Given the graph, $(\sum_i\varepsilon_i)^2-d$ has $L^2$ norm $\sqrt{2d(d-1)}$; $\widehat d$ is conditionally unbiased with conditional variance at most $(M-1)s(d)^2=2d(d-1)$. By the triangle inequality, the conditional $L^2$ norm of the difference is at most $\frac{2k}M\sqrt{2d(d-1)}$, which is at most $2k\sqrt{M-1}/M\to0$ after dividing by $s(d)$, so the same holds for $\widetilde d$.
\end{proof}

\subsection{failure conditioning at positive load ratio}
For $d\ge2$, a specified pair of duplicate edges already forms a core, so
\begin{equation*}
\Pr(F)\ge\Pr(S_1=S_2)=\binom Mk^{-1}.
\end{equation*}
By Lemma~\ref{lem:conc}, $\Pr(\mathcal G\notin\mathcal G_M\mid F)\le\Pr(\mathcal G\notin\mathcal G_M)\binom Mk\to0$. The error of the conditional CDF is at most the uniform error on the good graph plus this conditional exceptional probability, so at positive load ratio
\begin{equation*}
\sup_x\bigl|\Pr(X\le x\mid F)-\Phi(x)\bigr|\to0.
\end{equation*}
The same argument applies to any event that depends only on the mapping and has probability at least some fixed negative power of $M$. What is divided by $\Pr(F)$ here is the superpolynomially small exceptional probability; the CLT error on the good graph is simply averaged after conditioning, so the proof does not need a CLT convergence rate strong enough to overcome a rare failure. A typical floor failure is caused by a pair of duplicate edges, which only sets one $H_{ij}$ to $k$; the change in $\Omega_2$ is $O(1)$, far smaller than the deviation allowed by Lemma~\ref{lem:conc}, so it lies inside $\mathcal G_M$.

The conditional variance also tends to $s(d)^2$: on the good graph the ratio tends to one by \eqref{eq:condmoments} and Lemma~\ref{lem:conc}; on any graph the ratio is at most $M-1$, and the contribution of bad graphs is controlled by a superpolynomial probability.

\subsection{Gaussian limit under deterministic coefficients}
The following Gaussian limit holds uniformly over any deterministic coefficient vector $(s_i)\in\{\pm1\}^d$, so it applies to random-sign encoding by using it given iid coding coefficients and then averaging over the signs.

\begin{lemma}[Gaussian limit under deterministic coefficients]\label{lem:plantclt}
For deterministic coefficients $(s_i)\in\{\pm1\}^d$, let $C_v=\sum_{i:v\in S_i}s_i$, compute $\widehat d$ by \eqref{eq:dhat}, and obtain $X$ by \eqref{eq:X}. If $d/\sqrt M\to\infty$, then $X\Rightarrow N(0,1)$, with convergence uniform in the coefficients and without requiring $d/M$ to have a finite upper bound. Under the planted distribution $Q$, we likewise have $X\Rightarrow N(0,1)$, uniformly in the coefficients and the label pair.
\end{lemma}

\begin{proof}
Write $Y=\widehat d-d$, $W_{ij}=H_{ij}-k^2/M$, and $v_i=\mathbf1_{S_i}-(k/M)\mathbf1$. Exposing the edges in the order of their labels, the martingale difference is $\Delta_j=\frac{2}{k(1-k/M)}s_jv_j^\top\sum_{i<j}s_iv_i$, with $Y=\sum_j\Delta_j$. The sum of conditional variances is
\begin{gather*}
\mathcal V_d=s(d)^2+\frac{8}{k(1-k/M)(M-1)}\sum_{i<\ell}(d-\ell)s_is_\ell W_{i\ell},\\
\Var\mathcal V_d=\frac{64}{(M-1)^3}\sum_{\ell=2}^d(\ell-1)(d-\ell)^2=O\Bigl(\frac{d^4}{M^3}\Bigr).
\end{gather*}
Also $s(d)^4\approx4d^4/M^2$, so the variance of the normalized bracket $\mathcal V_d/s(d)^2$ is $O(1/M)$, converging to one in $L^2$. Given the first $j-1$ edges, the components of the vector $\sum_{i<j}s_iv_i$ sum to zero, so $v_j^\top\sum_{i<j}s_iv_i$ equals the sum of its components over the $k$ cells of $S_j$, whose fourth conditional moment is at most $k^3\cdot\frac kM$ times the sum of the fourth powers of all components. The component on cell $v$, $\sum_{i<j}s_i(\mathbf1\{v\in S_i\}-k/M)$, is a sum of $j-1$ independent, centered, bounded terms, each with variance at most $k/M$; by Rosenthal's inequality \cite{rosenthal1970}, its fourth moment is at most $C_k(j/M+j^2/M^2)$. Hence $\E\Delta_j^4\le C_k(j/M+j^2/M^2)$; summing over $j\le d$ and dividing by $s(d)^4$, the fourth-order Lyapunov remainder satisfies
\begin{equation*}
\frac{\sum_j\E\Delta_j^4}{s(d)^4}\le C_k\Bigl(\frac M{d^2}+\frac1d\Bigr)\to0,
\end{equation*}
and by the martingale CLT \cite{hall1980martingale}, $X\Rightarrow N(0,1)$. Both bounds depend only on $|s_i|=1$, so the convergence is uniform in the coefficients, and no boundedness of $d/M$ is needed.

Under the planted distribution, the statistic is invariant under label permutation, so we may fix the duplicate edge pair as 1, 2, with $S_1=S_2$. Write $Y_{\ge3}$ for $\widehat d$ computed from the remaining $d-2$ edges alone, minus $d-2$; then
\begin{equation*}
Y=Y_{\ge3}+2s_1s_2+\frac{2}{k(1-k/M)}(s_1+s_2)\sum_{j\ge3}s_j\Bigl(|S_1\cap S_j|-\frac{k^2}{M}\Bigr).
\end{equation*}
The last term has variance $O(d/M)$. The deterministic term $2s_1s_2$ and the last term, divided by $s(d)$, have mean squares $O(M/d^2)$ and $O(1/d)$ respectively, both tending to zero; $Y_{\ge3}$ is handled by the unconditional result. Combining with the total variation bound of Lemma~\ref{lem:plant} or \ref{lem:propL} returns to the failure-conditional distribution. The $d=o(M)$ and overload branches of the next subsection use exactly this conclusion.
\end{proof}

\subsection{\texorpdfstring{Proof of Theorem~\ref{thm:clt} (the Gaussian limit after failure) and Corollary~\ref{cor:gauss}}{Proof of the Gaussian limit after failure and the Gaussian rule}}
Assume $d/\sqrt M\to\infty$. For any subsequence, take a further subsequence along which $d/M$ converges in $[0,\infty]$.

When the limit is a positive finite value, use the previous subsection.

When the limit is zero, Lemma~\ref{lem:plant} gives that the TV distance between the failure-conditional distribution and the planted distribution $Q$ tends to zero, valid for all $2\le d=o(M)$; it is a purely graph-theoretic result, which still holds after adding independent signs. The conclusions of Lemma~\ref{lem:plantclt} under the random mapping and the planted distribution are both uniform in deterministic signs: the fourth-order Lyapunov remainder of the edge-exposure martingale is $O_k(M/d^2+1/d)$, and the normalized mean square of the correction term from the planted duplicate edge pair is also $O_k(M/d^2+1/d)$, both tending to zero when $d\gg\sqrt M$. Applying these uniform conclusions given iid coding coefficients and averaging over the signs yields the conditional Gaussian limit.

When the limit is infinite, eventually $d>M$. Each edge deleted in the full peeling can be assigned a cell of degree one at the moment of deletion; that cell is subsequently empty and cannot serve again as a witness for a deletion, so full peeling requires $d\le M$. In this branch $F$ occurs with certainty, and we directly use the conclusion of Lemma~\ref{lem:plantclt} under the random mapping.

All three branches give uniform convergence of the conditional CDF to $\Phi$, and a subsequence contradiction gives \eqref{eq:clt} for the full sequence. Corollary~\ref{cor:gauss} follows from the fact that the Gaussian quantile and the standardized chi-square quantile have the same limit, together with $\{U<d\}\subseteq\{\widehat d<t(d)\}$. \qed

\subsection{\texorpdfstring{Monotonicity of the acceptance lower bound $t(n)$ in Definition~\ref{def:gauss}}{Monotonicity of the acceptance lower bound of the Gaussian rule}}\label{app:gauss-mono}
Here $t(n)=n+q\,s(n)$ with $q=[\chi^2_{M-1,u}-(M-1)]/\sqrt{2(M-1)}$. When $q\ge0$, $t$ is clearly increasing. When $q<0$, differentiating for $n\ge2$, $\frac{d}{dn}s(n)\le\frac{\sqrt2(2n-1)}{2\sqrt{n(n-1)(M-1)}}\le\frac{3}{2\sqrt2}\cdot\frac{\sqrt2}{\sqrt{M-1}}$, so a sufficient condition is $|q|\sqrt{2/(M-1)}<\sqrt8/3$. The Chernoff bound for the chi-square lower tail \cite[Lemma~1]{laurent2000adaptive} gives $|q|\le\sqrt{2\log(1/u)}\le3.035$ (for $u\ge0.01$), and the condition holds when $M\ge64$. \qed

\section{Proofs of the discrete limit distribution and the Skellam rule}\label{app:inversion}
We prove Theorems~\ref{thm:cross}, \ref{thm:uniform}, and \ref{thm:coverage}. Write the iid coding coefficients as $\varepsilon_i$. By \eqref{eq:raw-expand}, $\widetilde d-d=\frac2k\sum_{i<j}H_{ij}\varepsilon_i\varepsilon_j$ holds for every sample. The final subsection discusses an extension with multiple sign counts per cell.

\subsection{Exact moments for a specified duplicate edge pair}
Let $S_1=S_2$, with the remaining $d-2$ edges independent and uniform. If the two coding coefficients have opposite signs, the edges cancel completely, leaving the remaining $d-2$ elements. If they are equal, the edges combine into a single edge with coefficient $\pm2$. For $d\ge2$,
\begin{equation*}
\E[\widehat d\mid S_1=S_2]=d,\qquad
\Var(\widehat d\mid S_1=S_2)=\frac{2d(d-1)}{M-1}+4-\frac4{M-1}.
\end{equation*}
The component with opposite coefficients has variance $2(d-2)(d-3)/(M-1)$, and the component with equal coefficients has variance $2(d-2)(d+5)/(M-1)$. Averaging these variances and adding $4$ from the difference between the component means $d\mp2$ gives the result.

\subsection{From failure conditioning to the planted distribution}
Lemma~\ref{lem:plant} gives $\|\mathcal L(\mathcal G\mid F)-Q\|_{TV}\to0$ for all $2\le d=o(M)$. Here $Q$ selects a pair of edge labels uniformly and then conditions on the two edges having the same mapping. For $2\le d\le C\sqrt M$, both $\lambda$ and $d^2/M^k$ are $O_k(M^{1-k})$. Since $b_k=\lceil3(k-2)/2\rceil$, we have $k-b_k-3=-\lceil k/2\rceil\le-2$ and hence $dM^{k-b_k-3}=O_k(d/M^2)$. The distance is therefore at most $O_k(d/M^2+M^{1-k})$. Adding independent signs and applying \eqref{eq:raw-expand} preserve this TV bound. The statistic is invariant under label permutation, so we may fix the duplicate edge pair as 1, 2. Their common edge and the remaining $d-2$ edges are independent and uniform. For $d=2$, the null distribution is exact under the planted distribution. We treat empty sums below as zero.

\subsection{Three negligible event classes for bounded \texorpdfstring{$c$}{c}}
By the union bound, the common edge of the duplicate edge pair intersects a remaining edge with probability $O_k(d/M)$. Any specified pair of remaining edges shares at least two cells with probability $O_k(M^{-2})$. Summing over all such pairs gives $O_k(d^2/M^2)$. In the intersection graph of the remaining edges, for any three specified edges, the probability that one specified edge intersects both others is
\begin{equation*}
\Bigl(1-\binom{M-k}{k}\Big/\binom Mk\Bigr)^2=O_k(M^{-2}),
\end{equation*}
This equality follows by conditioning on the edge shared by the two intersection tests. The other two edges are independent, and their intersection probability with the shared edge is independent of its position. Summing over ordered triples bounds the probability that the intersection graph is not a matching by $O_k(d^3/M^2)$. Thus, outside an event with probability bounded by \eqref{eq:tv-error}, the common edge of the duplicate edge pair is disjoint from the remaining edges, and all overlaps among the remaining edges are single-cell overlaps forming a matching.

\subsection{Poisson approximation for overlaps among the remaining edges}
For each pair of remaining edges, consider the indicator that they intersect. These indicators have a dependency graph defined by shared edge labels. Collections of indicators with disjoint label sets are independent. Each indicator has mean $p=1-\binom{M-k}k/\binom Mk=k^2/M+O_k(M^{-2})$ and $O(d)$ neighbors. By the calculation in the previous subsection, two indicators sharing an edge have joint probability exactly $p^2$. We apply the local dependence bound of Arratia, Goldstein, and Gordon \cite[Theorem~1]{arratia1989chen}. The neighborhood includes the indicator itself, and indicators outside it depend only on the remaining independent edges, so the third term is zero. The bound is a universal constant times the sum of products of means within neighborhoods plus the sum of joint probabilities over distinct neighbors. Both sums are $O_k((d^2+d^3)/M^2)$ here. Thus the number of intersecting edge pairs among the remaining edges has TV distance $O_k((d^2+d^3)/M^2)$ from $\Pois(\binom{d-2}2p)$. We replace the parameter by $\binom{d-2}2k^2/M=2\lambda_d$. The TV distance between two Poisson distributions is at most the difference between their parameters: we couple the larger-parameter variable to the smaller one by adding an independent Poisson increment, and the variables differ only when this increment is nonzero. The replacement adds an error of $O_k(d^2/M^2)$.

\subsection{Independence of the overlap sign marks and the Skellam limit}
On the matching event, each overlap uses a pair of iid signs shared with no other overlap. Their products are therefore mutually independent and symmetric, and independent of the duplicate edge pair's sign product $\varepsilon_1\varepsilon_2$. Given the mapping, $\widetilde d$ depends only on the number of overlaps among the remaining edges and these independent $\pm1$ marks. To formalize the coupling, we retain the true overlap count and generate that many new independent marks for each mapping. On the matching event, the old and new marked sums have the same conditional distribution. On all other mappings, we couple them arbitrarily, with a distributional error at most the exceptional event's probability. We then apply the same marking kernel to the Poisson approximation of the count, preserving the TV bound. By Poisson thinning \cite{kingman1993poisson}, $\Pois(2\lambda_d)$ independent $\pm1$ marks split by sign into two independent $\Pois(\lambda_d)$ counts. Their difference, $Z_d$, has a Skellam distribution \cite{skellam1946}. Adding the duplicate edge pair's contribution $2\varepsilon_1\varepsilon_2$ gives $J_{\lambda_d}$ in \eqref{eq:J}. Thus $\widetilde d$ is close to the null distribution $d+J_{\lambda_d}$ in \eqref{eq:aux}. We use the Poisson approximation for the unconditional count and the matching event to bound the coupling error.

The four steps give a total error of $O_k(d/M+d^2/M^2+d^3/M^2+d/M^2+M^{1-k})$. Since $d\le M$ and $k\ge3$, the terms $d/M$ and $d^3/M^2$ absorb all others, giving \eqref{eq:tv-error}. When $d/\sqrt M\to c$, we have $\lambda_d\to k^2c^2/4$. Continuity in the Poisson parameter gives convergence of $J_{\lambda_d}$ in total variation to $J_{k^2c^2/4}$, proving Theorem~\ref{thm:cross}(ii). The $\pm2$ and Skellam terms are independent, so $\Var J_\lambda=4+8\lambda/k^2$ and the limiting variance is $4+2c^2$. For (iii), as $\lambda\to0$, we have $\Pr(Z\ne0)\le2\lambda\to0$ and $J_\lambda\Rightarrow2\varepsilon_1\varepsilon_2$. As $\lambda\to\infty$, we have $Z/\sqrt{2\lambda}\Rightarrow\mathcal N(0,1)$ and $\sqrt{8\lambda}/k=\sqrt2\,c$. The $\pm2$ term divided by $\sqrt2\,c$ tends to zero, so $J_{k^2c^2/4}/(\sqrt2\,c)\Rightarrow\mathcal N(0,1)$. The estimator $\widehat d$ has the same weak limit: the AMS term in \eqref{eq:raw-decomp} has coefficient $k/M$, and $(\sum_i\varepsilon_i)^2=O_p(d)=O_p(\sqrt M)$. \qed

\subsection{Sparse limit for fixed small \texorpdfstring{$d$}{d}}
\begin{proposition}[Limit when $d=o(\sqrt M)$]\label{prop:smalld}
Fix $k\ge3$ and let $2\le d=o(\sqrt M)$. Conditional on failure, both $\widehat d-d$ and $\widetilde d-d$ converge weakly to the distribution assigning equal probability to $\pm2$. The Gaussian rule's conditional miss rate tends to zero for $d=2,3$, and to $1/2$ for fixed $d\ge4$ or $d\to\infty$.
\end{proposition}

\begin{proof}
Theorem~\ref{thm:cross} gives the equally weighted $\pm2$ limit for the fluctuation of $\widetilde d$ when $d=o(\sqrt M)$. By \eqref{eq:raw-decomp}, the correction for $\widehat d$ is $O_p(d^2/M)=o_p(1)$, so its fluctuation has the same weak limit. Consider the event that the remaining edges are pairwise disjoint and disjoint from the duplicate edge pair. On this event, $\widetilde d$ equals $d+2\varepsilon_1\varepsilon_2$ exactly, and the event has conditional probability tending to one. The Gaussian rule has $t(d)=d-O(s(d))$ with $s(d)\to0$. The negative component leads to rejection of the true candidate, and the positive component leads to acceptance. For $d=2,3$, the negative component gives an observation below $t(2)$, so the rule returns $+\infty$, which is not a miss. For $d\ge4$, the negative component still leads to acceptance of candidate $d-2$, causing a miss. For $d=2$, failure is equivalent to a duplicate edge pair, and $\Pr(\widehat d=0\mid F)=\Pr(\widehat d=4\mid F)=1/2$ for every $M>k$. Thus multiplying $\widehat d$ by any fixed constant to obtain an upper bound leaves a post-failure risk of at least $1/2$.
\end{proof}

\subsection{Proof of the uniform CDF approximation}
\begin{proof}
First consider $d/\sqrt M\to\infty$. Appendix~\ref{app:clt} gives a conditional Gaussian limit for $\widetilde d$ standardized by $s(d)$. The null distribution is the distribution of $d+J_{\lambda_d}$. As $\lambda_d\to\infty$, we can express $Z_d$ in the Skellam term as a sum of iid increments whose number grows with $\lambda_d$. The CLT gives $(2/k)Z_d/\sqrt{8\lambda_d/k^2}\Rightarrow\mathcal N(0,1)$. The ratio of $8\lambda_d/k^2=2(d-2)(d-3)/M$ to $s(d)^2$ tends to one, and the $\pm2$ term divided by $s(d)$ tends to zero. Hence $J_{\lambda_d}/s(d)\Rightarrow\mathcal N(0,1)$. Both standardized CDFs converge to the continuous CDF $\Phi$. Weak convergence to a continuous limit implies uniform CDF convergence (P\'olya's theorem), so their Kolmogorov distance tends to zero.

To prove uniformity, suppose that there exist $\delta>0$ and a sequence $(M,d)$ along which the distance is at least $\delta$. Take a subsequence with $d/\sqrt M\to c\in[0,\infty]$. If $c<\infty$, this subsequence eventually satisfies $d\le C\sqrt M$ for some constant, and Theorem~\ref{thm:cross} gives a TV distance tending to zero. If $c=\infty$, the previous paragraph gives convergence to zero. Both cases contradict the lower bound $\delta$.
\end{proof}

\subsection{Proof of the conditional miss rate of the Skellam rule}
\begin{proof}
Write $G_n$ for the CDF of the null distribution at candidate $n$ and $Y$ for the observation. The rule rejects $n$ if and only if $G_n(Y)\le u$. We use the standard argument for test inversion and the conservativeness of a discrete p-value \cite{neyman1937}. If $Y\sim G_n$, then $\Pr(G_n(Y)\le u)\le u$. If no lattice point satisfies $G_n(y)\le u$, the rejection event is empty. Otherwise, let $y^*$ be the largest lattice point satisfying $G_n(y)\le u$. The rejection event is then $\{Y\le y^*\}$, with probability $G_n(y^*)\le u$. Since the rejection region $\{Y\le y^*\}$ is a lower set, the rejection probabilities under the true distribution and $G_n$ differ by at most their Kolmogorov distance. Taking $n=d$ and applying Theorem~\ref{thm:uniform}, we obtain a rejection probability of at most $u+o(1)$ for the true candidate, uniformly over all $d\ge2$. The event $\{U<d\}$ implies rejection of the true candidate, giving \eqref{eq:coverage}. The argument applies to arbitrary acceptance sets and allows candidate rejection probabilities to vary nonmonotonically with $n$.
\end{proof}

Using $\le u$ as the rejection criterion with the left limit $G_n(Y-)=\Pr(Y^{\prime}<Y)$ (where $Y^{\prime}\sim G_n$) gives a different rule. The equally weighted $\pm2$ distribution has left limit zero at $-2$. This rejects its smallest support point, which has probability $1/2$, and the conditional miss rate tends to $1/2$ when $d\ge4$ and $d=o(\sqrt M)$. At that point, $G_n$ equals $1/2>u$, so the Skellam rule accepts the candidate.

For $d=3$, the observation may lie outside the support of $G_3$, and $\widetilde d\ge1$ always holds. If all three coefficients have the same sign, $\sum_vC_v^2\ge3k$. Otherwise, label the edge whose sign differs from the other two as edge $3$. The sum of squares is $3k+2(H_{12}-H_{13}-H_{23})$. Among the $k$ cells of edge $3$, at least $H_{13}+H_{23}-H_{12}$ are shared with edge $1$ or edge $2$. Thus $H_{13}+H_{23}-H_{12}\le k$, and the sum of squares is at least $k$. Hence $G_3(\widetilde d)\ge G_3(1)=1/2>u$, and the true candidate is always accepted.

\subsection{Combining the two rules using the first-round observation}\label{app:switch}
The combined rule uses the Gaussian rule's upper bound when $\widetilde d$ is at or above the switch point and the Skellam rule's upper bound otherwise. Assume that the ratio of the switch point to $\sqrt{M}$ tends to infinity as $M\to\infty$. We show that the combined rule satisfies \eqref{eq:coverage}.

\begin{proof}
Suppose that there exist $\delta>0$ and a sequence $(M,d)$ along which the combined rule satisfies $\Pr(U<d\mid F)\ge u+\delta$. Take a subsequence with $d/\sqrt M\to c\in[0,\infty]$.

For $c<\infty$, Theorem~\ref{thm:uniform} shows that the distance between the post-failure distribution of $\widetilde d-d$ and the distribution of $J_{\lambda_d}$ tends to zero. Since $\lambda_d$ is bounded, $\widetilde d=d+O_p(1)=O_p(\sqrt M)$. The probability of exceeding the switch point therefore tends to zero. The combined rule differs from the Skellam rule only on this event, so Theorem~\ref{thm:coverage} bounds the miss probability by $u+o(1)$.

For $c=\infty$, a miss implies that the rule selected from the observation rejects the true candidate $d$. The Skellam rule rejects $d$ if and only if $(\widetilde d-d)/s(d)$ is below the lower $u$-quantile of the standardized $G_d$. By the proof of Theorem~\ref{thm:uniform}, $J_{\lambda_d}/s(d)\Rightarrow\mathcal N(0,1)$, so this quantile tends to $\Phi^{-1}(u)$. The Gaussian rule rejects $d$ if and only if $(\widehat d-d)/s(d)<(t(d)-d)/s(d)$. The right-hand side also tends to $\Phi^{-1}(u)$, and $(\widehat d-\widetilde d)/s(d)\to0$ (Appendix~\ref{app:ams}). By Theorem~\ref{thm:clt}, the conditional distribution of the standardized statistic converges to the continuous distribution $\mathcal N(0,1)$. The two events for rejecting $d$ therefore differ only in a boundary region whose probability tends to zero. The combined rule's rejection event is contained in their union. Its probability is thus at most the Skellam rule's rejection probability plus $o(1)$, namely $u+o(1)$.

Both cases contradict $\delta>0$.
\end{proof}

\subsection{Extension to multiple sign counts}\label{app:multi}
Suppose each cell stores $L$ copies of the count, using independent sign hashes and a shared mapping. Conditional on $\mathcal G$, the estimates $\widehat d$ from these copies are mutually independent, each with mean $d$. Averaging independent estimates is the AMS approach to reducing variance \cite{alon1999spacecomplexity}. From $\E[\widehat d^2\mid\mathcal G]\le3d^2$ and $e^{-x}\le1-x+x^2/2$, we obtain
\begin{equation*}
\E\Bigl[\exp\Bigl(-\frac{\widehat d}{5d}\Bigr)\Bigm|\mathcal G\Bigr]\le\frac{43}{50},
\end{equation*}
and the exponential Markov inequality (Chernoff bound) applied to $U=\frac2L\sum_\ell\widehat d^{(\ell)}$ gives $\Pr(U<d\mid\mathcal G)\le(\frac{43}{50}e^{1/10})^L$. This bound holds for every deterministic mapping, hence for arbitrarily rare failure. For $L=272$, the right-hand side is below $10^{-6}$. When only an extra scalar AMS accumulator is stored, the $L$ count sums are independent given $d$, and their joint distribution is a product of binomial distributions. This permits an exact discrete candidate test with the finite-sample guarantee $\Pr(U<d\mid F)\le u$.

Repeating counts reduces absolute error and preserves the ratio of the conditional scale to the unconditional scale. Fix $d\ge4$ and $L$, and let $M\to\infty$. If we apply Gaussian inversion to the average of $L$ copies using the unconditional standard deviation $s(n)/\sqrt L$, the conditional miss rate tends to $2^{-L}\sum_{j<L/2}\binom Lj$. This equals $1/2,1/4,5/16,93/256$ for $L=1,2,4,8$, respectively. Calibration with more copies still requires the correct post-failure distribution, motivating our use of a single count and the discrete limit distribution in the main text.

\section{Exact evaluation and search cutoffs}\label{app:search}
This appendix describes the computation of the Skellam rule (Definition~\ref{def:skellam}): the exact evaluation of $G_n(\widetilde d)$, the statement and proof of two search cutoffs, and the numerical implementation of the exact CDF.

The guarantee of the Skellam rule does not require the accepted candidates to form an interval, and the implementation takes the largest accepted candidate as specified by the definition. The candidates have no theoretical upper bound, so an explicit cutoff is needed: the cutoff of Proposition~\ref{prop:cut11} is about $1.035$ times the observation $\widetilde d$ at $M=16384$ and $u=0.01$. The search checks candidates downward from the cutoff and returns the first accepted one, matching an unbounded scan. The number of candidates requiring an exact CDF and the length of the Poisson window are both proportional to $c$, so the cost of a single sizing is roughly proportional to $\lambda_d$. The Poisson table depends only on $(M,k)$ and the candidate $n$, and can be reused across syncs by the same receiver; measured results appear in Appendix~\ref{app:timing}.

\subsection{\texorpdfstring{Computing $G_n(\widetilde d)$ with integer thresholds}{Computing G_n with integer thresholds}}
Expanding \eqref{eq:aux} according to the two cases $\varepsilon_1\varepsilon_2=\pm1$,
\begin{equation*}
G_n(\widetilde d)=\frac12\Pr\Bigl(Z_n\le\frac k2\bigl(\widetilde d-n-2\bigr)\Bigr)+\frac12\Pr\Bigl(Z_n\le\frac k2\bigl(\widetilde d-n+2\bigr)\Bigr).
\end{equation*}
Let $R=\sum_vC_v^2$ (an integer), so $\widetilde d=R/k$. A random variable following $G_n$ is at most the observation if and only if $Z_n\le\frac k2(R/k-n\mp2)$, that is, $Z_n\le\lfloor(R-kn\mp2k)/2\rfloor$. Hence $G_n(\widetilde d)$ equals
\begin{equation*}
\frac12\Pr\Bigl(Z_n\le\Bigl\lfloor\frac{R-kn-2k}2\Bigr\rfloor\Bigr)+\frac12\Pr\Bigl(Z_n\le\Bigl\lfloor\frac{R-kn+2k}2\Bigr\rfloor\Bigr),
\end{equation*}
where both floors are computed over integers, so floating-point arithmetic does not change the decision at lattice boundaries.

\subsection{Search cutoff from the Chernoff bound}
\begin{proposition}[Chernoff cutoff]\label{prop:cut10}
If $e^{-9M/80}<u$, then every candidate satisfying
\begin{equation}
n\ge\max\Bigl\{4\Bigl(\widetilde d+2\Bigr),\ \frac M{2k},\ 4\Bigr\}\label{eq:cut10}
\end{equation}
is rejected. If $e^{-9M/80}\ge u$, the rule returns $+\infty$ directly.
\end{proposition}

At $M=64$, $e^{-9M/80}\approx7.5\times10^{-4}$, which holds for both $u=0.01$ and $0.05$; the branch returning $+\infty$ disappears as $M\to\infty$ for any fixed $u$, and does not affect Theorem~\ref{thm:coverage}.

\begin{proof}
The lower tail of the null distribution is governed by the component with shift $-2$, so it suffices to show that this component satisfies $\Pr(n-2+\frac2kZ_n\le R/k)\le e^{-9M/80}$ under \eqref{eq:cut10}. From $n\ge4(R/k+2)$, $R/k-n+2\le-\frac34n$, so we need $Z_n\le-\frac{3k}{8}n$. The sum of the two Poisson parameters of $Z_n$ satisfies $2\lambda_n\le k^2n^2/(2M)$. By the moment generating function of the Skellam distribution $\E e^{-tZ_n}=\exp\{2\lambda_n(\cosh t-1)\}$ and the exponential Markov inequality (Chernoff bound), for $t>0$,
\begin{equation*}
\Pr\Bigl(Z_n\le-\frac{3kn}8\Bigr)\le\exp\Bigl\{-\frac{3kn}8t+2\lambda_n(\cosh t-1)\Bigr\}.
\end{equation*}
Take $t=M/(2kn)$; since $n\ge M/(2k)$, $t\le1$. Using $\cosh t-1\le\frac35t^2$ (for $0\le t\le1$, by comparing the power series term by term with $t^2$ and using $\cosh1-1<3/5$), the exponent is at most $-\frac{3M}{16}+\frac{k^2n^2}{2M}\cdot\frac35\cdot\frac{M^2}{4k^2n^2}=-\frac{3M}{16}+\frac{3M}{40}=-\frac{9M}{80}$. The condition $n\ge4$ ensures that the expression for $\lambda_n$ is valid.
\end{proof}

\subsection{Smaller cutoff from the Bernstein bound}
The coefficient $4$ in Proposition~\ref{prop:cut10} makes the scan range about $3U$ candidates. The Bernstein bound for the Skellam distribution gives a smaller cutoff.

\begin{proposition}[Bernstein cutoff]\label{prop:cut11}
If $M>4\log(1/u)$, then every candidate satisfying
\begin{equation}
n\ge\frac{\widetilde d+2+\frac{4}{3k}\log(1/u)}{1-2\sqrt{\log(1/u)/M}}\label{eq:cut11}
\end{equation}
is rejected.
\end{proposition}

At $M=16384$ and $u=0.01$, the coefficient of the observation $\widetilde d$ in \eqref{eq:cut11} is about $1.035$, reducing the scanned candidates from about $3U$ to one or two hundred; the number of candidates requiring an exact CDF is unchanged. The search checks downward one by one, starting from the smaller of the two cutoffs minus one.

\begin{proof}
For $0<t<3$, $\log\E e^{-tZ_n}=2\lambda_n(\cosh t-1)=2\lambda_n\sum_{j\ge1}t^{2j}/(2j)!\le\lambda_nt^2\sum_{j\ge0}(t/3)^{j}=\lambda_nt^2/(1-t/3)$, because $2/(2j)!\le3^{-(2j-2)}$. This shows that $-Z_n$ is a sub-gamma variable with variance factor $2\lambda_n$ and scale $1/3$; the standard Bernstein argument \cite{boucheron2013concentration} gives
\begin{equation*}
\Pr(Z_n\le-x)\le\exp\Bigl[-\frac{x^2}{4\lambda_n+2x/3}\Bigr],\qquad x>0.
\end{equation*}
A sufficient condition for the right side to be at most $u$ is $x\ge2\sqrt{\lambda_n\log(1/u)}+\frac23\log(1/u)$. Taking $x=\frac k2(n-R/k-2)$ and using $\lambda_n\le k^2n^2/(4M)$, the sufficient condition becomes $\frac k2(n-R/k-2)\ge kn\sqrt{\log(1/u)/M}+\frac23\log(1/u)$, which rearranges to \eqref{eq:cut11}; $M>4\log(1/u)$ ensures the denominator is positive. When $\lambda_n=0$, the lower tail is zero and the conclusion holds directly.
\end{proof}

Numerically, the cutoff is computed at risk $u(1-10^{-9})$, below the $u-10^{-12}$ used by the acceptance test in the implementation, and one is added after rounding up. Hence rounding errors cannot place the cutoff above a candidate that should be accepted.

\subsection{Per-candidate Chernoff prescreening}
For each candidate, if $j=\lfloor(R-kn+2k)/2\rfloor<0$, first compute the Chernoff bound of the Skellam lower tail, namely the exponential Markov bound of the moment generating function above at the optimal parameter $t=\operatorname{asinh}r$, where $r=-j/(2\lambda_n)$:
\begin{equation*}
\Pr(Z_n\le j)\le\exp\Bigl\{-\Bigl[(-j)\operatorname{asinh}r-2\lambda_n\bigl(\sqrt{1+r^2}-1\bigr)\Bigr]\Bigr\}.
\end{equation*}
Of the two components, the one with shift $-2$ has the larger CDF; if its Chernoff bound is already at most $u(1-10^{-9})$, the average is at most $u$, and the candidate is rejected directly. Prescreening rejects only candidates already proven to be rejected, and turning prescreening off and computing exact CDFs for all candidates gives identical results one by one.

\subsection{Windowed Poisson table for the exact CDF}
$\Pr(Z_n\le j)=\sum_m\Pr(P_2=m)\Pr(P_1\le j+m)$, with $P_1,P_2$ independent $\Pois(\lambda_n)$. The implementation builds one Poisson table per candidate, containing only the window $[\lambda_n-w,\lambda_n+w]$ with $w=14\sqrt{\lambda_n}+40$: starting from the mode, the pmf is recurred outward on both sides by $p(i+1)=p(i)\lambda_n/(i+1)$, normalized within the window, and accumulated to obtain the CDF. By the Chernoff bound for the Poisson tail, the Poisson mass outside the window on each side is below $10^{-40}$, smaller than double-precision rounding; normalization also avoids the relative error of applying $\log\Gamma$ term by term at large $\lambda_n$. The initial implementation built the table from zero up to $\lambda_n+14\sqrt{\lambda_n}$; at $M_1=16384$ and load $3$, $\lambda_n$ is about $2\times10^5$ to $4\times10^5$, a single sizing taking about $1.3$ s and exhausting memory in the statistical experiments; after windowing a single sizing takes about $36$ ms, with identical results one by one.

The Poisson table depends only on $(M,k)$ and $n$, and can be shared across syncs and risk levels. The implementation maintains a cache keyed by $(M,k)$; when full, it evicts the least recently used tables until half the capacity is free; the cache capacity is measured in entries, each entry being two float64.

The cutoffs, prescreening, and windowed table above do not change the output of the rule. We compared against a reference implementation with 70 decimal digits case by case (4536 general cases and 286 bounded-$c$ cases, covering lattice points, parity boundaries, and $+\infty$), and compared the windowed implementation with cutoffs against the full-table implementation on 9779 cases, with identical results one by one.

\section{Experimental Setup and Supplementary Results}\label{app:exp}
We first describe the validation goals and experimental procedures, then present the setup and supplementary results in the order of Table~\ref{tab:exp-grid}.

\subsection{Validation Goals and Experiment Overview}\label{app:exp-overview}
The experiments in Section~\ref{sec:exp} validate the following conclusions. The statistical experiments (\exptag{STAT}) measure the lower-tail shape of the post-failure statistic (Theorems~\ref{thm:cross} and~\ref{thm:clt}) and the conditional miss rate and extra capacity of the two rules at finite size (Theorem~\ref{thm:coverage} and Corollary~\ref{cor:gauss}). For some configurations, we sample from the planted distribution in place of the failure-conditioned distribution. The planted sampling control (\exptag{PLANT}) tests the finite-size error of this substitution (Lemma~\ref{lem:plant}). The decoding-equivalence test (\exptag{DEC}) verifies on the full IBLT that random-sign encoding preserves decoding (Proposition~\ref{prop:decode}). The conditional-unbiasedness test (\exptag{UNB}) verifies that the post-failure estimate is conditionally unbiased given the mapping for every sign composition (Propositions~\ref{prop:moments} and~\ref{prop:composition}). The worst-configuration test (\exptag{WORST}) uses independent samples to test the conditional miss rate of the Skellam rule on the configurations closest to the risk upper bound, supplementing the statistical experiments. The protocol experiments (\exptag{PROTO}) measure the reliability and communication cost of the two-round protocol (Theorem~\ref{thm:protocol}). The computation-cost measurement (\exptag{COST}) measures the time and memory required for sizing and decoding.

\begin{table}[!htbp]
\centering
\caption{Experimental setup. Tags identify experiments in the main text and appendix. The last column lists the subsections reporting the results.}
\label{tab:exp-grid}
\small
\setlength{\tabcolsep}{3pt}
\begin{tabular}{@{}>{\raggedright\arraybackslash}p{.11\linewidth}>{\raggedright\arraybackslash}p{.17\linewidth}>{\raggedright\arraybackslash}p{.30\linewidth}>{\raggedright\arraybackslash}p{.24\linewidth}>{\raggedright\arraybackslash}p{.11\linewidth}@{}}
\toprule
Experiment & Goal & Parameter range & Samples and stopping rule & Location\\
\midrule
\exptag{STAT} statistical experiments & Lower-tail shape (Theorems~\ref{thm:cross}, \ref{thm:clt}); conditional miss rate and extra capacity (Theorem~\ref{thm:coverage}) & $k=3,4,5$, $M\in\{256,1024,4096,16384\}$, $c\in\{0.5,\dots,64\}$, load $\in\{0.95,\dots,3\}$ and $d=2,3,6$, load at most 3; $M=64$ boundary under the same rule & 189 + 42 configurations, about 20{,}000 failures each; fixed number of trials & \ref{sec:exp-setup}, \ref{sec:exp-dist}, \ref{sec:exp-cover}; \ref{app:sampling-impl}, \ref{app:limit}, \ref{app:cover-more}\\
\exptag{PLANT} planted sampling control & Finite-size error of the planted distribution (Lemma~\ref{lem:plant}) & Critical-load group: $k=3,4,5$, $M=64$--1024, load 0.85--0.95, direct failure probability at least $10^{-3}$; low-load group: for $k=3$, $M=1024$--16384, load 0.04--0.61; for $k=4,5$, $M=256,1024$, load 0.04--0.71 & Critical-load group: 21 configurations, 65{,}000 failures each for direct and planted; low-load group: 12 configurations, direct trial count fixed by $\lambda$, about 65{,}000 failures for $k=3$ and about 20{,}000 to 65{,}000 for $k=4,5$, 65{,}000 planted each & \ref{sec:exp-setup}; \ref{app:planted}\\
\exptag{DEC} decoding-equivalence test & Decoding equivalence (Proposition~\ref{prop:decode}) & Grid: $k\in\{3,4,5,6,10\}$, $M\in\{64,1024,16384\}$, load 0.5--3, five sign compositions; near threshold: $M\in\{1024,16384,65536\}$, load 0.90--1.05; rejected candidates: $M=16384$, load 0.5--1, nine sign compositions & Grid: 375 configurations, 500,000,000 paired trials in total; near threshold: 13,000,000; rejected candidates: 6,750,000 & \ref{sec:exp-decode}; \ref{app:decode-exp}\\
\exptag{UNB} conditional-unbiasedness test & Conditional unbiasedness (Propositions~\ref{prop:moments}, \ref{prop:composition}) & $k=3$--6, $M=64$--262144, critical window and its left side with load 0.4--0.95; Figures~\ref{fig:encoding}, \ref{fig:composition}: $M=16384$, failure probability about 0.1 & 203 configurations; Figure~\ref{fig:encoding}: 4,000,000 mappings, Figure~\ref{fig:composition}: 4 configurations & \ref{sec:exp-decode}; \ref{app:dist}\\
\exptag{WORST} worst-configuration test & Risk guarantee on the worst configurations (Theorem~\ref{thm:coverage}) & 12 configurations (Table~\ref{tab:confirm}); another 10 configurations for Figure~\ref{fig:scale} under the same design & 56{,}000--330{,}000 failures per configuration; stop when the block-ordered cumulative failure count reaches the target & \ref{sec:exp-cover}, \ref{sec:exp-e2e}; \ref{app:cover-more}\\
\exptag{PROTO} protocol experiments & Protocol reliability and communication cost (Theorem~\ref{thm:protocol}) & $k=3,4,5$, $M_1=256$--16384, load $\in\{0.95,1,1.05,1.5,3\}$, one-sided and balanced, $u\in\{0.01,0.05\}$ & 120 configurations, 13,940,000 trials in total per level of $u$ & \ref{sec:exp-e2e}; \ref{app:protocol}\\
\exptag{COST} computation-cost measurement & Time and memory of sizing and decoding & $M=16384$, $k=3,4,5$, $c$ from 4 to 314; protocol load 1.0--3 & 10 repetitions per item; candidate-check counts run separately, 30 trials per configuration & \ref{sec:exp-cost}; \ref{app:timing}\\
\bottomrule
\end{tabular}
\end{table}

Before launch, we fix the configurations, seeds, sampling sources, stopping rules, and budgets for the statistical experiments, worst-configuration test, planted sampling control, and protocol experiments using a single feasibility pilot. We fix the selection rule for the worst-configuration test before the statistical experiments begin, then derive the specific configurations and sample sizes from that rule after the statistical experiments finish (Appendix~\ref{app:cover-more}). We conduct the conditional-unbiasedness test and decoding-equivalence test in separate independent runs, and measure computation cost separately from the other runs (Appendix~\ref{app:timing}).

\subsection{Sampling Method}\label{app:sampling-impl}
For the statistical experiments, we choose the sampling source based on the cost estimated in the pilot. We use direct sampling when the estimated cost of obtaining 20{,}000 failures is at most 20 CPU core-hours, setting the number of trials to $\lceil2\times10^4/\widehat{\Pr}(F)\rceil$. For the remaining configurations, we use 20{,}000 planted samples. Figure~\ref{fig:exp-grid} shows each configuration's location in $c$ and load and its sampling source. We run the six points at load 3 ($d=\mathrm{round}(3\alpha_cM)$, with rounding making the load slightly above 3) separately with independent seeds after the statistical experiments finish. These points use the same sampling method and target failure count as the other configurations.

\begin{figure}[!htbp]
\centering\includegraphics[width=\linewidth]{../figure/fig-exp-grid.pdf}
\caption{The 189 configurations of the statistical experiments. Point shapes indicate the sampling source. Vertical lines mark the grouping boundaries of $c$, and horizontal shading marks load 0.95--1.05.}
\label{fig:exp-grid}
\end{figure}

Planted sampling selects a pair of difference keys uniformly at random and assigns them the same uniformly random $k$-cell set. We generate the remaining keys' mappings independently and uniformly, and draw the signs independently. This construction always produces a peeling failure. Direct sampling requires an average of $1/\Pr(F)$ mappings per failure, whereas planted sampling generates a mapping drawn from $Q$ directly. Lemma~\ref{lem:plant} gives its asymptotic relation to the true failure-conditioned distribution, and Appendix~\ref{app:planted} gives the finite-size comparison. Under iid random signs, the statistic is label-invariant, so we equivalently fix the last two labels as the duplicate edge pair in these experiments.

For the conditional-unbiasedness test, we combine direct sampling and splitting on an NVIDIA H100 with CPU sampling. The GPU direct sampler processes one graph per thread block and performs peeling in synchronized rounds. Each round first identifies degree-one cells, then uses atomic operations to claim and delete the corresponding edges, ensuring that each edge is deleted once. We use Philox to generate random numbers, organizing them by seed and trial index so that the random stream is independent of the GPU block partition. An independent CPU implementation re-peels the same exported batch of mappings and compares the final residual core graph by graph. For the initial direct sampler, we find no differences on 14,100 graphs, and the failure-probability statistical controls on all 15 configurations pass the preset criterion. Synchronized peeling and sequential CPU peeling yield the same final core. The CPU implementation handles path statistics that depend on the peeling order.

\subsection{Finite-Size Error of the Planted Distribution}\label{app:planted}
The planted sampling control (\exptag{PLANT}) compares finite-size failure samples from the planted distribution and the failure-conditioned distribution (Lemma~\ref{lem:plant}). We compare three quantities: the fraction of pure duplicate edge pairs among direct failures, the full-range KS distance between the two samples' post-failure $X$ distributions, and the miss/u of the Skellam rule on each sample. We compare the KS distance with the asymptotic 95\% critical value: 0.0075 with 65{,}000 failures on each side, and 0.011 with 20{,}000 on the direct side. This comparison is conservative because $X$ is discrete. We divide the control into two groups. The low-load group contains 12 configurations from the planted configurations of the statistical experiments and directly tests the planted samples used in the main text. The critical-load group contains 21 configurations at load 0.85 to 0.95, where direct failures are common, and tests how the error of the planted approximation varies with $M$. Figure~\ref{fig:planted} summarizes both groups.

Most of the 105 planted configurations in the statistical experiments lie far from the critical load. The expected number of duplicate edge pairs $\lambda=\binom d2/\binom Mk$ is therefore small, and Lemmas~\ref{lem:plant} and~\ref{lem:propL} imply that the failure probability is of the same order as $\lambda$. The low-load group consists of the 12 configurations for which direct sampling is affordable. Before sampling, we fix the direct-sampling trial count using $\lambda$ (Table~\ref{tab:v-low}). The four configurations with $k=3$ cover $c=1$ to 64 and load 0.04 to 0.61, with about 65{,}000 failures on each side. At $u=0.05$, all equivalence tests with bound $\pm0.1$ and one-sided level $0.05/4$ pass. At $u=0.01$, all difference intervals contain zero. The KS distances range from 0.0048 to 0.0079, with three below the critical value. The distance for $(3,4096,1024)$ is 0.0079 ($p=0.034$), so no difference is detectable after correction for 12 tests. The eight configurations with $k=4,5$ cover $c=0.5$ to 16 and load 0.04 to 0.71, where failures are rarer. Each has about 20{,}000 direct failures, except $(k,M,d)=(4,256,16)$ and $(5,256,128)$, which have about 65{,}000. Each has 65{,}000 planted samples. At $u=0.05$, all point estimates of the differences lie within $\pm0.04$, and all 95\% intervals contain zero. The equivalence tests with bound $\pm0.1$ and one-sided level $0.05/8$ pass for 5 configurations. For the other 3 configurations ($(4,1024,512)$, $(5,256,32)$, $(5,256,64)$), the interval endpoints lie within $\pm0.13$. At $u=0.01$, the difference intervals for 7 configurations contain zero. The difference for $(4,1024,512)$ is $-0.16$ (interval $[-0.30,-0.02]$), indicating higher miss in the planted sample. The KS distances range from 0.0048 to 0.0071, all below their respective critical values. Across the 12 configurations, the fraction of pure duplicate edge pairs among direct failures ranges from 0.988 to 1.000, consistent with the mechanism in Lemma~\ref{lem:plant}.

\begin{table}[!htbp]
\centering
\caption{Direct and planted sampling in the low-load group (Skellam rule, $u=0.05$). miss/u lists direct / planted. Difference gives the 95\% interval for direct minus planted. KS is the full-range KS distance of $X$, with 95\% critical value 0.0075 for 65{,}000 direct failures and 0.011 for 20{,}000.}
\label{tab:v-low}
\small
\begin{tabular}{@{}rrrrrrcccc@{}}
\toprule
& & & & & direct & pure duplicate edge pairs & & & \\
$k$ & $M$ & $d$ & $c$ & load & failure count & fraction & miss/u & difference & KS\\
\midrule
3 & 1024 & 32 & 1 & 0.04 & 65061 & 1.000 & 0.94 / 0.94 & $[-0.05,+0.04]$ & 0.0059\\
3 & 4096 & 512 & 8 & 0.15 & 64977 & 1.000 & 0.95 / 0.96 & $[-0.06,+0.03]$ & 0.0048\\
3 & 4096 & 1024 & 16 & 0.31 & 65136 & 0.999 & 0.92 / 0.94 & $[-0.07,+0.03]$ & 0.0079\\
3 & 16384 & 8192 & 64 & 0.61 & 65530 & 0.999 & 0.98 / 0.99 & $[-0.06,+0.04]$ & 0.0051\\
\addlinespace
4 & 256 & 8 & 0.5 & 0.04 & 20071 & 0.998 & 0.54 / 0.52 & $[-0.03,+0.07]$ & 0.0071\\
4 & 256 & 16 & 1 & 0.08 & 65839 & 0.995 & 0.68 / 0.69 & $[-0.05,+0.03]$ & 0.0057\\
4 & 1024 & 128 & 4 & 0.16 & 20067 & 0.997 & 0.95 / 0.95 & $[-0.06,+0.07]$ & 0.0049\\
4 & 1024 & 256 & 8 & 0.32 & 20265 & 0.994 & 0.98 / 0.97 & $[-0.06,+0.08]$ & 0.0048\\
4 & 1024 & 512 & 16 & 0.65 & 20038 & 0.988 & 0.94 / 0.97 & $[-0.10,+0.04]$ & 0.0062\\
5 & 256 & 32 & 2 & 0.18 & 20137 & 0.999 & 1.00 / 0.98 & $[-0.05,+0.09]$ & 0.0068\\
5 & 256 & 64 & 4 & 0.36 & 19885 & 0.997 & 0.91 / 0.95 & $[-0.11,+0.03]$ & 0.0048\\
5 & 256 & 128 & 8 & 0.71 & 65470 & 0.995 & 0.99 / 1.01 & $[-0.07,+0.02]$ & 0.0063\\
\bottomrule
\end{tabular}
\end{table}

Direct sampling is infeasible for the remaining planted configurations. Among the 77 planted configurations with $k=4,5$, 76 have $\lambda$ below $10^{-6}$, with a minimum of about $10^{-19}$. The feasibility pilot yields no failures in 68 of 75 configurations. Collecting 65{,}000 failures requires at least $6.5\times10^{10}$ trials, and about half the configurations require more than $10^{15}$, exceeding the feasible range of direct sampling on CPU and GPU. For these configurations, Lemma~\ref{lem:plant} controls the distance between the planted distribution and the failure-conditioned distribution, with an error term that tends to zero when $d=o(M)$. The approximation is most accurate precisely where direct sampling is hardest.

The critical-load group tests the range of applicability of the planted approximation. The statistical experiments include five planted configurations at load 0.95, all with $M\ge4096$. This group uses smaller $M$ to compare the two samples where larger stopping sets cause most failures. We determine the sample size using an equivalence test on the miss/u difference at $u=0.05$, with bound $\pm0.1$, one-sided level $0.05/21$, and 80\% power when the difference is zero. This gives 65{,}000 failures on each side. Table~\ref{tab:v} reports the results for these 21 configurations.

\begin{table}[!htbp]
\centering
\caption{Direct and planted sampling in the critical-load group (Skellam rule). miss/u lists direct / planted. KS is the full-range KS distance of $X$, with 95\% critical value 0.0075.}
\label{tab:v}
\small
\begin{tabular}{@{}rrrrcccc r@{}}
\toprule
$k$ & $M$ & load & pure duplicate edge pairs & \multicolumn{2}{c}{$u=0.01$} & \multicolumn{2}{c}{$u=0.05$} & KS\\
\cmidrule(lr){5-6}\cmidrule(lr){7-8}
& & & fraction & miss/u & difference & miss/u & difference & \\
\midrule
3 & 64 & 0.86 & 0.09 & 0.44 / 0.48 & $-0.04\pm0.07$ & 0.87 / 0.88 & $-0.01\pm0.04$ & 0.0026\\
3 & 64 & 0.90 & 0.05 & 0.49 / 0.52 & $-0.03\pm0.08$ & 0.82 / 0.87 & $-0.05\pm0.04$ & 0.0085\\
3 & 64 & 0.95 & 0.02 & 0.44 / 0.51 & $-0.07\pm0.07$ & 0.87 / 0.94 & $-0.07\pm0.05$ & 0.0088\\
4 & 64 & 0.85 & 0.01 & 0.56 / 0.49 & $+0.08\pm0.08$ & 0.89 / 0.92 & $-0.03\pm0.05$ & 0.0105\\
4 & 64 & 0.89 & 0.01 & 0.54 / 0.65 & $-0.11\pm0.08$ & 0.89 / 0.96 & $-0.07\pm0.05$ & 0.0097\\
4 & 64 & 0.95 & 0.00 & 0.39 / 0.55 & $-0.16\pm0.07$ & 0.81 / 0.93 & $-0.13\pm0.04$ & 0.0131\\
5 & 64 & 0.85 & 0.00 & 0.58 / 0.69 & $-0.11\pm0.09$ & 0.96 / 1.07 & $-0.11\pm0.05$ & 0.0130\\
5 & 64 & 0.89 & 0.00 & 0.57 / 0.69 & $-0.12\pm0.09$ & 1.00 / 1.13 & $-0.13\pm0.05$ & 0.0142\\
5 & 64 & 0.96 & 0.00 & 0.55 / 0.72 & $-0.16\pm0.09$ & 0.89 / 1.04 & $-0.15\pm0.05$ & 0.0118\\
3 & 256 & 0.85 & 0.63 & 0.75 / 0.62 & $+0.13\pm0.09$ & 0.89 / 0.88 & $+0.01\pm0.04$ & 0.0063\\
3 & 256 & 0.90 & 0.13 & 0.75 / 0.66 & $+0.09\pm0.09$ & 0.95 / 0.88 & $+0.07\pm0.05$ & 0.0080\\
3 & 256 & 0.95 & 0.02 & 0.77 / 0.64 & $+0.14\pm0.09$ & 0.92 / 0.91 & $+0.01\pm0.05$ & 0.0060\\
4 & 256 & 0.90 & 0.01 & 0.71 / 0.70 & $+0.01\pm0.09$ & 0.93 / 0.98 & $-0.05\pm0.05$ & 0.0045\\
4 & 256 & 0.95 & 0.00 & 0.72 / 0.73 & $-0.01\pm0.09$ & 0.93 / 0.97 & $-0.05\pm0.05$ & 0.0052\\
5 & 256 & 0.90 & 0.00 & 0.71 / 0.75 & $-0.04\pm0.09$ & 0.92 / 0.94 & $-0.02\pm0.05$ & 0.0052\\
5 & 256 & 0.95 & 0.00 & 0.73 / 0.80 & $-0.07\pm0.09$ & 0.95 / 0.94 & $+0.00\pm0.05$ & 0.0027\\
3 & 1024 & 0.85 & 0.97 & 0.79 / 0.78 & $+0.01\pm0.10$ & 0.92 / 0.91 & $+0.02\pm0.05$ & 0.0033\\
3 & 1024 & 0.90 & 0.96 & 0.83 / 0.86 & $-0.03\pm0.10$ & 0.93 / 0.95 & $-0.02\pm0.05$ & 0.0037\\
3 & 1024 & 0.95 & 0.08 & 0.84 / 0.84 & $+0.01\pm0.10$ & 0.96 / 0.96 & $-0.00\pm0.05$ & 0.0091\\
4 & 1024 & 0.95 & 0.00 & 0.89 / 0.91 & $-0.02\pm0.10$ & 0.97 / 0.94 & $+0.03\pm0.05$ & 0.0054\\
5 & 1024 & 0.95 & 0.00 & 0.85 / 0.79 & $+0.06\pm0.10$ & 0.96 / 0.95 & $+0.01\pm0.05$ & 0.0039\\
\bottomrule
\end{tabular}
\end{table}

For all five tested $M=1024$ configurations, the equivalence tests with bound $\pm0.1$ pass at $u=0.05$, and the difference intervals at both levels of $u$ contain zero. These include configurations where the fraction of pure duplicate edge pairs among direct failures is 0.97 and others where it is close to zero. For seven of the nine $M=64$ configurations, the tests do not establish that the difference lies within the bound. The planted sample has higher miss in all seven, with the difference interval extending as far as $-0.20$. The tests also do not establish the bound for three of the seven $M=256$ configurations: at $k=3$ and load 0.90 the direct sample has higher miss, with difference $+0.07$; at $k=4$ and loads 0.90 and 0.95 the difference is $-0.05$ in both. For the two configurations with $k=3$, $M=256$, and loads 0.85 and 0.95, both difference intervals contain zero at $u=0.05$. At $u=0.01$, the direct sample's miss/u is higher by 0.13 (interval $[+0.04,+0.22]$) and 0.14 (interval $[+0.04,+0.23]$), respectively.

The KS distance decreases as $M$ grows. For the five $M=1024$ configurations, it ranges from 0.0033 to 0.0091, with four below the critical value. The configuration $(k,M,d)=(3,1024,796)$ exceeds the critical value because of differences in the middle of the distribution. The maximum difference in the lower tail between the 0.2\% and 15\% quantiles of the pooled sample is only 0.0039. For the seven $M=256$ configurations, the distance ranges from 0.0027 to 0.0080, with six below the critical value. For the nine $M=64$ configurations, it ranges from 0.0026 to 0.0142, with eight above the critical value. At $M=1024$, the configuration with a pure duplicate edge pair fraction of 0.97 and those with fractions close to zero have similarly small KS distances. Thus, the distributions of $X$ can be close even when the failure mechanisms differ. The fraction of pure duplicate edge pairs alone is insufficient to predict the distribution error of $X$, which also depends on overlap fluctuations on the $s(d)$ scale.

\begin{figure}[!htbp]
\centering\includegraphics[width=\linewidth]{../figure/fig-planted.pdf}
\caption{Direct versus planted sampling. The upper three panels show the critical-load group, grouped by $M$; the bottom panel shows the low-load group. The vertical axis labels $k$ and load. Left: Skellam rule miss/u difference (direct minus planted) with 95\% intervals and shading at $\pm0.1$. Right: full-range KS distance of $X$; gray ticks mark the 95\% critical value for each sample size.}
\label{fig:planted}
\end{figure}

\subsection{Random-Sign Encoding Does Not Change Decoding}\label{app:decode-exp}
The decoding-equivalence test (\exptag{DEC}) uses the two-round protocol's actual IBLT encoder and decoder. We build and subtract the two parties' tables, then compare first-round decoding under standard encoding and random-sign encoding. The test consists of three groups of paired trials, all using the same construction and verification procedure.

In each trial, we generate distinct keys and fresh seeds for the cell hash, sign hash, and checksum hash. Both encodings share the same keys and hashes. Before decoding, we verify cell by cell that the three fields after cancellation of the common elements match those obtained by encoding the difference directly. An independent reference uses only the unsigned incidence graph of the true difference and computes its residual 2-core by the degree-one rule. Whenever the actual decoder accepts a key, we check the true residual degree, key, and direction, then update the reference state. These checks cover trajectories ending in both success and failure. After termination, we verify that the recovered keys are exactly the difference outside the 2-core, with all directions correct. We reconstruct the three fields from the true residual, compare them cell by cell, and check that all fields are zero if and only if the residual core is empty.

The first group is a grid spanning the full load range: $k\in\{3,4,5,6,10\}$, $M\in\{64,1024,16384\}$, load $d/(\alpha_cM)\in\{0.5,0.95,1,1.05,3\}$, and positive-direction ratio $d_+/d$ with target values $0,1/4,1/2,3/4,1$. We round $d$ and $d_+$ to integers from the target values. We use the corresponding peeling threshold for each $k$, giving 375 configurations in total. The number of common elements is $\min\{1000,M\}$. For $M=64,1024,16384$, we run 3,000,000, 900,000, and 100,000 trials per configuration, respectively. This gives 500,000,000 trials in total, with 100,000,000 for each $k$.

Table~\ref{tab:decode-paired} reports the results by $k$. Across all configurations, we observe zero paired success/failure disagreements, non-pure acceptances, incorrect recoveries, or residual-core disagreements. The two decoders together check 211,516,840,988 acceptance operations, all satisfying the true degree-one and direction conditions.

The second group samples loads more densely near the threshold and provides the results in Table~\ref{tab:decode} of the main text. The configurations use $k\in\{3,4,5,6,10\}$, $M\in\{1024,16384,65536\}$, 10 load values between 0.90 and 1.05, and $d_-/d\in\{0,1/2\}$. For $M=1024,16384,65536$, we run 100,000, 20,000, and 10,000 trials per configuration, respectively, for a total of 13,000,000. All checks pass. At $M=16384$, the failure probability ranges from 0.020 to 0.060 at load 0.99 and from 0.993 to 0.998 at load 1.01. At $M=65536$, almost no trials fail at load 0.99, and almost all fail at load 1.01. After failure, the decoder still recovers the difference outside the 2-core. At $M=65536$ and load 1.02, the average recovery fraction ranges from 0.18 ($k=10$) to 0.50 ($k=3$).

The third group of paired trials counts rejected candidates separately for each encoding. These are cells with count $\pm1$ that fail the checksum or cell-ownership check. The configurations use $M=16384$, $k\in\{3,4,5,6,10\}$, load $0.5,0.95,1$, and nine values of $d_-/d$ from 0 to 1 in steps of $1/8$. We run 50,000 trials per configuration, for a total of 6,750,000, and all checks pass. Figure~\ref{fig:decode} reports the results for $k=3$ to 6. Standard encoding has no rejected candidates under one-sided differences because the count equals the degree. As the difference approaches balanced, more opposite-sign cancellation occurs and more candidates are rejected. Under random-sign encoding, the number of rejected candidates differs by at most 0.3\% across the nine sign compositions and matches standard encoding in the balanced case. This agrees with Proposition~\ref{prop:composition}: the coding coefficient is an iid random sign independent of $(d_+,d_-)$. The results for $k=10$ follow the same pattern, with both encodings yielding about 1.27 rejected candidates per difference key at load 1. Across all three groups, the two decoders together reject $2.8\times10^{11}$ non-pure candidates with count $\pm1$, with zero false acceptances. The cell-ownership check rejects zero candidates, and the checksum rejects every incorrect candidate.

\begin{figure}[!htbp]
\centering\includegraphics[width=\linewidth]{../figure/fig-decode.pdf}
\caption{Rejected candidates per first-round decoding, divided by $d$ ($M=16384$; both encodings share the same mappings for each configuration). Solid lines with filled markers show standard encoding; dashed lines with open markers show random-sign encoding.}
\label{fig:decode}
\end{figure}

\begin{table}[!htbp]
\centering
\caption{First-round paired decoding results.}
\label{tab:decode-paired}
\begin{tabular}{rrrr}
\toprule
$k$ & paired trials & successes & failures\\
\midrule
3 & 100,000,000 & 36,532,515 & 63,467,485\\
4 & 100,000,000 & 35,900,798 & 64,099,202\\
5 & 100,000,000 & 33,530,250 & 66,469,750\\
6 & 100,000,000 & 32,553,378 & 67,446,622\\
10 & 100,000,000 & 30,233,256 & 69,766,744\\
\bottomrule
\end{tabular}
\end{table}

\subsection{The Post-Failure Estimate Is Conditionally Unbiased for Every Sign Composition}\label{app:dist}
Of the 203 configurations in the conditional-unbiasedness test (\exptag{UNB}), 196 have at least 1000 failure samples. Under random-sign encoding, $|\E[X\mid F]|$ is at most 0.032 for these configurations, with 93\% within two standard errors of zero, consistent with the exact unbiasedness in Proposition~\ref{prop:moments}. The conditional variance to the left of the critical window is 0.972--1.049. Among the 84 direct-sampling configurations in the statistical experiments, $|\E[X\mid F]|$ under random-sign encoding is at most 0.019, with 98\% within two standard errors of zero. In the statistical experiments at load 1.0 and failure probability about 0.67, standard encoding with one-sided differences has conditional means of $+0.20$, $+0.02$, and $-0.08$ for $k=3,4,5$, respectively (standard error 0.007).

Table~\ref{tab:mean-by-k} reports the values in Figure~\ref{fig:composition} of the main text. For each $(k,M,d)$, we record five sign compositions under both encodings. The ten datasets share the same mappings, and both encodings use all failure records. The four configurations have failure probabilities of 0.099--0.102. All 20 values under random-sign encoding have absolute value at most 0.005 and lie within two standard errors of zero. We generate the data for Figure~\ref{fig:encoding} using the sampling code from the interactive experiment page Peeling Lens. The code computes $X$ for three encodings on each random mapping, drawing random signs from an independent generator. This configuration has $c\approx104$. Under standard encoding, $\E[X\mid F]=+0.599\pm0.001$ for one-sided differences and $-0.002$ for balanced differences. Under random-sign encoding, it is $+0.001\pm0.001$, with conditional variance 1.007. The failure samples and the full sample have similar distributions, consistent with the large-$c$ behavior in Theorem~\ref{thm:clt}.

\begin{table}[!htbp]
\centering
\caption{Values in Figure~\ref{fig:composition}: mean post-failure $X$ ($M=16384$, failure probability about 0.1). Standard errors are 0.0018 for $k=3$ and 0.0031--0.0032 for $k=4,5,6$.}
\label{tab:mean-by-k}
\begin{tabular}{rrlrrrrr}
\toprule
& & & \multicolumn{5}{c}{$d_-/d$}\\
\cmidrule(lr){4-8}
$k$ & $d$ & encoding & 0 (one-sided) & 0.1 & 0.25 & 0.4 & 0.5 (balanced)\\
\midrule
3 & 13306 & standard & $+0.638$ & $+0.405$ & $+0.158$ & $+0.028$ & $+0.003$\\
  &       & random sign & $+0.002$ & $-0.001$ & $0.000$ & $-0.004$ & $+0.002$\\
\addlinespace
4 & 12567 & standard & $+0.084$ & $+0.052$ & $+0.016$ & $-0.002$ & $-0.003$\\
  &       & random sign & $-0.001$ & $-0.002$ & $+0.005$ & $+0.001$ & $-0.001$\\
\addlinespace
5 & 11418 & standard & $-0.218$ & $-0.137$ & $-0.054$ & $-0.009$ & $+0.002$\\
  &       & random sign & $+0.003$ & $-0.005$ & $0.000$ & $-0.002$ & $-0.004$\\
\addlinespace
6 & 10363 & standard & $-0.387$ & $-0.244$ & $-0.096$ & $-0.021$ & $-0.004$\\
  &       & random sign & $+0.002$ & $+0.003$ & $+0.005$ & $+0.001$ & $+0.001$\\
\bottomrule
\end{tabular}
\end{table}

\subsection{Distance between the Post-Failure Lower Tail and the Limit Distribution}\label{app:limit}
We use the failure records from the statistical experiments (\exptag{STAT}). For Figure~\ref{fig:crosscdf}, we compute $\widetilde d=\sum_vC_v^2/k$ from the raw integer sums of squares, retain all empirical values, and standardize them on the same scale as $G_d$ in the main text. We compute the KS distance immediately to the left and right of every value taken by either distribution. By \eqref{eq:raw-decomp} in Appendix~\ref{app:ams}, the difference between $\widehat d$ and $\widetilde d$ is about $0.02s(d)$ at $M=16384$. This perturbation is small relative to the standard deviation, yet it shifts the discrete atoms. We therefore use $\widetilde d$, as in Theorem~\ref{thm:cross}, to validate the distribution.

Table~\ref{tab:limit-ks} reports the full KS distances computed from the raw integer records of the statistical experiments. The column for the discrete limit distribution compares $\widetilde d$ with $G_d$ for $c=1,2,4$ and $k=3,4,5$, giving nine configurations for each $M$. We use direct and planted samples for $M\le1024$ and planted samples for $M\ge4096$. The Gaussian control in the same column uses $(\widetilde d-d)/s(d)$. The Gaussian column on the right compares $X=(\widehat d-d)/s(d)$ with $N(0,1)$ for the 75 direct-sampling configurations with $c\ge8$, including the critical window and overload. Each configuration has about 20{,}000 failures. At this fixed sample size, the per-configuration 95\% DKW error bound is about 0.010. The discrete approximation has visible finite-size error for small tables, which generally decreases as $M$ grows. The Gaussian quantile error also decreases as $M$ grows.

\begin{table}[!htbp]
\centering
\caption{Full post-failure KS distance. The discrete limit uses $\widetilde d$; the Gaussian uses $\widehat d$. Quantile error is in units of $s(d)$ and gives the maximum absolute error across configurations.}
\label{tab:limit-ks}
\small
\begin{tabular}{rcccccc}
\toprule
& \multicolumn{3}{c}{discrete limit ($c=1,2,4$)} & \multicolumn{3}{c}{Gaussian ($c\ge8$, direct sampling)}\\
\cmidrule(lr){2-4}\cmidrule(lr){5-7}
$M$ & KS & KS of Gaussian ($c=1$) & configs & KS & 1\% quantile error & configs\\
\midrule
256 & 0.019--0.034 & 0.237--0.252 & 9 & 0.010--0.021 & 0.18 & 23\\
1024 & 0.008--0.023 & 0.224--0.244 & 9 & 0.005--0.015 & 0.10 & 23\\
4096 & 0.004--0.016 & 0.215--0.230 & 9 & 0.005--0.014 & 0.06 & 17\\
16384 & 0.003--0.012 & 0.214--0.230 & 9 & 0.003--0.010 & 0.05 & 12\\
\bottomrule
\end{tabular}
\end{table}

The conditional variance reflects the same structure. Among the 84 direct-sampling configurations in the statistical experiments, the conditional variance under random-sign encoding is 0.98--1.06 for $c\ge8$ and 0.985--1.020 within the critical window. For small $c$, it is well above one: about 1.12 at $c=4$, about 1.5 at $c=2$, and above 10 when $c<1$. These values agree with $1+2/c^2$, obtained by converting the limit distribution's variance $4+2c^2$ to the $X$ scale. The excess comes from the duplicate edge pair's $\pm2$ shift, which is non-negligible relative to $s(d)$.

\subsection{Conditional Miss Rate and Extra Capacity at Finite Size}\label{app:cover-more}
Of the 189 configurations in the statistical experiments, 84 use direct sampling and 105 use planted sampling. The Skellam rule's maximum miss/u on the direct-sampling configurations is 1.04 at both levels of $u$, occurring within the critical window. Its maximum on the planted configurations is 1.03, computed under the planted distribution. Each configuration has about 20{,}000 failures, corresponding to about 200 misses at $u=0.01$. The standard error of miss/u is about 0.07, so a value slightly above 1, such as 1.04, is statistically indistinguishable from 1. On the low-load configurations with $c<8$, the Skellam rule's median miss/u at $u=0.01$ is at most 0.76. Every sample yields at least one accepted candidate. If no candidate is accepted, $U=+\infty$, which does not count as a miss. For $c\ge8$, the Gaussian rule gives miss/u of 0.83--1.19 across all configurations. At $c=1$, $u=0.01$, it is about 11 for $k=3$ and about 8.5 for $k=4,5$. When $c$ is small and the observation is near $d-2$, observations for $d=2,3$ fall below every candidate's acceptance threshold. The Gaussian rule then returns $+\infty$, giving no miss. On the direct-sampling configurations with $c\ge8$ in the statistical experiments, the Skellam rule's median miss/u at the two levels of $u$ is 0.79--0.86 and 0.95--0.97.

At $u=0.01$ and $c\ge8$, the Skellam rule's median extra capacity is 25\%, 11.3\%, 5.4\%, and 2.6\% for $M=256$, 1024, 4096, 16384, respectively. The corresponding 90\% quantiles are at most 41\%, 18\%, 8.5\%, and 4.1\%. The Gaussian rule's medians are 23.5\%, 10.8\%, 5.3\%, and 2.6\%, differing from the Skellam rule by at most 1.5 percentage points. At $c=1,2,4$, the median of $U-d$ is 5, 8, 15 for the Skellam rule and 3, 6, 14 for the Gaussian rule. For fixed small $d$, the 90\% quantile of the upper bound at $d=2$ exceeds $d$ by 4--5. Among the 42 configurations with $M=64$, the Skellam rule's maximum miss/u is 1.08 at $u=0.05$ ($k=4$, $d=6$, 95\% Wilson interval 1.02--1.15) and 1.04 at $u=0.01$.

Table~\ref{tab:gauss-c} reports the Gaussian rule's miss/u at fixed $c$ for $M=16384$, corresponding to the triangles with $c\le16$ in Figure~\ref{fig:coverage} of the main text. The excess grows rapidly as $c$ decreases. The miss/u is about 8--11 at $c=1$ and above 33 at $c=0.5$ for $u=0.01$. For $c\le4$, the values at $M=4096$ and 16384 differ by at most 8\%, so the excess persists as the table size grows.

\begin{table}[!htbp]
\centering
\caption{Gaussian rule miss/u at fixed $c$ ($M=16384$).}
\label{tab:gauss-c}
\begin{tabular}{rrrrrrr}
\toprule
& \multicolumn{3}{c}{$u=0.01$} & \multicolumn{3}{c}{$u=0.05$}\\
\cmidrule(lr){2-4}\cmidrule(lr){5-7}
$c$ & $k=3$ & $k=4$ & $k=5$ & $k=3$ & $k=4$ & $k=5$\\
\midrule
0.5 & 36.4 & 33.2 & 33.2 & 9.37 & 8.81 & 9.39\\
1 & 11.3 & 8.74 & 8.27 & 4.09 & 4.24 & 4.60\\
2 & 3.06 & 2.48 & 2.60 & 2.00 & 1.70 & 1.94\\
4 & 1.51 & 1.41 & 1.50 & 1.21 & 1.19 & 1.20\\
8 & 1.12 & 1.03 & 1.08 & 1.07 & 1.08 & 1.06\\
16 & 0.98 & 0.93 & 0.99 & 0.98 & 1.01 & 1.00\\
\bottomrule
\end{tabular}
\end{table}


The worst-configuration test (\exptag{WORST}) uses samples independent of the statistical experiments. We fix the selection rule and sample sizes before sampling. We first select the five configurations whose miss/u exceeds 1 in exploration. We fill the remaining slots by ranking configurations in descending order of the larger one-sided 95\% exact upper bound divided by $u$ over the two levels of $u$ in the statistical experiments. Four highly ranked configurations have only planted samples, and direct sampling cannot obtain failures at feasible cost. Following the rule, we skip these in favor of the next configurations. Their miss/u under $Q$ is at most 1.02. For each configuration and each level of $u$, we use the previously estimated 95\% Wilson upper bound (at most 1) as the planning value. We choose the failure count to give a passing probability of at least $1-0.2/24$ for each judgment and take the larger count over the two levels. At the planning value, this gives at least 80\% probability that all 24 judgments pass. A judgment passes if its one-sided simultaneous confidence upper bound is at most $1.1u$. We use the Clopper--Pearson interval with a Bonferroni correction over the 24 judgments. All 24 pass, with a maximum upper bound of $1.024u$. Of the 24 two-sided 95\% intervals, 22 lie entirely below 1 and the other 2 cover 1. Only one test configuration uses $M=256$. We sample in blocks and analyze all blocks through the first point at which the cumulative failure count in block order reaches the target. Stopping depends only on the failure count. For the two test configurations with $k=3$, $M=1024$, the Gaussian rule's miss/u is 11.1 ($u=0.01$) and 4.1 ($u=0.05$) at $c=1$, and 2.9 and 1.8 at $c=2$. Table~\ref{tab:confirm} reports the load, failure count, miss/u, and one-sided simultaneous confidence upper bound for each test configuration. The segments in Figure~\ref{fig:confirm} of the main text are two-sided 95\% intervals.

\begin{table}[!htbp]
\centering
\caption{Skellam rule miss/u on the \exptag{WORST} configurations, with one-sided simultaneous confidence upper bounds in parentheses (per-item level $0.05/24$).}
\label{tab:confirm}
\begin{tabular}{lrrcc}
\toprule
configuration & load & failures & $u=0.01$ & $u=0.05$\\
\midrule
$k=3$, $M=1024$, $d=32$ & 0.04 & 55{,}509 & 0.53 (0.62) & 0.97 (1.02)\\
$k=3$, $M=1024$, $d=64$ & 0.08 & 289{,}001 & 0.90 (0.95) & 0.91 (0.94)\\
$k=3$, $M=1024$, $d=192$ & 0.23 & 289{,}394 & 0.79 (0.83) & 0.95 (0.97)\\
$k=4$, $M=256$, $d=128$ & 0.65 & 105{,}689 & 0.74 (0.82) & 0.93 (0.97)\\
$k=4$, $M=1024$, $d=751$ & 0.95 & 298{,}439 & 0.85 (0.90) & 0.95 (0.98)\\
$k=3$, $M=4096$, $d=3352$ & 1.00 & 333{,}476 & 0.93 (0.98) & 0.97 (1.00)\\
$k=5$, $M=4096$, $d=3018$ & 1.05 & 332{,}350 & 0.90 (0.95) & 0.99 (1.01)\\
$k=5$, $M=1024$, $d=755$ & 1.05 & 332{,}353 & 0.83 (0.87) & 0.96 (0.98)\\
$k=4$, $M=16384$, $d=16384$ & 1.29 & 293{,}250 & 0.95 (1.01) & 0.98 (1.01)\\
$k=3$, $M=4096$, $d=5029$ & 1.50 & 332{,}350 & 0.91 (0.96) & 0.98 (1.00)\\
$k=4$, $M=16384$, $d=18980$ & 1.50 & 290{,}808 & 0.96 (1.02) & 0.98 (1.00)\\
$k=4$, $M=16384$, $d=37959$ & 3.00 & 289{,}952 & 0.95 (1.01) & 0.97 (1.00)\\
\bottomrule
\end{tabular}
\end{table}

Of the 12 configurations in the upper row of Figure~\ref{fig:scale}, two ($k=3$, $M=4096$ and $k=4$, $M=16384$) are test configurations. The other 10 follow the same design, with sample size and stopping rule fixed before sampling, and each has about 290{,}000 direct-sampling failures. Under random-sign encoding, the common elements cancel in the difference table and sign composition does not affect the post-failure distribution. The $|b|\le1$ used by the protocol under balanced differences also leaves the upper bound unchanged. The protocol experiments, statistical experiments, and these direct samples therefore estimate the same conditional miss rate, and we pool them by configuration in the figure. For these 10 configurations, the Skellam rule's one-sided upper bound is at most $1.025u$ at $u=0.01$ and $1.003u$ at $u=0.05$.

\subsection{Reliability and Communication Cost of the Two-Round Protocol}\label{app:protocol}
In each trial, we generate the common elements and difference elements separately. For one-sided differences, all differences lie on the $A$ side; for balanced differences, half lie on each side. Common elements cancel in the difference table and participate as usual in table construction and cell-by-cell subtraction. Trials that succeed in the first round require no upper bound, so only first-round-failure trials enter the comparison. We evaluate the two levels of $u$ separately on the same batch of trials. We impose a minimum of 64 cells on the second-round table length specified in Section~\ref{sec:flow}. The second round decodes independently on the new table, consistent with Theorem~\ref{thm:protocol}. All second-round recovery results match the true difference element by element. The rounds use independent cell hashes and share the sign hash and checksum hash. We count 20 bytes per cell (4 bytes for the transmitted count and 8 bytes each for keySum and checksum), plus an 8-byte global counter for random-sign encoding. All upper bounds use the same keys and hashes. For each first-round-failure trial and each bound, we record $U_b$, the second-round size, and the second-round outcome, pairing comparisons by trial. Both levels of $u$ use the same number of trials per configuration, set to $\lceil2\times10^4/\widehat{\Pr}(F_1)\rceil$ with a cap of 3,000,000. The 13 configurations with no first-round failure in the pilot receive a fixed 100,000 trials. Among them, the eight configurations with $k=4,5$ at load 0.95 and $M_1\ge4096$ still have no first-round failure in these trials. These configurations contribute only to the failure-rate and byte statistics. We compute intervals for the success-rate difference using the Newcombe hybrid score method for paired proportion differences.

Table~\ref{tab:e2e} summarizes the protocol results from Section~\ref{sec:exp-e2e} of the main text by $M_1$ and $u$. At $M_1=256$, the final two-round failure rate is at most 0.3\%. The conditional miss rate on balanced configurations agrees with the statistical experiments in Section~\ref{sec:exp-cover}. Second-round failure after $U_b\ge d$ is 0.06--0.08\% at $M_1=256$. The oracle uses a smaller second-round table. In paired comparisons by trial, the Skellam rule's success rate is up to 0.35 percentage points higher than the oracle's on some $M_1=256$ configurations.

\begin{table}[!htbp]
\centering
\caption{Two-round protocol summary ($k=3,4,5$). miss/u gives the median / maximum over balanced configurations with at least 1000 first-round failures. Final failure is the maximum over configurations. Second-round size and total bytes are medians relative to the oracle. The oracle is independent of $u$ and appears only in the upper half.}
\label{tab:e2e}
\begin{tabular}{@{}llrrrr@{}}
\toprule
$M_1$, $u$ & upper bound & miss/u & final failure & second-round size & total bytes\\
\midrule
256, 0.01 & Skellam & 0.73 / 0.87 & 0.24\% & +25.6\% & +14.7\%\\
& Gaussian & 1.00 / 1.12 & 0.25\% & +23.9\% & +13.7\%\\
& oracle & -- & 0.51\% & 0 & 0\\
1024, 0.01 & Skellam & 0.88 / 1.00 & 0.06\% & +11.4\% & +6.5\%\\
& Gaussian & 1.04 / 1.17 & 0.08\% & +11.1\% & +6.3\%\\
& oracle & -- & 0.09\% & 0 & 0\\
4096, 0.01 & Skellam & 0.92 / 1.01 & 0.02\% & +5.4\% & +3.1\%\\
& Gaussian & 1.01 / 1.09 & 0.03\% & +5.3\% & +3.0\%\\
& oracle & -- & 0.03\% & 0 & 0\\
16384, 0.01 & Skellam & 0.96 / 1.07 & 0.01\% & +2.6\% & +1.5\%\\
& Gaussian & 1.02 / 1.13 & 0.01\% & +2.6\% & +1.5\%\\
& oracle & -- & 0.02\% & 0 & 0\\
\addlinespace
256, 0.05 & Skellam & 0.93 / 1.01 & 0.29\% & +16.9\% & +9.7\%\\
& Gaussian & 1.01 / 1.07 & 0.27\% & +16.3\% & +9.4\%\\
1024, 0.05 & Skellam & 0.95 / 1.05 & 0.13\% & +7.8\% & +4.5\%\\
& Gaussian & 1.02 / 1.09 & 0.11\% & +7.7\% & +4.4\%\\
4096, 0.05 & Skellam & 0.97 / 1.04 & 0.02\% & +3.8\% & +2.2\%\\
& Gaussian & 1.01 / 1.07 & 0.02\% & +3.8\% & +2.1\%\\
16384, 0.05 & Skellam & 1.00 / 1.05 & 0.01\% & +1.9\% & +1.1\%\\
& Gaussian & 1.01 / 1.06 & 0.02\% & +1.9\% & +1.1\%\\
\bottomrule
\end{tabular}
\end{table}

We use the calibration rule from \cite{wu2026selfsizing} as the baseline. For $k=3$, that work groups results by first-round size and takes the worst 1\% quantile $q_{0.01}$ of post-failure $\widehat d/d$ over several loads and sign compositions (0.8045, 0.8997, 0.9482 for $M_1=256,1024,4096$). We use its upper-bound component $U=\lceil\widehat d/q_{0.01}\rceil$ with the same second-round table length as our two rules. The comparison excludes that protocol's capacity coefficient, deduction of elements recovered in the first round, and joint peeling across the two rounds. Table~\ref{tab:vldb} compares this rule under both encodings, with and without accounting for $|b|$.

\begin{table}[!htbp]
\centering
\caption{Calibration rule versus the Skellam rule ($k=3$, $u=0.01$). The conditional miss rate /u gives the median / maximum for balanced configurations and the range for one-sided configurations. Final failure is the maximum. Total bytes are the median relative to the oracle.}
\label{tab:vldb}
\small
\begin{tabular}{@{}rlcccc@{}}
\toprule
$M_1$ & rule & balanced & one-sided & final failure & total bytes\\
\midrule
256 & standard encoding, calibration factor & 0.92 / 1.02 & 0.09--0.90 & 0.27\% & +13.9\%\\
& standard encoding, calibration factor, with $|b|$ & 0.92 / 1.02 & 0 & 0.26\% & +13.9\%\\
& random sign, calibration factor & 0.87 / 0.97 & 0.85--0.93 & 0.23\% & +13.9\%\\
& random sign, Skellam rule & 0.72 / 0.80 & 0 & 0.24\% & +14.6\%\\
1024 & standard encoding, calibration factor & 0.88 / 0.94 & 0.05--1.03 & 0.07\% & +6.4\%\\
& standard encoding, calibration factor, with $|b|$ & 0.88 / 0.94 & 0 & 0.07\% & +6.4\%\\
& random sign, calibration factor & 0.98 / 1.03 & 0.80--1.04 & 0.08\% & +6.4\%\\
& random sign, Skellam rule & 0.89 / 0.92 & 0 & 0.06\% & +6.5\%\\
4096 & standard encoding, calibration factor & 0.84 / 0.86 & 0.36--0.89 & 0.02\% & +3.1\%\\
& standard encoding, calibration factor, with $|b|$ & 0.84 / 0.86 & 0 & 0.02\% & +3.1\%\\
& random sign, calibration factor & 0.90 / 0.94 & 0.77--0.95 & 0.02\% & +3.1\%\\
& random sign, Skellam rule & 0.97 / 1.01 & 0 & 0.01\% & +3.1\%\\
\bottomrule
\end{tabular}
\end{table}

With the same calibration factor, the conditional miss rate on one-sided configurations ranges from $0.05u$ to $1.03u$ under standard encoding and from $0.77u$ to $1.04u$ under random-sign encoding. At fixed load 0.95--1.0, the rates for one-sided and balanced differences under standard encoding are $0.05$--$0.40u$ and $0.86$--$0.94u$, respectively. Under random-sign encoding, they are $0.80$--$0.95u$ and $0.85$--$1.01u$. Accounting for $|b|$ gives zero one-sided miss under both encodings. In paired comparisons by trial, with the calibration factor and the choice of whether to account for $|b|$ fixed, the success-rate difference between random-sign encoding and standard encoding lies within $\pm1\times10^{-3}$. Intervals exclude zero for 2 of the 30 comparisons, both at $M_1=256$ and load 1.0. The median total-byte difference is about 0.01\%.

\subsection{Computation Cost}\label{app:timing}
For the computation-cost measurement (\exptag{COST}), we use a single-threaded Go implementation on a dual-socket Linux host with Intel Xeon Silver 4410Y processors (nominal 2.0 GHz, 24 physical cores and 48 logical CPUs). Before timing, we pause the host's task queue and wait for all tasks in the current batch to finish, then pin the program to one core. We record that core's frequency and system load every 10 seconds. The frequency remains at 2.001 GHz throughout the run.

We store Poisson tables in a main-memory cache managed by the program. Each cache entry corresponds to one value in a candidate $n$'s Poisson table and stores its probability and cumulative probability as two double-precision numbers, totaling 16 bytes. The cache holds 16,000,000 entries (about 256 MB). When the cache exceeds its limit, the program evicts tables in least-recently-used order until the occupancy plus the size of the incoming table is at most half the capacity, then inserts the new table. We time sizing under three cache conditions. In cold, each sizing starts with an empty cache. In session, the same receiver performs consecutive synchronizations at fixed $(M,k)$, each possibly involving one or two protocol rounds and each processing a new observation. The cache persists between synchronizations. In warm, we immediately recompute the same observation once.

For sizing timings, we use the first 200 failure records per configuration from the statistical experiments at $M=16384$. We repeat each of the four modes, cold, session, warm, and Gaussian rule, 10 times, rotating their order across repetitions. For protocol timings, we use $M_1=16384$ and 100,000 common elements. We repeat the same 40 trials per configuration 10 times and rotate the branch order across repetitions. We count first-round decoding time per trial. Figure~\ref{fig:sizing} and Table~\ref{tab:proto} report results from two runs on the same host with the same settings. The protocol results differ by less than 1\% from the first run.

Appendix~\ref{app:search} predicts sizing cost roughly proportional to $\lambda_d$, or approximately $c^2$. In Figure~\ref{fig:sizing}, $M$ is fixed and load is proportional to $c$. The cold time grows roughly as load to the power 1.9, or $c^{1.9}$, from 0.06--0.12 ms at load about 0.08 ($c=8$) to 69--84 ms at load 3. The session time is about one fifth of cold, and warm is about one sixth to one seventh. At load 3, the Poisson tables requested cumulatively in session exceed the cache limit. The program evicts some old tables, and the session time rises to 41--55 ms. Gaussian rule sizing takes 0.06--0.29 $\mu$s. In the full protocol, the total times for the Gaussian rule and oracle differ by less than 1\%. Peak session memory for Poisson tables is 49--86 MB at load 1.05 and 128--174 MB at load 1.5, reaching the 256 MB limit at load 3.

\begin{figure}[!htbp]
\centering\includegraphics[width=\linewidth]{../figure/fig-sizing.pdf}
\caption{Skellam rule time per sizing ($M=16384$, single thread, cache limit 16,000,000 entries). (a) cold; (b) session (solid) and warm (dashed).}
\label{fig:sizing}
\end{figure}

The decoder treats cells with count $\pm1$ as candidate pure cells and verifies them with the checksum. Under standard encoding with one-sided differences, all keys have the same sign. A cell with count $\pm1$ therefore contains exactly one key, and every candidate passes. With both positive and negative difference elements, three keys with signs $+1,+1,-1$, for example, can occupy the same cell and also give count $+1$. The checksum rejects such candidates. We measure pure-cell checks, checksum rejections, and sign hashes per recovered key in a separate run (30 trials per configuration, 3 repetitions). Random-sign encoding checks an average of 1.77 candidates ($k=3$) and 2.53 candidates ($k=4$). Standard encoding checks 1.00 under one-sided differences and matches random-sign encoding under balanced differences. The first-round decoding time ratio (random sign / standard) is 0.99--1.06 for balanced differences and 1.26--2.42 for one-sided differences at load 1.0--1.5. At load 3, the one-sided ratio is 3.1 because standard encoding stops almost immediately under one-sided differences. Both absolute times are only 0.04--0.16 ms, less than 1\% of total protocol time.

For large parameters, we can replace the exact Skellam CDF with the Gaussian approximation. Across the 33 configurations in the exploration phase with Poisson mean above 200, this changes $U$ for at most 1.8\% of samples and leaves the miss count unchanged. We use the exact rule for all reported results.

\end{document}
